% Leçon 12 — Expression écrite : caractères liés à l'importance de la pratique du sport — Chinois 1ère A — Cours signé OnBuch+
% Programme officiel MINESEC. Compiler avec : tectonic ZHA12-ecrit-caracteres-lies-a-l-importance-de-la-pratiqu.tex
\newcommand{\DOCMATIERE}{Chinois}
\newcommand{\DOCNIVEAU}{1ère A}
\newcommand{\DOCMODULE}{Module 2 : Environnement et santé}
\newcommand{\DOCLECON}{ZHA12}
\newcommand{\DOCTITRE}{Leçon 12 — Expression écrite : caractères liés à l'importance de la pratique du sport}
% OnBuch+ — préambule « cours signés » ENRICHI (compilé avec tectonic / XeLaTeX)
% Système visuel « Pop » d'OnBuch+ : fond crème (jamais blanc), bordures encre
% épaisses, ombres « dures » décalées, titres Archivo Black en capitales.
% Version MATHÉMATIQUES : graphes (pgfplots/tikz), boîtes pédagogiques
% (définition, propriété, méthode, attention, le savais-tu, exercice résolu,
% exercice, TP, bilan), page de garde et signature OnBuch+ ancrée en bas.
%
% Macros attendues dans main.tex AVANT \input{preamble} :
%   \DOCMATIERE  \DOCNIVEAU  \DOCMODULE  \DOCLECON (numéro)  \DOCTITRE
\documentclass[11pt]{article}

\usepackage{fontspec}
\usepackage[french]{babel} % français = langue principale ; le chinois via \zh{..} / textebox (xeCJK)
\usepackage{xeCJK}
\usepackage{pifont} % \ding{51} \ding{55} utilisés par les modèles
\usepackage[a4paper,margin=2cm,top=3.4cm,bottom=2.8cm,headheight=1.4cm,footskip=1.3cm]{geometry}
\usepackage{amsmath,amssymb,amsfonts}
\usepackage{mathtools} % \xmapsto, \coloneqq... souvent utilisés par les modèles
\usepackage{cancel} % simplifications barrées (bilans de forces)
% \abs{z} (module, valeur absolue) et \norm{u} (norme), utilisés par les modèles
\providecommand{\abs}{}\let\abs\relax\DeclarePairedDelimiter{\abs}{\lvert}{\rvert}
\providecommand{\norm}{}\let\norm\relax\DeclarePairedDelimiter{\norm}{\lVert}{\rVert}
\usepackage[table]{xcolor} % [table] : fournit aussi \rowcolors (alternance de lignes)
\usepackage{pagecolor}
\usepackage{tikz}
\usetikzlibrary{arrows.meta,positioning,calc,shapes.geometric,shapes.misc,shapes.arrows,decorations.pathmorphing,decorations.pathreplacing,decorations.markings,patterns,shadows,mindmap,trees,fit,backgrounds,matrix,intersections,angles,quotes}
\usepackage{pgfplots}
\usepackage[version=4]{mhchem} % équations nucléaires \ce{^{235}_{92}U}
\usepackage[european,siunitx]{circuitikz} % schémas électriques
\pgfplotsset{compat=1.18}
\usepgfplotslibrary{fillbetween,groupplots} % aires entre courbes (calcul intégral)
\usepackage{enumitem}
\usepackage{fancyhdr}
% [most] charge toutes les bibliothèques tcolorbox (listings, minted, raster,
% documentation...) et gonfle inutilement la mémoire TeX sur un document long
% (~300 boîtes) : on ne charge que ce qui est réellement utilisé.
\usepackage[skins,breakable]{tcolorbox}
\usepackage{graphicx}
\usepackage{colortbl}
\usepackage{booktabs}
\usepackage{tabularx}
\usepackage{diagbox} % case d'en-tête diagonale des tableaux à double entrée
% Colonnes extensibles centrée / gauche / droite (C, L, R), souvent utilisées par les modèles
\newcolumntype{C}{>{\centering\arraybackslash}X}
\newcolumntype{L}{>{\raggedright\arraybackslash}X}
\newcolumntype{R}{>{\raggedleft\arraybackslash}X}
\usepackage{array}
\usepackage{multicol}
\usepackage{multirow}
\usepackage{siunitx}
% parse-numbers=false : accepte aussi \SI{2,5 \times 10^{-3}}{...}
\sisetup{locale=FR,output-decimal-marker={,},group-minimum-digits=4,parse-numbers=false}
\DeclareSIUnit\ohm{\ensuremath{\symup{\Omega}}} % U+2126 absent de la police
\DeclareSIUnit\division{div} % réglages d'oscilloscope (ms/div, V/div)
\DeclareSIUnit\tr{tr} % tours (vitesse de rotation, ex. \SI{1500}{\tr\per\minute})
\DeclareSIUnit\molar{M} % concentration molaire (ex. \SI{3}{\molar}, \SI{1}{\micro\molar})
\DeclareSIUnit\an{an} % année (le modèle confond parfois avec \year, un compteur TeX) — ex. \SI{5}{\kilo\meter\per\an}
\DeclareSIUnit\FCFA{FCFA} % franc CFA (ex. \SI{500}{\FCFA\per\kilo\gram})
\DeclareSIUnit\degreeBrix{\textdegree Brix} % teneur en sucre d'un sirop/jus (réfractométrie)
\DeclareSIUnit\month{mois} % le modèle utilise parfois \month, un compteur TeX, comme unité de durée
\usepackage{qrcode}
% \degree : commande fréquemment utilisée par les modèles pour un angle ou une
% température en dehors de \SI{}{} (non fournie par les paquets chargés).
\providecommand{\degree}{^{\circ}}
% \qed (fin de démonstration), utilisé par les modèles sans amsthm
\providecommand{\qed}{\hfill$\square$}
% \barre : double barre des valeurs interdites dans un tableau de variations (tkz-tab)
\providecommand{\barre}{\ensuremath{\|}}
% environnement proof (amsthm non chargé) utilisé par les modèles
\ifdefined\proof\else\newenvironment{proof}[1][Démonstration]{\par\noindent\textit{#1.}\ }{\qed\par}\fi
\usepackage{lastpage}
\usepackage{hyperref}
\hypersetup{hidelinks}

% ============================================================
% Palette « Pop » OnBuch+ — mêmes hex que la classe P (plus_ui.dart)
% ============================================================
\definecolor{poporange}{HTML}{F59321}
\definecolor{popdark}{HTML}{E07A0C}
\definecolor{popcream}{HTML}{FFF6E8}
\definecolor{popink}{HTML}{1C1714}
\definecolor{poppurple}{HTML}{7A5AE0}
\definecolor{popgreen}{HTML}{2E9E6A}
\definecolor{poppink}{HTML}{E0457B}
\definecolor{popblue}{HTML}{2F7DE1}
\definecolor{popgold}{HTML}{C98A12}
\definecolor{popmuted}{HTML}{8A7F76}
\definecolor{popline}{HTML}{EADFD2}
\definecolor{popgray}{HTML}{8A7F76}
\colorlet{popgrey}{popgray}
% Alias tolérants (le modèle nomme parfois les couleurs différemment)
\colorlet{popwhite}{white}
\colorlet{popblack}{popink}
\colorlet{popred}{poppink}
\colorlet{popyellow}{popgold}
% Teintes claires (fonds de tableaux/encadrés) — le modèle nomme parfois
% les couleurs avec un suffixe L (ex. popblueL) sans les avoir définies.
\colorlet{poporangeL}{poporange!20}
\colorlet{popdarkL}{popdark!20}
\colorlet{poppurpleL}{poppurple!20}
\colorlet{popgreenL}{popgreen!20}
\colorlet{poppinkL}{poppink!20}
\colorlet{popblueL}{popblue!20}
\colorlet{popgoldL}{popgold!20}
\colorlet{popredL}{popred!20}
\colorlet{popgrayL}{popgray!20}
% Teintes claires pour fonds de tableaux / zones
\colorlet{popblueL}{popblue!10!white}
\colorlet{poporangeL}{poporange!14!white}
\colorlet{popgreenL}{popgreen!12!white}
\colorlet{poppurpleL}{poppurple!12!white}
\colorlet{poppinkL}{poppink!12!white}
\colorlet{popgoldL}{popgold!14!white}

\pagecolor{popcream}

% ============================================================
% Typographie
% ============================================================
\defaultfontfeatures{Ligatures=TeX}
\setmainfont{PlusJakartaSans-Medium.ttf}[Path=../../../fonts/,BoldFont=PlusJakartaSans-Bold.ttf]
% Symboles, API et emojis absents de la police principale : repli sur DejaVu Sans (caractères actifs)
\newfontface\symfont{DejaVuSans.ttf}[Path=../../../fonts/]
\makeatletter
\newcommand\symactive[1]{\catcode#1=13 \begingroup\lccode`\~=#1 \lowercase{\endgroup\def~}{{\symfont\char#1}}}
\newcommand\symblank[1]{\catcode#1=13 \begingroup\lccode`\~=#1 \lowercase{\endgroup\def~}{}}
\newcommand\symhyph[1]{\catcode#1=13 \begingroup\lccode`\~=#1 \lowercase{\endgroup\def~}{\mbox{-}}}
\newcount\symc
\newcommand\symrange[2]{\symc=#1 \@whilenum\symc<#2 \do{\expandafter\symactive\expandafter{\the\symc}\advance\symc by 1 }}
\newcommand\symblankrange[2]{\symc=#1 \@whilenum\symc<#2 \do{\expandafter\symblank\expandafter{\the\symc}\advance\symc by 1 }}
\makeatother
\symrange{"0250}{"02FF}\symactive{"2713}\symactive{"2714}\symactive{"2717}\symactive{"2718}\symactive{"2610}\symactive{"2611}\symactive{"2612}
\symrange{"2600}{"27BF}
\catcode"2705=13 \begingroup\lccode`\~="2705 \lowercase{\endgroup\def~}{{\symfont\char"2713}}\symblank{"26BD}\symblank{"26A1}
\symrange{"2460}{"257F}\symactive{"223C}\symactive{"2248}\symactive{"2260}
\symhyph{"2011}\symhyph{"2E1A}\symblankrange{"1F300}{"1FAFF}\symblank{"FE0F}
% Caractères chinois : WenQuanYi Zen Hei (sous-ensemble GB2312 + pinyin), gras/italique simulés
\setCJKmainfont{WenQuanYiZenHei-subset.ttf}[Path=../../../fonts/,AutoFakeBold=3,AutoFakeSlant=0.2]
\xeCJKsetup{CJKspace=false}
% Pinyin : voyelles à caron (ǎ ǐ ǒ ǔ ǖ ǘ ǚ ǜ) absentes de la police principale -> caractères actifs
\newfontface\pyfont{WenQuanYiZenHei-subset.ttf}[Path=../../../fonts/]
\newcommand\pyactive[1]{\catcode#1=13 \begingroup\lccode`\~=#1 \lowercase{\endgroup\def~}{{\pyfont\char#1}}}
\pyactive{"01CD}\pyactive{"01CE}\pyactive{"01CF}\pyactive{"01D0}\pyactive{"01D1}\pyactive{"01D2}\pyactive{"01D3}\pyactive{"01D4}\pyactive{"01D5}\pyactive{"01D6}\pyactive{"01D7}\pyactive{"01D8}\pyactive{"01D9}\pyactive{"01DA}\pyactive{"01DB}\pyactive{"01DC}
% Maths en sans-serif (Fira Math) avec chiffres et lettres droites de Plus
% Jakarta, pour que les formules se fondent dans le texte.
\usepackage[math-style=ISO,bold-style=ISO]{unicode-math}
\setmathfont{FiraMath-Regular.otf}[Path=../../../fonts/]
\setmathfont{PlusJakartaSans-Medium.ttf}[Path=../../../fonts/,range={up/num,up/{Latin,latin},\mathup}]
\usepackage{icomma} % virgule décimale sans espace en mode math
% Caractères mathématiques Unicode absents des polices de texte (Plus Jakarta,
% Archivo Black) : rendus par leur équivalent mathématique, en texte comme en
% formule — sinon ils s'affichent en carré vide (ex. « Arithmétique dans ℤ »).
\usepackage{newunicodechar}
\newunicodechar{ℝ}{\ensuremath{\mathbb{R}}}
\newunicodechar{ℤ}{\ensuremath{\mathbb{Z}}}
\newunicodechar{ℂ}{\ensuremath{\mathbb{C}}}
\newunicodechar{ℕ}{\ensuremath{\mathbb{N}}}
\newunicodechar{ℚ}{\ensuremath{\mathbb{Q}}}
\newunicodechar{𝔻}{\ensuremath{\mathbb{D}}}
\newunicodechar{α}{\ensuremath{\alpha}}
\newunicodechar{β}{\ensuremath{\beta}}
\newunicodechar{γ}{\ensuremath{\gamma}}
\newunicodechar{δ}{\ensuremath{\delta}}
\newunicodechar{ε}{\ensuremath{\varepsilon}}
\newunicodechar{θ}{\ensuremath{\theta}}
\newunicodechar{λ}{\ensuremath{\lambda}}
\newunicodechar{μ}{\ensuremath{\mu}}
\newunicodechar{σ}{\ensuremath{\sigma}}
\newunicodechar{τ}{\ensuremath{\tau}}
\newunicodechar{φ}{\ensuremath{\varphi}}
\newunicodechar{ψ}{\ensuremath{\psi}}
\newunicodechar{ω}{\ensuremath{\omega}}
\newunicodechar{Σ}{\ensuremath{\Sigma}}
% majuscules produites par \MakeUppercase dans les titres (α-aminés -> Α)
\newunicodechar{Α}{\ensuremath{\alpha}}
\newunicodechar{Β}{\ensuremath{\beta}}
\newunicodechar{Γ}{\ensuremath{\Gamma}}
\newunicodechar{Δ}{\ensuremath{\Delta}}
\newunicodechar{Ε}{\ensuremath{\varepsilon}}
\newunicodechar{Θ}{\ensuremath{\Theta}}
\newunicodechar{Λ}{\ensuremath{\Lambda}}
\newunicodechar{Μ}{\ensuremath{\mu}}
\newunicodechar{Τ}{\ensuremath{\tau}}
\newunicodechar{Φ}{\ensuremath{\Phi}}
\newunicodechar{Ψ}{\ensuremath{\Psi}}
\newunicodechar{Ω}{\ensuremath{\Omega}}
\newunicodechar{↦}{\ensuremath{\mapsto}}
\newunicodechar{∈}{\ensuremath{\in}}
\newunicodechar{∉}{\ensuremath{\notin}}
\newunicodechar{∧}{\ensuremath{\wedge}}
\newunicodechar{⇔}{\ensuremath{\Leftrightarrow}}
\newunicodechar{⟺}{\ensuremath{\Longleftrightarrow}}
\newunicodechar{⇒}{\ensuremath{\Rightarrow}}
\newunicodechar{⟹}{\ensuremath{\Longrightarrow}}
\newunicodechar{∘}{\ensuremath{\circ}}
\newunicodechar{∪}{\ensuremath{\cup}}
\newunicodechar{∩}{\ensuremath{\cap}}
\newunicodechar{≡}{\ensuremath{\equiv}}
\newunicodechar{⊂}{\ensuremath{\subset}}
\newunicodechar{⊥}{\ensuremath{\perp}}
\newunicodechar{∃}{\ensuremath{\exists}}
\newunicodechar{∀}{\ensuremath{\forall}}
\newunicodechar{∛}{\ensuremath{\sqrt[3]{\,}}}
\newunicodechar{∜}{\ensuremath{\sqrt[4]{\,}}}
\newunicodechar{⃗}{\ensuremath{{}^{\to}}}
\newunicodechar{ⁿ}{\ifmmode^{n}\else\textsuperscript{\ensuremath{n}}\fi}
\newunicodechar{ˣ}{\ifmmode^{x}\else\textsuperscript{\ensuremath{x}}\fi}
\newunicodechar{ᵇ}{\ifmmode^{b}\else\textsuperscript{\ensuremath{b}}\fi}
\newunicodechar{ʳ}{\ifmmode^{r}\else\textsuperscript{\ensuremath{r}}\fi}
\newunicodechar{ᵘ}{\ifmmode^{u}\else\textsuperscript{\ensuremath{u}}\fi}
\newunicodechar{ʸ}{\ifmmode^{y}\else\textsuperscript{\ensuremath{y}}\fi}
\newunicodechar{ᵃ}{\ifmmode^{a}\else\textsuperscript{\ensuremath{a}}\fi}
\newunicodechar{ᵉ}{\ifmmode^{e}\else\textsuperscript{\ensuremath{e}}\fi}
\newunicodechar{ᵏ}{\ifmmode^{k}\else\textsuperscript{\ensuremath{k}}\fi}
\newunicodechar{ᵖ}{\ifmmode^{p}\else\textsuperscript{\ensuremath{p}}\fi}
\newunicodechar{ᴺ}{\ifmmode^{N}\else\textsuperscript{\ensuremath{N}}\fi}
\newunicodechar{ᶜ}{\ifmmode^{c}\else\textsuperscript{\ensuremath{c}}\fi}
\newunicodechar{ʰ}{\ifmmode^{h}\else\textsuperscript{\ensuremath{h}}\fi}
\newunicodechar{ᵠ}{\ifmmode^{q}\else\textsuperscript{\ensuremath{q}}\fi}
\newunicodechar{ₙ}{\ifmmode_{n}\else\textsubscript{\ensuremath{n}}\fi}
\newunicodechar{ₐ}{\ifmmode_{a}\else\textsubscript{\ensuremath{a}}\fi}
\newunicodechar{ᵢ}{\ifmmode_{i}\else\textsubscript{\ensuremath{i}}\fi}
\newunicodechar{ᵣ}{\ifmmode_{r}\else\textsubscript{\ensuremath{r}}\fi}
\newunicodechar{ₖ}{\ifmmode_{k}\else\textsubscript{\ensuremath{k}}\fi}
\newunicodechar{ᵦ}{\ifmmode_{\beta}\else\textsubscript{\ensuremath{\beta}}\fi}
\newunicodechar{ₓ}{\ifmmode_{x}\else\textsubscript{\ensuremath{x}}\fi}
\newfontfamily\popdisplay{ArchivoBlack-Regular.ttf}[Path=../../../fonts/]
\newfontfamily\poplogo{SpaceGrotesk-Bold.ttf}[Path=../../../fonts/]
\newfontfamily\popsemi{PlusJakartaSans-SemiBold.ttf}[Path=../../../fonts/]
\newfontfamily\popxbold{PlusJakartaSans-ExtraBold.ttf}[Path=../../../fonts/]
\newfontfamily\popxboldls{PlusJakartaSans-ExtraBold.ttf}[Path=../../../fonts/,LetterSpace=3.0]

\setlist[enumerate]{leftmargin=1.7em,itemsep=5pt,topsep=4pt}
\setlist[itemize]{leftmargin=1.7em,itemsep=4pt,topsep=4pt,label={\color{poporange}\textbullet}}
\setlength{\parindent}{0pt}
\setlength{\parskip}{5pt}
\renewcommand{\arraystretch}{1.25}

% pgfplots : style Pop par défaut
\pgfplotsset{
  every axis/.append style={axis line style={popink,line width=1pt},tick style={popink},
    grid style={popline,line width=0.5pt},label style={font=\small},tick label style={font=\footnotesize}},
  popcurve/.style={poporange,line width=1.8pt,no markers,smooth},
}
% Même style disponible dans un \draw TikZ simple (hors pgfplots)
\tikzset{popcurve/.style={poporange,line width=1.8pt}}
\tikzset{popgrid/.style={popline,very thin}}
\colorlet{popcurve}{poporange} % le nom du style est parfois utilisé comme couleur

% ============================================================
% Pill / squiggle
% ============================================================
\newcommand{\poppill}[3]{%
  \tikz[baseline=-0.55ex]\node[draw=popink,line width=1.2pt,rounded corners=2.6pt,%
    fill=#2,inner xsep=6.5pt,inner ysep=3pt]{%
    \popxboldls\fontsize{7}{7}\selectfont\color{#3}\MakeUppercase{#1}};%
}
\newcommand{\popsquiggle}[1]{%
  \tikz[baseline=-0.4ex]{
    \draw[#1,line width=2.6pt,line cap=round]
      plot [smooth] coordinates {(0,0.09) (0.34,0.01) (0.68,0.21) (1.02,0.01) (1.36,0.10)};
  }%
}
% Mot-clé surligné (marqueur orange sous le texte)
\newcommand{\cle}[1]{{\popxbold\color{popdark}#1}}

% ============================================================
% En-tête / pied
% ============================================================
\pagestyle{fancy}
\fancyhf{}
\renewcommand{\headrulewidth}{0pt}
\renewcommand{\footrulewidth}{0pt}
\lhead{\raisebox{-0.15ex}{\poplogo\fontsize{13}{13}\selectfont\color{popink}OnBuch\color{poporange}+}}
\rhead{\poppill{\DOCMATIERE\ \textbullet\ \DOCNIVEAU\ \textbullet\ Leçon \DOCLECON}{white}{popink}}
\lfoot{\popxbold\fontsize{7.6}{7.6}\selectfont\color{popmuted}\MakeUppercase{Cours signé OnBuch+}\ \textbullet\ {\popsemi\DOCTITRE}}
\rfoot{\tikz[baseline=-0.6ex]\node[rounded corners=8.5pt,fill=popink,minimum height=17pt,minimum width=17pt,inner xsep=5pt,inner sep=0]{\popxbold\fontsize{8}{8}\selectfont\color{popcream}\thepage\,/\,\pageref*{LastPage}};}

% ============================================================
% Titres
% ============================================================
\newcommand{\chaptitre}[2]{%
  \begin{tcolorbox}[enhanced,colback=poporange,colframe=popink,boxrule=2.6pt,arc=11pt,
      drop shadow=popink,left=15pt,right=15pt,top=12pt,bottom=11pt,before skip=3pt,after skip=18pt]
    \poppill{#2}{popink}{popcream}\par\vspace{9pt}
    {\raggedright\hyphenpenalty=10000\exhyphenpenalty=10000\popdisplay\fontsize{22}{24}\selectfont\color{popcream}\MakeUppercase{#1}\par}
  \end{tcolorbox}
}
% Section numérotée : pastille orange avec le numéro + titre Archivo
\newcounter{popsec}
\newcommand{\coursec}[1]{%
  \stepcounter{popsec}\setcounter{popsub}{0}%
  \par\vspace{16pt}\noindent
  \begin{minipage}{\linewidth}
  \tikz[baseline=-0.7ex]\node[draw=popink,line width=1.6pt,fill=poporange,rounded corners=4pt,minimum size=22pt,inner sep=0]{\popdisplay\fontsize{12}{12}\selectfont\color{popcream}\thepopsec};%
  \hspace{8pt}{\popdisplay\fontsize{15}{17}\selectfont\color{popink}\MakeUppercase{#1}}\par
  \vspace{4pt}\noindent{\color{poporange}\rule{36pt}{3.6pt}}
  \end{minipage}\par\nopagebreak\vspace{8pt}
}
\newcounter{popsub}
\newcommand{\courssub}[1]{%
  \stepcounter{popsub}%
  \par\vspace{9pt}\noindent{\popxbold\fontsize{11.5}{13}\selectfont\color{popink}{\color{poporange}\thepopsec.\thepopsub}\hspace{6pt}#1}\par\nopagebreak\vspace{4pt}
}

% ============================================================
% Boîtes pédagogiques — même recette : fond blanc, bordure épaisse colorée,
% ombre dure encre, badge plein. Titre optionnel [#1].
% ============================================================
\tcbset{popbase/.style={breakable,enhanced,colback=white,boxrule=2.3pt,arc=9pt,
  shadow={3pt}{-3pt}{0pt}{popink},left=14pt,right=14pt,top=11pt,bottom=11pt,before skip=11pt,after skip=15pt,
  fontupper=\popsemi\fontsize{10.6}{14.5}\selectfont\color{popink},
  fontlower=\popsemi\fontsize{10.6}{14.5}\selectfont\color{popink}}}
% Coche dessinée (le glyphe ✓ n'existe pas dans les polices du document)
\newcommand{\popcheck}[1]{\tikz[baseline=-0.5ex]\draw[#1,line width=1.6pt,line cap=round,line join=round](-0.13,0.0)--(-0.04,-0.09)--(0.13,0.1);}
\providecommand{\coche}{\popcheck{popgreen}}
\providecommand{\resultat}[1]{\par\smallskip\noindent\fcolorbox{popgreen}{popgreenL}{\parbox{\dimexpr\linewidth-2\fboxsep-2\fboxrule\relax}{\textbf{Résultat.} #1}}\par\smallskip}
\newcommand{\popboxhead}[3]{\poppill{#1}{#2}{white}\ifx\relax#3\relax\else\hspace{6pt}{\popxbold\fontsize{10.6}{12}\selectfont\color{popink}#3}\fi\par\vspace{7pt}}

\newtcolorbox{aretenir}[1][]{popbase,colframe=popgreen,before upper={\popboxhead{À retenir}{popgreen}{#1}}}
% Texte en chinois (document de lecture, dialogue, extrait) : cadre
% d'étude. Un titre optionnel (source, auteur) s'affiche dans l'en-tête.
\newtcolorbox{textebox}[1][]{popbase,colframe=popblue,before upper={\popboxhead{文本 · Texte}{popblue}{#1}},after upper={}}
% Marques ✓ / ✗ (forme correcte / fautive) souvent utilisées par les modèles
\providecommand{\cross}{\textcolor{poppink}{\ensuremath{\times}}}
\providecommand{\xmark}{\cross}\providecommand{\crossmark}{\cross}\providecommand{\cmark}{\ensuremath{\checkmark}}
% Ligne de réponse / justification à compléter (inventée par les modèles)
\providecommand{\justification}[1]{\par\noindent\textit{Justification~:} \dotfill\par}
% \textipa (package tipa non chargé) : la police ne couvre pas l'API -> simple passage
\providecommand{\textipa}[1]{#1}
% Soulignement (ulem non chargé) : \uline, \ul
\providecommand{\uline}[1]{\underline{#1}}\providecommand{\ul}[1]{\underline{#1}}
% \glossaire{\item ...} inventée par les modèles : liste de glossaire
\providecommand{\glossaire}[1]{\begin{itemize}#1\end{itemize}}
% \alert (beamer) inventée par les modèles : texte mis en évidence
\providecommand{\alert}[1]{\textbf{\textcolor{poppink}{#1}}}
% variantes ASCII d'ordinaux écrites en macro
\providecommand{\iere}{\textsuperscript{re}}\providecommand{\ieme}{\textsuperscript{e}}\providecommand{\ere}{\textsuperscript{re}}\providecommand{\eme}{\textsuperscript{e}}\providecommand{\er}{\textsuperscript{er}}\providecommand{\ie}{\textsuperscript{e}}
% Ordinaux français écrits en macro par les modèles : 1\ière, 3\ième
\newcommand{\ière}{\textsuperscript{re}}\newcommand{\ième}{\textsuperscript{e}}
% \exemple (inventée par les modèles) : étiquette « Exemple : »
\providecommand{\exemple}{\textbf{Exemple~:}\ }
% \square utilisable aussi hors mode math (cases à cocher)
\AtBeginDocument{\let\squareMath\square\renewcommand{\square}{\ensuremath{\squareMath}}}
% Mot ou phrase en chinois à l'intérieur d'un paragraphe français
\newcommand{\zh}[1]{\ifmmode\text{#1}\else{#1}\fi}
\let\eng\zh \let\esp\zh \let\all\zh % alias tolérés
% flèches d'intonation inventées par les modèles
\providecommand{\textsearrow}{}\renewcommand{\textsearrow}{\ensuremath{\searrow}}\providecommand{\textnearrow}{}\renewcommand{\textnearrow}{\ensuremath{\nearrow}}
% \fr{..} : mot français (inventé par les modèles)
\providecommand{\fr}[1]{\textit{#1}}
% zhbox : encadré de texte chinois inventé par les modèles
\newenvironment{zhbox}{\begin{quote}}{\end{quote}}
% \vs : « versus » inventé par les modèles
\providecommand{\vs}{\textit{vs}\ }
% \blank : case à compléter inventée par les modèles
\providecommand{\blank}[1][]{\underline{\hspace{1.6em}}}
\newtcolorbox{exemplebox}[1][]{popbase,colframe=poporange,before upper={\popboxhead{Exemple}{poporange}{#1}}}
\newtcolorbox{definition}[1][]{popbase,colframe=popblue,before upper={\popboxhead{Définition}{popblue}{#1}}}
\newtcolorbox{propriete}[1][]{popbase,colframe=poppurple,before upper={\popboxhead{Propriété}{poppurple}{#1}}}
\newtcolorbox{methode}[1][]{popbase,colframe=popgold,before upper={\popboxhead{Méthode}{popgold}{#1}}}
\newtcolorbox{attention}[1][]{popbase,colframe=poppink,before upper={\popboxhead{Attention}{poppink}{#1}}}
\newtcolorbox{experience}[1][]{popbase,colframe=popblue,colback=popblueL,before upper={\popboxhead{Expérience}{popblue}{#1}}}
% « Le savais-tu ? » : carte violette pleine (contexte, culture, Cameroun)
\newtcolorbox{savaistu}[1][]{popbase,colback=poppurple,colframe=popink,
  fontupper=\popsemi\fontsize{10.4}{14.2}\selectfont\color{white},
  before upper={\poppill{Le savais-tu\,?}{white}{poppurple}\ifx\relax#1\relax\else\hspace{6pt}{\popxbold\color{white}#1}\fi\par\vspace{7pt}}}
% Exercice résolu : énoncé en haut, \tcblower puis la solution sur fond crème
\newcounter{popexo}\newcounter{popexr}
\newtcolorbox{exoresolu}[1][]{popbase,colframe=poporange,colbacklower=popcream,
  segmentation style={popink,line width=1.2pt,dashed},
  before upper={\stepcounter{popexr}\popboxhead{Exercice résolu \thepopexr}{poporange}{#1}},
  before lower={\poppill{Solution}{popink}{popcream}\par\vspace{6pt}}}
% Exercice (non résolu en place) — la correction est dans la partie Corrigés
\newtcolorbox{exercice}[1][]{popbase,colframe=popink,
  before upper={\stepcounter{popexo}\popboxhead{Exercice \thepopexo}{popink}{#1}}}
% Corrigé
\newtcolorbox{corrige}[1][]{popbase,colframe=popgreen,colback=popgreenL,
  before upper={\popboxhead{Corrigé}{popgreen}{#1}}}
% Objectifs / prérequis / bilan
\newtcolorbox{objectifs}{popbase,colframe=popink,colback=poporangeL,
  before upper={\popboxhead{Objectifs de la leçon}{popink}{}}}
\newtcolorbox{prerequis}{popbase,colframe=popink,colback=popblueL,
  before upper={\popboxhead{Prérequis}{popblue}{}}}
\newtcolorbox{bilan}{popbase,colframe=popink,colback=white,boxrule=2.8pt,
  before upper={\popboxhead{Fiche bilan}{poporange}{L'essentiel à connaître pour l'examen}}}
% Figure encadrée avec légende
\newtcolorbox{popfigure}[1][]{enhanced,breakable=false,colback=white,colframe=popline,boxrule=1.4pt,arc=8pt,
  left=10pt,right=10pt,top=10pt,bottom=8pt,before skip=10pt,after skip=12pt,halign=center,
  lower separated=false,halign lower=center,
  fontlower=\popsemi\fontsize{9}{11.5}\selectfont\color{popmuted},
  #1}
\newcommand{\legende}[1]{\par\vspace{6pt}{\popsemi\fontsize{9}{11.5}\selectfont\color{popmuted}#1}}

% ============================================================
% Signature OnBuch+
% ============================================================
\newcommand{\poplogoinline}[1]{{\poplogo\fontsize{#1}{#1}\selectfont\color{popink}OnBuch\color{poporange}+}}
\newcommand{\signatureonbuch}{%
  \begin{tcolorbox}[enhanced,colback=popink,colframe=popink,boxrule=2.2pt,arc=10pt,
      drop shadow=poporange,left=14pt,right=14pt,top=10pt,bottom=10pt,before skip=0pt,after skip=0pt]
    \tikz[baseline=-0.7ex]\node[circle,fill=popgreen,draw=popcream,line width=1.3pt,minimum size=22pt,inner sep=0]{\popcheck{white}};%
    \hspace{9pt}\begin{minipage}[c]{0.62\linewidth}
      {\popxbold\fontsize{12}{13}\selectfont\color{popcream}Signé\ \textbullet\ {\poplogo OnBuch\color{poporange}+}}\par\vspace{2pt}
      {\raggedright\popsemi\fontsize{8}{10.5}\selectfont\color{popline}\MakeUppercase{Équipe pédagogique OnBuch+}\\\MakeUppercase{Conforme au programme officiel MINESEC}\par}
    \end{minipage}\hfill
    {\poplogo\fontsize{22}{22}\selectfont\color{popcream}OnBuch\color{poporange}+}
  \end{tcolorbox}%
}

% ============================================================
% Page de garde
% #1 titre · #2 matière · #3 niveau · #4 module · #5 n° leçon · #6 durée ·
% #7 description / objectifs (texte ou itemize)
% ============================================================
\newcommand{\pagegarde}[7]{%
  \thispagestyle{empty}
  \newgeometry{a4paper,margin=1.7cm,top=1.6cm,bottom=1.4cm}%
  \begin{tikzpicture}[remember picture,overlay]
    \fill[poporange!18] ([xshift=-3.2cm,yshift=-3.6cm]current page.north east) circle (4.6cm);
    \fill[poppurple!14] ([xshift=2.4cm,yshift=7.6cm]current page.south west) circle (3.8cm);
  \end{tikzpicture}%
  {\poplogo\fontsize{30}{30}\selectfont\color{popink}OnBuch\color{poporange}+}\hfill
  \raisebox{0.6ex}{\poppill{Cours signé \textbullet\ Premium}{popink}{popcream}}\par
  \vspace{1pt}\popsquiggle{poppurple}\par
  \vspace{8pt}
  \begin{tcolorbox}[enhanced,colback=poporange,colframe=popink,boxrule=2.8pt,arc=13pt,
      drop shadow=popink,left=18pt,right=18pt,top=12pt,bottom=13pt,after skip=10pt]
    \poppill{#2 \textbullet\ #3}{popink}{popcream}\hspace{5pt}\poppill{#4}{white}{popink}\par\vspace{8pt}
    {\popdisplay\fontsize{14}{14}\selectfont\color{popink}LEÇON #5}\par\vspace{4pt}
    {\raggedright\hyphenpenalty=10000\exhyphenpenalty=10000\popdisplay\fontsize{25}{27}\selectfont\color{popcream}\MakeUppercase{#1}\par}
    \vspace{8pt}
    {\popxbold\fontsize{9.5}{9.5}\selectfont\color{popink}\MakeUppercase{Durée conseillée : #6}}
  \end{tcolorbox}
  \begin{tcolorbox}[enhanced,colback=white,colframe=popink,boxrule=2.2pt,arc=10pt,
      drop shadow=popink,left=16pt,right=16pt,top=10pt,bottom=10pt,after skip=8pt]
    \poppill{Dans cette leçon}{poppurple}{white}\par\vspace{5pt}
    \popsemi\fontsize{9.3}{12.6}\selectfont\color{popink}#7
  \end{tcolorbox}
  \noindent\begin{minipage}[t]{0.487\linewidth}
    \begin{tcolorbox}[enhanced,colback=popgreen,colframe=popink,boxrule=2.2pt,arc=10pt,
        drop shadow=popink,left=11pt,right=11pt,top=10pt,bottom=10pt,height=3.1cm,valign=center]
      \tcbox[colback=white,colframe=popink,boxrule=1.3pt,arc=5pt,boxsep=0pt,nobeforeafter,
        left=3pt,right=3pt,top=3pt,bottom=3pt,valign=center]{\qrcode[height=1.9cm]{https://play.google.com/store/apps/details?id=com.luvvix.onbuch&hl=fr}}%
      \hspace{8pt}\begin{minipage}[c]{0.5\linewidth}
        {\popdisplay\fontsize{11}{12.5}\selectfont\color{white}\MakeUppercase{Télécharge OnBuch}}\par\vspace{4pt}
        {\raggedright\popsemi\fontsize{8.4}{11}\selectfont\color{white}Gratuit \textbullet\ Google Play\\Tuteur IA, annales, concours, résultats\par}
      \end{minipage}
    \end{tcolorbox}
  \end{minipage}\hfill
  \begin{minipage}[t]{0.487\linewidth}
    \begin{tcolorbox}[enhanced,colback=poppurple,colframe=popink,boxrule=2.2pt,arc=10pt,
        drop shadow=popink,left=11pt,right=11pt,top=10pt,bottom=10pt,height=3.1cm,valign=center]
      \tikz[baseline=-0.7ex]\node[circle,fill=white,draw=popink,line width=1.3pt,minimum size=1.6cm,inner sep=0]{\popdisplay\fontsize{24}{24}\selectfont\color{poppurple}+};%
      \hspace{8pt}\begin{minipage}[c]{0.6\linewidth}
        {\popdisplay\fontsize{11}{12.5}\selectfont\color{white}\MakeUppercase{Passe à OnBuch+}}\par\vspace{4pt}
        {\raggedright\popsemi\fontsize{8.4}{11}\selectfont\color{white}Tous les cours signés, Tuteur IA illimité, fiches PDF — onglet OnBuch+ de l'app\par}
      \end{minipage}
    \end{tcolorbox}
  \end{minipage}
  \vspace{8pt}
  \popcontenu
  \vfill
  \signatureonbuch
  \restoregeometry
  \newpage
}

% Rangée « Ce cours contient » (page de garde)
\newcommand{\poptile}[3]{% #1 couleur #2 gros texte #3 légende
  \begin{tcolorbox}[enhanced,colback=#1,colframe=popink,boxrule=2pt,arc=9pt,shadow={2.5pt}{-2.5pt}{0pt}{popink},
      left=6pt,right=6pt,top=9pt,bottom=9pt,halign=center,valign=center,height=1.75cm,width=\linewidth,nobeforeafter]
    {\popdisplay\fontsize{12.5}{13}\selectfont\color{white}\MakeUppercase{#2}}\par\vspace{3pt}
    {\centering\popsemi\fontsize{7.8}{9.5}\selectfont\color{white}#3\par}
  \end{tcolorbox}}
\newcommand{\popcontenu}{%
  \par\noindent{\popxboldls\fontsize{7.5}{7.5}\selectfont\color{popmuted}CE COURS SIGNÉ CONTIENT}\par\vspace{7pt}
  \noindent\begin{minipage}[t]{0.235\linewidth}\poptile{popblue}{Cours}{complet, illustré, pas à pas}\end{minipage}\hfill
  \begin{minipage}[t]{0.235\linewidth}\poptile{poporange}{Méthodes}{et exercices résolus}\end{minipage}\hfill
  \begin{minipage}[t]{0.235\linewidth}\poptile{poppink}{Exercices}{type examen + corrigés}\end{minipage}\hfill
  \begin{minipage}[t]{0.235\linewidth}\poptile{popgreen}{Fiche bilan}{l'essentiel à retenir}\end{minipage}\par}

% Sommaire
\newtcolorbox{sommaire}{enhanced,colback=white,colframe=popink,boxrule=2.2pt,arc=10pt,
  drop shadow=popink,left=18pt,right=18pt,top=13pt,bottom=13pt,before skip=6pt,after skip=20pt,
  before upper={\poppill{Sommaire}{poppurple}{white}\par\vspace{9pt}\popsemi\fontsize{10.6}{14.5}\selectfont\color{popink}}}

% Signature de fin de cours — poussée tout en bas de la dernière page
\newcommand{\findecours}{%
  \par\vfill\nopagebreak
  \begin{center}\popsquiggle{poporange}\end{center}\vspace{4pt}
  \signatureonbuch
}
\AtBeginDocument{\def\llbracket{\mathopen{[\mkern-3.5mu[}}\def\rrbracket{\mathclose{]\mkern-3.5mu]}}} % intervalles d'entiers (glyphe ⟦ absent de la police)
% Glyphes absents de Fira Math : redéfinis à partir de glyphes disponibles
\AtBeginDocument{%
  \def\setminus{\mathbin{\backslash}}\let\smallsetminus\setminus
  \def\vdots{\mathord{\vbox{\baselineskip3.2pt\lineskiplimit0pt\kern4pt\hbox{.}\hbox{.}\hbox{.}}}}%
  \def\ddots{\mathinner{\mkern1mu\raise6.5pt\hbox{.}\mkern2mu\raise3.6pt\hbox{.}\mkern2mu\raise0.7pt\hbox{.}\mkern1mu}}%
  \def\triangle{\mathord{\scalebox{0.9}{$\symup{\Delta}$}}}%
  \def\nabla{\mathord{\raisebox{\depth}{\scalebox{1}[-1]{$\symup{\Delta}$}}}}%
  \def\checkmark{\popcheck{popdark}}%
  \def\star{\mathbin{\ast}}\def\bigstar{\mathord{\ast}}%
  \def\diamond{\mathbin{\scalebox{0.6}{$\Diamond$}}}%
  \def\bigcup{\mathop{\mathchoice{\raisebox{-0.3em}{\scalebox{1.7}{$\cup$}}}{\scalebox{1.25}{$\cup$}}{\cup}{\cup}}\displaylimits}%
  \def\bigcap{\mathop{\mathchoice{\raisebox{-0.3em}{\scalebox{1.7}{$\cap$}}}{\scalebox{1.25}{$\cap$}}{\cap}{\cap}}\displaylimits}%
  \def\lesseqqgtr{\mathrel{\lessgtr}}\def\bigoplus{\mathop{\oplus}\displaylimits}%
}
\providecommand{\ssi}{si et seulement si} % abréviation fréquente
\AtBeginDocument{\providecommand{\wideparen}{\overparen}} % arc de cercle

%% FIGFIX-BEGIN v4 — correctifs globaux des figures (docs/onbuch-generation/tools/figfix)
\usepackage{adjustbox}\usepackage{etoolbox}
% toute figure TikZ/pgfplots plus large que la ligne est réduite à la largeur disponible
\BeforeBeginEnvironment{tikzpicture}{\begin{adjustbox}{max width=\linewidth}}
\AfterEndEnvironment{tikzpicture}{\end{adjustbox}}
% légendes pgfplots lisibles par-dessus les courbes
\pgfplotsset{every axis/.append style={legend style={fill=white,fill opacity=0.92,text opacity=1,draw=black!15,inner sep=3pt}}}
% étiquettes d'arêtes (syntaxe edge["..."]) avec fond blanc
\tikzset{every edge quotes/.append style={fill=white,inner sep=1.2pt}}
%% FIGFIX-END
\begin{document}

\pagegarde{Leçon 12 — Expression écrite : caractères liés à l'importance de la pratique du sport}{Chinois}{1ère A}{Module 2 : Environnement et santé}{ZHA12}{2 h}{Cette leçon d'expression écrite apprend à rédiger un texte simple en chinois sur l'importance de la pratique du sport. Elle entraîne l'écriture correcte des caractères du thème (traits, radicaux, caractères proches) et la traduction dans les deux sens (thème et version).
\vspace{4pt}
\begin{multicols}{2}\small
\begin{itemize}
\item À la fin de cette leçon, je sais écrire correctement les caractères liés au sport et à la santé (运动, 健康, 身体, 锻炼...)
\item je sais distinguer les caractères proches du thème (体/休, 练/炼, 运/远)
\item je sais analyser un modèle de texte simple et repérer sa structure (introduction, arguments, conclusion)
\item je sais utiliser les connecteurs utiles : 因为...所以..., 不但...而且..., 首先, 其次, 总之
\item je sais produire un texte simple de 5 à 8 phrases sur l'importance du sport
\item je sais traduire du français vers le chinois (thème) des phrases simples sur le sport
\end{itemize}\end{multicols}}

\begin{sommaire}
\begin{multicols}{2}
\begin{enumerate}
\item Modèle de texte commenté
\item Structure et connecteurs utiles
\item Écriture correcte des caractères du thème
\item Thème : français → chinois
\item Version : chinois → français
\item Production guidée et grille d'évaluation
\item Activité d'intégration
\item Méthodes et exercices résolus
\item Exercices
\item Corrigés des exercices
\item Fiche bilan
\end{enumerate}
\end{multicols}
\end{sommaire}

\begin{objectifs}
\begin{itemize}
\item À la fin de cette leçon, je sais écrire correctement les caractères liés au sport et à la santé (运动, 健康, 身体, 锻炼...)
\item je sais distinguer les caractères proches du thème (体/休, 练/炼, 运/远)
\item je sais analyser un modèle de texte simple et repérer sa structure (introduction, arguments, conclusion)
\item je sais utiliser les connecteurs utiles : 因为...所以..., 不但...而且..., 首先, 其次, 总之
\item je sais produire un texte simple de 5 à 8 phrases sur l'importance du sport
\item je sais traduire du français vers le chinois (thème) des phrases simples sur le sport
\item je sais traduire du chinois vers le français (version) un court texte sur le sport
\item je sais m'auto-évaluer à l'aide d'une grille d'évaluation (caractères, grammaire, connecteurs, sens)
\end{itemize}
\end{objectifs}

\begin{prerequis}
\begin{itemize}
\item Vocabulaire de base du corps et de la santé (身体, 健康, 病)
\item Vocabulaire des sports et activités (跑步, 打球, 游泳, 足球)
\item Structure de base sujet-verbe-objet et verbes modaux (应该, 可以, 要)
\item La structure de cause 因为...所以... vue en début d'année
\item Les radicaux fréquents (亻, 辶, 火, 身)
\end{itemize}
\end{prerequis}

\newpage

\coursec{Modèle de texte commenté}

\courssub{Situation d'accroche et problématique}

Au lycée de Yaoundé, les élèves de 1ère A participent chaque semaine aux activités sportives : football, course, basket-ball. Le médecin scolaire leur rappelle que le sport est bon pour la santé. En Chine aussi, les élèves font de la gymnastique collective chaque matin dans la cour de l'école. Comment écrire en chinois un petit texte pour expliquer pourquoi il est important de pratiquer un sport ? Cette leçon vous donne les caractères, les connecteurs et la méthode pour y arriver.

\textbf{Problématique :} comment un court texte chinois peut-il défendre une idée (la thèse) avec des arguments bien enchaînés ? Nous allons observer un texte modèle, comprendre sa construction en trois parties, puis repérer les mots qui portent l'argumentation.

\courssub{Lecture du texte modèle}

Lisons d'abord le texte en entier, sans chercher à tout traduire mot à mot. Écoutez le rythme des phrases : elles sont courtes, claires, et chacune joue un rôle précis.

\begin{textebox}[运动很重要 — Le sport est important]
运动很重要。每天运动对我们的身体很好。因为运动可以让我们更健康，所以我们都应该锻炼。运动不但对身体好，而且让我们心情愉快。总之，我们应该多运动。\\
\emph{Yùndòng hěn zhòngyào. Měitiān yùndòng duì wǒmen de shēntǐ hěn hǎo. Yīnwèi yùndòng kěyǐ ràng wǒmen gèng jiànkāng, suǒyǐ wǒmen dōu yīnggāi duànliàn. Yùndòng búdàn duì shēntǐ hǎo, érqiě ràng wǒmen xīnqíng yúkuài. Zǒngzhī, wǒmen yīnggāi duō yùndòng.}\\
« Le sport est très important. Faire du sport chaque jour est très bon pour notre corps. Parce que le sport peut nous rendre plus en bonne santé, nous devrions tous nous entraîner. Le sport est non seulement bon pour le corps, mais il nous rend aussi de bonne humeur. En résumé, nous devrions faire plus de sport. »
\end{textebox}

\textbf{Glossaire du texte :}

\begin{center}
\begin{tabularx}{\linewidth}{|X|X|X|X|}
\hline
\rowcolor{popblueL} Caractères & Pinyin & Nature & Traduction \\
\hline
运动 & yùndòng & nom / verbe & sport ; faire du sport \\
\hline
重要 & zhòngyào & adjectif & important \\
\hline
每天 & měitiān & complément de temps & chaque jour \\
\hline
对 & duì & préposition & envers, pour (quelqu'un) \\
\hline
我们 & wǒmen & pronom & nous \\
\hline
身体 & shēntǐ & nom & corps, santé physique \\
\hline
因为 & yīnwèi & conjonction & parce que \\
\hline
可以 & kěyǐ & verbe modal & pouvoir, permettre \\
\hline
让 & ràng & verbe & faire, rendre, laisser \\
\hline
更 & gèng & adverbe & davantage, plus \\
\hline
健康 & jiànkāng & adjectif / nom & en bonne santé ; santé \\
\hline
所以 & suǒyǐ & conjonction & donc, c'est pourquoi \\
\hline
都 & dōu & adverbe & tous, tout \\
\hline
应该 & yīnggāi & verbe modal & devoir (il faut) \\
\hline
锻炼 & duànliàn & verbe & s'entraîner, faire de l'exercice \\
\hline
不但 & búdàn & conjonction & non seulement \\
\hline
而且 & érqiě & conjonction & mais aussi, et de plus \\
\hline
心情 & xīnqíng & nom & humeur, état d'esprit \\
\hline
愉快 & yúkuài & adjectif & joyeux, agréable \\
\hline
总之 & zǒngzhī & connecteur & en résumé, bref \\
\hline
多 & duō & adverbe / adjectif & plus, beaucoup \\
\hline
\end{tabularx}
\end{center}

\courssub{Observation : la structure en trois parties}

Relisons le texte phrase par phrase et demandons-nous : que fait chaque phrase ?

\begin{itemize}
\item \zh{运动很重要。} — L'auteur annonce son idée principale : c'est la \cle{thèse}.
\item \zh{每天运动对我们的身体很好。} — Première raison générale : le sport est bon pour le corps.
\item \zh{因为运动可以让我们更健康，所以我们都应该锻炼。} — \cle{Argument 1} : cause (\zh{因为}) et conséquence (\zh{所以}).
\item \zh{运动不但对身体好，而且让我们心情愉快。} — \cle{Argument 2} : deux avantages accumulés (\zh{不但...而且...}).
\item \zh{总之，我们应该多运动。} — \cle{Conclusion} : le connecteur \zh{总之} referme le texte et rappelle ce qu'il faut faire.
\end{itemize}

\begin{definition}[Les trois parties d'un texte argumentatif simple]
Un texte argumentatif chinois de niveau HSK 2-3 suit presque toujours le plan : \textbf{thèse} (une phrase qui affirme) $\rightarrow$ \textbf{arguments} (deux ou trois phrases qui justifient, reliées par des connecteurs) $\rightarrow$ \textbf{conclusion} (une phrase qui résume et propose une action). En français, on dirait : introduction, développement, conclusion.
\end{definition}

\begin{popfigure}
\begin{tikzpicture}[node distance=0.9cm, every node/.style={align=center}]
\node[draw=popblue, fill=popblueL, rounded corners, thick, minimum width=6.5cm, minimum height=1cm, align=center] (these){\textbf{Thèse}\\ 运动很重要。};
\node[draw=popgreen, fill=popgreenL, rounded corners, thick, minimum width=6.5cm, minimum height=1cm, below=of these, align=center] (arg1){\textbf{Argument 1 : cause $\rightarrow$ conséquence}\\ 因为...所以...};
\node[draw=poporange, fill=poporangeL, rounded corners, thick, minimum width=6.5cm, minimum height=1cm, below=of arg1, align=center] (arg2){\textbf{Argument 2 : accumulation}\\ 不但...而且...};
\node[draw=poppurple, fill=poppurpleL, rounded corners, thick, minimum width=6.5cm, minimum height=1cm, below=of arg2, align=center] (concl){\textbf{Conclusion}\\ 总之，我们应该多运动。};
\draw[-{Stealth}, very thick, popdark] (these) -- (arg1);
\draw[-{Stealth}, very thick, popdark] (arg1) -- (arg2);
\draw[-{Stealth}, very thick, popdark] (arg2) -- (concl);
\end{tikzpicture}
\legende{Organigramme du texte argumentatif : de la thèse à la conclusion}
\end{popfigure}

\begin{methode}[Analyser un texte modèle en quatre étapes]
\begin{enumerate}
\item \textbf{Souligner la phrase de thèse} : c'est la phrase qui affirme l'idée principale (ici : \zh{运动很重要。}).
\item \textbf{Encadrer les connecteurs} : repérer \zh{因为...所以...}, \zh{不但...而且...}, \zh{总之}. Ce sont les « charnières » du texte.
\item \textbf{Numéroter les arguments} : écrire 1, 2, 3 dans la marge devant chaque raison avancée.
\item \textbf{Repérer la conclusion} : elle commence souvent par \zh{总之} ou \zh{所以} et contient un verbe modal comme \zh{应该}.
\end{enumerate}
\end{methode}

\courssub{Les verbes modaux qui portent l'argumentation}

Dans ce texte, trois petits mots font tout le travail de persuasion : \zh{可以}, \zh{应该} et \zh{让}. Observons-les.

\begin{exemplebox}[Les trois moteurs de l'argumentation]
\zh{运动可以让我们更健康。} \emph{Yùndòng kěyǐ ràng wǒmen gèng jiànkāng.} « Le sport peut nous rendre plus en bonne santé. »\\
\zh{我们都应该锻炼。} \emph{Wǒmen dōu yīnggāi duànliàn.} « Nous devrions tous nous entraîner. »\\
\zh{运动让我们心情愉快。} \emph{Yùndòng ràng wǒmen xīnqíng yúkuài.} « Le sport nous met de bonne humeur. »
\end{exemplebox}

\begin{propriete}[Emploi de 可以, 应该, 让]
\begin{itemize}
\item \zh{可以} + verbe : exprime la \textbf{possibilité} (« pouvoir »). Structure : Sujet + 可以 + Verbe. Exemple : \zh{运动可以让我们更健康。}
\item \zh{应该} + verbe : exprime le \textbf{devoir, le conseil} (« devoir, il faudrait »). Structure : Sujet + 应该 + Verbe. Exemple : \zh{我们应该每天锻炼。} \emph{Wǒmen yīnggāi měitiān duànliàn.} « Nous devrions nous entraîner chaque jour. »
\item \zh{让} + personne + adjectif/verbe : exprime le \textbf{fait de rendre quelqu'un...} (« faire que, rendre »). Structure : Sujet + 让 + Personne + Adjectif. Exemple : \zh{运动让我们心情愉快。}
\end{itemize}
Ces trois mots se placent \textbf{avant le verbe principal} et ne prennent jamais de \zh{了}.
\end{propriete}

\begin{center}
\begin{tabularx}{\linewidth}{|X|X|X|}
\hline
\rowcolor{popblueL} Structure & Exemple (caractères + pinyin) & Traduction \\
\hline
Sujet + 可以 + Verbe & 运动可以让我们健康。Yùndòng kěyǐ ràng wǒmen jiànkāng. & Le sport peut nous rendre en bonne santé. \\
\hline
Sujet + 应该 + Verbe & 我们应该多运动。Wǒmen yīnggāi duō yùndòng. & Nous devrions faire plus de sport. \\
\hline
Sujet + 让 + Personne + Adj. & 运动让我们愉快。Yùndòng ràng wǒmen yúkuài. & Le sport nous rend joyeux. \\
\hline
对 + Personne + 好 & 运动对身体很好。Yùndòng duì shēntǐ hěn hǎo. & Le sport est bon pour le corps. \\
\hline
\end{tabularx}
\end{center}

\begin{attention}[Erreurs fréquentes des francophones]
\begin{itemize}
\item Ne dites pas \zh{我们应该锻炼每天} : en chinois, le complément de temps \zh{每天} se place \textbf{avant le verbe} : \zh{我们应该每天锻炼。}
\item \zh{让} ne signifie pas « laisser » dans ce texte, mais « rendre, faire que ». \zh{让我们更健康} = « nous rendre plus en bonne santé », pas « nous laisser plus en bonne santé ».
\item N'oubliez pas la paire : \zh{因为} appelle presque toujours \zh{所以} dans la même phrase, contrairement au français où « parce que... donc » ensemble est une faute. En chinois, c'est correct et même recommandé !
\end{itemize}
\end{attention}

\courssub{Carte mentale du vocabulaire du thème}

\begin{popfigure}
\begin{tikzpicture}[mindmap, grow cyclic, every node/.style={concept, align=center}, root concept/.append style={fill=popblue, text=white, font=\bfseries}, level 1/.append style={level distance=3.2cm, sibling angle=90}]
\node[root concept, align=center]{运动\\ yùndòng\\ sport}
  child[concept color=popgreenL] { node[align=center]{健康\\ jiànkāng\\ santé} }
  child[concept color=poporangeL] { node[align=center]{身体\\ shēntǐ\\ corps} }
  child[concept color=poppurpleL] { node[align=center]{锻炼\\ duànliàn\\ s'entraîner} }
  child[concept color=poppinkL] { node[align=center]{心情\\ xīnqíng\\ humeur} };
\end{tikzpicture}
\legende{Carte mentale lexicale autour de 运动 (le sport)}
\end{popfigure}

\begin{exoresolu}[Repérer la structure d'une phrase]
\textbf{Énoncé :} Dans la phrase \zh{因为运动可以让我们更健康，所以我们都应该锻炼。}, soulignez la cause, encadrez la conséquence et identifiez les deux verbes modaux.
\tcblower
\textbf{Solution détaillée :}
\begin{enumerate}
\item La cause est introduite par \zh{因为} : \zh{运动可以让我们更健康} (« le sport peut nous rendre plus en bonne santé »).
\item La conséquence est introduite par \zh{所以} : \zh{我们都应该锻炼} (« nous devrions tous nous entraîner »).
\item Les deux verbes modaux sont \zh{可以} (possibilité, dans la cause) et \zh{应该} (devoir, dans la conséquence). Le verbe \zh{让} complète le dispositif argumentatif en exprimant « rendre ».
\end{enumerate}
\end{exoresolu}

\begin{savaistu}[La gymnastique matinale en Chine]
En Chine, la gymnastique matinale collective, \zh{早操} (\emph{zǎocāo}), est obligatoire dans la plupart des écoles : chaque matin, tous les élèves se rassemblent dans la cour et font des exercices en musique avant les cours. Cela ressemble aux séances de sport collectives des lycées du Cameroun, mais en Chine, c'est un rituel quotidien ! Le mot \zh{操} (\emph{cāo}) signifie « exercice, gymnastique » ; il contient la clé de la main \zh{扌}, car on bouge son corps avec les bras et les mains.
\end{savaistu}

\begin{aretenir}[L'essentiel de la section]
\begin{itemize}
\item Un texte argumentatif simple = \textbf{thèse} + \textbf{arguments} (\zh{因为...所以...} ; \zh{不但...而且...}) + \textbf{conclusion} (\zh{总之}).
\item Trois mots portent l'argumentation : \zh{可以} (pouvoir), \zh{应该} (devoir), \zh{让} (rendre).
\item Le complément de temps se place avant le verbe : \zh{每天运动}.
\item Méthode d'analyse : souligner la thèse, encadrer les connecteurs, numéroter les arguments, repérer la conclusion.
\end{itemize}
\end{aretenir}

\coursec{Structure et connecteurs utiles}

Dans la section précédente, nous avons lu et commenté un modèle de texte sur l'importance du sport. Vous avez remarqué que ce texte n'est pas une simple suite de phrases : il est \emph{organisé}. Des petits mots relient les idées entre elles, annoncent les arguments et concluent le propos. Ces mots, ce sont les \cle{connecteurs}. En chinois comme en français, un bon texte argumentatif repose sur un squelette logique clair. Cette section vous donne tous les outils pour construire ce squelette : les connecteurs de cause, d'addition, d'organisation, les formules d'opinion et les structures d'importance. À la fin, vous saurez assembler un paragraphe argumenté de cinq phrases, prêt pour l'examen.

\courssub{Observer avant de comprendre : les connecteurs en action}

\begin{textebox}[Un petit texte bien connecté]
我觉得运动对每个人都很重要。\\
\emph{Wǒ juéde yùndòng duì měi gè rén dōu hěn zhòngyào.}\\
« Je pense que le sport est très important pour tout le monde. »\\[4pt]
因为运动对身体好，所以我们应该每天运动。\\
\emph{Yīnwèi yùndòng duì shēntǐ hǎo, suǒyǐ wǒmen yīnggāi měitiān yùndòng.}\\
« Parce que le sport est bon pour le corps, nous devons faire du sport chaque jour. »\\[4pt]
运动不但让我们健康，而且让我们心情愉快。\\
\emph{Yùndòng búdàn ràng wǒmen jiànkāng, érqiě ràng wǒmen xīnqíng yúkuài.}\\
« Le sport nous rend non seulement en bonne santé, mais aussi joyeux. »\\[4pt]
首先，我们可以在学校踢球；其次，我们可以和朋友一起跑步。\\
\emph{Shǒuxiān, wǒmen kěyǐ zài xuéxiào tī qiú ; qícì, wǒmen kěyǐ hé péngyou yìqǐ pǎobù.}\\
« D'abord, nous pouvons jouer au ballon à l'école ; ensuite, nous pouvons courir avec nos amis. »\\[4pt]
总之，多运动对我们的生活很有好处。\\
\emph{Zǒngzhī, duō yùndòng duì wǒmen de shēnghuó hěn yǒu hǎochù.}\\
« En résumé, faire plus de sport est très bénéfique pour notre vie. »
\end{textebox}

Observez les mots en gras dans votre esprit : \zh{我觉得}, \zh{因为……所以……}, \zh{不但……而且……}, \zh{首先}, \zh{其次}, \zh{总之}, \zh{对……很重要}, \zh{对……有好处}. Ce sont exactement les huit outils de cette leçon. Chacun a une fonction précise, comme les pièces d'un meuble : on ne peut pas les intervertir.

\courssub{Le connecteur de cause-conséquence : 因为……所以……}

\begin{propriete}[Règle de 因为……所以……]
Structure : \zh{因为 + cause，所以 + conséquence。}\\[3pt]
Contrairement au français, où « parce que... donc... » est une faute de syntaxe, le chinois \emph{admet et même préfère} la présence des deux éléments dans la même phrase : \zh{因为} introduit la cause, \zh{所以} introduit la conséquence. On peut omettre l'un des deux (dire seulement \zh{因为……} ou \zh{……所以……}), mais jamais inverser leur ordre : la cause vient toujours \emph{avant} la conséquence.
\end{propriete}

\begin{exemplebox}[Exemples traduits]
\zh{因为运动对身体好，所以我们应该多运动。}\\
\emph{Yīnwèi yùndòng duì shēntǐ hǎo, suǒyǐ wǒmen yīnggāi duō yùndòng.}\\
« Parce que le sport est bon pour le corps, nous devons faire plus de sport. »\\[4pt]
\zh{因为昨天下雨了，所以我们没有踢球。}\\
\emph{Yīnwèi zuótiān xià yǔ le, suǒyǐ wǒmen méiyǒu tī qiú.}\\
« Parce qu'il a plu hier, nous n'avons pas joué au ballon. »\\[4pt]
\zh{因为他每天跑步，所以他很健康。}\\
\emph{Yīnwèi tā měitiān pǎobù, suǒyǐ tā hěn jiànkāng.}\\
« Parce qu'il court tous les jours, il est en bonne santé. »
\end{exemplebox}

\begin{attention}[Erreur fréquente des francophones]
Ne traduisez jamais « donc » par un mot séparé placé n'importe où. En chinois, \zh{所以} se place \emph{directement devant le sujet ou le verbe} de la seconde proposition, jamais en fin de phrase. On dit \zh{所以我们应该多运动} et non \zh{我们应该多运动所以}. Autre piège : ne dites pas « parce que... donc » en français dans vos versions ; en français, il faut choisir : « Comme il court chaque jour, il est en bonne santé » ou « Il court chaque jour, donc il est en bonne santé ».
\end{attention}

\courssub{Le connecteur d'addition : 不但……而且……}

\begin{propriete}[Règle de 不但……而且……]
Structure : \zh{不但 + idée 1，而且 + idée 2。}\\[3pt]
Ce connecteur exprime l'addition renforcée : « non seulement... mais aussi... ». Si les deux propositions ont le \emph{même sujet}, le sujet se place avant \zh{不但}. Si les sujets sont \emph{différents}, chaque sujet se place après son connecteur : \zh{不但运动对身体好，而且运动也让我们快乐。}
\end{propriete}

\begin{exemplebox}[Exemples traduits]
\zh{运动不但对身体好，而且让我们心情愉快。}\\
\emph{Yùndòng búdàn duì shēntǐ hǎo, érqiě ràng wǒmen xīnqíng yúkuài.}\\
« Le sport est non seulement bon pour le corps, mais il nous rend aussi joyeux. »\\[4pt]
\zh{他不但会踢足球，而且会打篮球。}\\
\emph{Tā búdàn huì tī zúqiú, érqiě huì dǎ lánqiú.}\\
« Il sait non seulement jouer au football, mais aussi au basket-ball. »\\[4pt]
\zh{跑步不但不花钱，而且很方便。}\\
\emph{Pǎobù búdàn bù huā qián, érqiě hěn fāngbiàn.}\\
« Courir non seulement ne coûte rien, mais en plus c'est très pratique. »
\end{exemplebox}

\courssub{Les marqueurs d'organisation : 首先、其次、总之}

Pour énumérer des arguments, le chinois utilise des marqueurs d'ordre, placés en tête de phrase et suivis d'une virgule :

\begin{exemplebox}[Exemples traduits]
\zh{首先，运动让我们健康；其次，运动让我们快乐。}\\
\emph{Shǒuxiān, yùndòng ràng wǒmen jiànkāng ; qícì, yùndòng ràng wǒmen kuàilè.}\\
« D'abord, le sport nous rend en bonne santé ; ensuite, il nous rend heureux. »\\[4pt]
\zh{总之，我们应该多运动。}\\
\emph{Zǒngzhī, wǒmen yīnggāi duō yùndòng.}\\
« En résumé, nous devons faire plus de sport. »
\end{exemplebox}

Notez le point-virgule chinois \zh{；} qui relie deux arguments parallèles, exactement comme en français. \zh{总之} annonce toujours la conclusion : il ne peut pas ouvrir un texte.

\courssub{Les formules d'opinion : 我觉得…… / 我认为……}

\begin{propriete}[Exprimer son opinion]
\zh{我觉得 + phrase} : « je pense que... » — registre courant, oral et écrit.\\
\zh{我认为 + phrase} : « je considère que... » — registre plus soutenu, idéal à l'écrit.\\[3pt]
Les deux se placent en tête de phrase et sont suivis directement d'une proposition complète, sans mot de liaison (pas de « que » en chinois).
\end{propriete}

\zh{我觉得喀麦隆的足球很好。} \emph{Wǒ juéde Kāmàilóng de zúqiú hěn hǎo.} « Je pense que le football camerounais est très bon. »

\zh{我认为运动对学习有好处。} \emph{Wǒ rènwéi yùndòng duì xuéxí yǒu hǎochù.} « Je considère que le sport est bénéfique pour les études. »

\courssub{Les structures d'importance : 对……很重要 / 对……有好处}

\begin{propriete}[Parler d'importance et de bénéfice]
Structure 1 : \zh{A 对 B 很重要} = « A est très important pour B ».\\
Structure 2 : \zh{A 对 B 有好处} = « A est bon pour B, A profite à B ».\\[3pt]
Le mot-outil \zh{对} (duì, « envers, pour ») introduit toujours le bénéficiaire. L'ordre est figé : on ne peut pas dire \zh{很重要对身体}. Le bénéficiaire vient juste après \zh{对}, puis vient l'adjectif ou l'expression.
\end{propriete}

\zh{运动对学生很重要。} \emph{Yùndòng duì xuésheng hěn zhòngyào.} « Le sport est très important pour les élèves. »

\zh{多喝水对身体有好处。} \emph{Duō hē shuǐ duì shēntǐ yǒu hǎochù.} « Boire beaucoup d'eau est bon pour le corps. »

\courssub{Tableau récapitulatif des connecteurs}

\begin{center}\begin{tabularx}{\linewidth}{|X|X|X|X|}\hline
\rowcolor{popblueL} Caractères & Pinyin & Fonction & Traduction \\ \hline
因为……所以…… & yīnwèi... suǒyǐ... & cause-conséquence & parce que... donc... \\ \hline
不但……而且…… & búdàn... érqiě... & addition renforcée & non seulement... mais aussi... \\ \hline
首先 & shǒuxiān & premier argument & d'abord \\ \hline
其次 & qícì & argument suivant & ensuite \\ \hline
总之 & zǒngzhī & conclusion & en résumé \\ \hline
我觉得 & wǒ juéde & opinion (courant) & je pense que \\ \hline
我认为 & wǒ rènwéi & opinion (soutenu) & je considère que \\ \hline
对……很重要 & duì... hěn zhòngyào & importance & très important pour \\ \hline
对……有好处 & duì... yǒu hǎochù & bénéfice & bon pour \\ \hline
\end{tabularx}\end{center}

\begin{popfigure}
\begin{tikzpicture}[node distance=0.9cm]
\node[draw=popblue, fill=popblueL, rounded corners, align=center, minimum width=4.6cm] (c1) {因为……所以……};
\node[draw=popblue, fill=popblueL, rounded corners, align=center, minimum width=4.6cm, below=of c1] (c2) {不但……而且……};
\node[draw=popblue, fill=popblueL, rounded corners, align=center, minimum width=4.6cm, below=of c2] (c3) {首先 / 其次 / 总之};
\node[draw=popblue, fill=popblueL, rounded corners, align=center, minimum width=4.6cm, below=of c3] (c4) {我觉得 / 我认为};
\node[draw=popblue, fill=popblueL, rounded corners, align=center, minimum width=4.6cm, below=of c4] (c5) {对……很重要 / 有好处};
\node[draw=poporange, fill=poporangeL, rounded corners, align=center, minimum width=4.6cm, right=3.2cm of c1] (f1) {parce que / donc\\(à choisir en français !)};
\node[draw=poporange, fill=poporangeL, rounded corners, align=center, minimum width=4.6cm, right=3.2cm of c2] (f2) {non seulement\\mais aussi};
\node[draw=poporange, fill=poporangeL, rounded corners, align=center, minimum width=4.6cm, right=3.2cm of c3] (f3) {d'abord / ensuite\\en résumé};
\node[draw=poporange, fill=poporangeL, rounded corners, align=center, minimum width=4.6cm, right=3.2cm of c4] (f4) {je pense que\\je considère que};
\node[draw=poporange, fill=poporangeL, rounded corners, align=center, minimum width=4.6cm, right=3.2cm of c5] (f5) {très important pour\\bon pour};
\draw[-{Stealth}, thick, popdark] (c1) -- (f1);
\draw[-{Stealth}, thick, popdark] (c2) -- (f2);
\draw[-{Stealth}, thick, popdark] (c3) -- (f3);
\draw[-{Stealth}, thick, popdark] (c4) -- (f4);
\draw[-{Stealth}, thick, popdark] (c5) -- (f5);
\end{tikzpicture}
\legende{Correspondances chinois-français des connecteurs du thème « le sport ».}
\end{popfigure}

\courssub{Le plan type en cinq phrases}

Pour toute production écrite sur l'importance du sport, retenez ce plan en cinq phrases, simple et efficace :

\begin{enumerate}
\item \textbf{Thèse} (opinion) : \zh{我觉得运动很重要。}
\item \textbf{Cause} : \zh{因为运动对身体有好处，所以我每天运动。}
\item \textbf{Addition} : \zh{运动不但让我健康，而且让我快乐。}
\item \textbf{Exemple personnel} : \zh{首先，我每天跑步；其次，我常和朋友踢球。}
\item \textbf{Conclusion} : \zh{总之，运动对每个人都很重要。}
\end{enumerate}

\begin{methode}[Construire un paragraphe argumenté]
\begin{enumerate}
\item Annoncez votre opinion avec \zh{我觉得} ou \zh{我认为}.
\item Justifiez avec une cause et une conséquence : \zh{因为……所以……}.
\item Ajoutez un second argument avec \zh{不但……而且……}.
\item Illustrez avec vos exemples personnels introduits par \zh{首先} et \zh{其次}.
\item Concluez avec \zh{总之} et une structure d'importance : \zh{对……很重要} ou \zh{对……有好处}.
\item Relisez : vérifiez l'ordre des mots (cause avant conséquence, \zh{对} devant le bénéficiaire) et la ponctuation chinoise.
\end{enumerate}
\end{methode}

\begin{exoresolu}[Assembler les phrases en un paragraphe]
Énoncé : remettez dans l'ordre les phrases suivantes pour former un paragraphe cohérent, puis traduisez le résultat.\\
(a) \zh{因为运动对身体有好处，所以我每天运动。}\\
(b) \zh{总之，运动对每个人都很重要。}\\
(c) \zh{我觉得运动很重要。}\\
(d) \zh{运动不但让我健康，而且让我快乐。}
\tcblower
Solution détaillée : on commence par l'opinion (c), puis la justification (a), puis l'argument supplémentaire (d), et enfin la conclusion (b). Ordre : c — a — d — b.\\[4pt]
\zh{我觉得运动很重要。因为运动对身体有好处，所以我每天运动。运动不但让我健康，而且让我快乐。总之，运动对每个人都很重要。}\\[3pt]
\emph{Wǒ juéde yùndòng hěn zhòngyào. Yīnwèi yùndòng duì shēntǐ yǒu hǎochù, suǒyǐ wǒ měitiān yùndòng. Yùndòng búdàn ràng wǒ jiànkāng, érqiě ràng wǒ kuàilè. Zǒngzhī, yùndòng duì měi gè rén dōu hěn zhòngyào.}\\[3pt]
Traduction : « Je pense que le sport est très important. Comme le sport est bon pour le corps, je fais du sport chaque jour. Le sport me rend non seulement en bonne santé, mais aussi heureux. En résumé, le sport est très important pour tout le monde. »
\end{exoresolu}

\begin{attention}[Derniers pièges à éviter]
\begin{itemize}
\item Ne mettez pas deux sujets autour de \zh{对} : on dit \zh{运动对身体好}, jamais \zh{运动对身体它好}.
\item \zh{总之} ne s'utilise qu'en conclusion ; pour ouvrir, utilisez \zh{首先}.
\item En version (chinois → français), n'écrivez jamais « parce que... donc » : reformulez avec « comme... » ou supprimez l'un des deux connecteurs.
\item \zh{我觉得} exprime une opinion personnelle ; ne l'utilisez pas pour un fait objectif comme \zh{我觉得地球是圆的}, préférez une phrase affirmative simple.
\end{itemize}
\end{attention}

Vous disposez maintenant de tous les connecteurs nécessaires. Dans la section suivante, nous passerons à l'écriture correcte des caractères du thème, trait par trait, avant de vous entraîner à la traduction.

\coursec{Écriture correcte des caractères du thème}

Dans la section précédente, nous avons construit la \cle{structure} d'un texte sur le sport et appris les \cle{connecteurs} utiles (\zh{因为}…\zh{所以}…, \zh{首先}…\zh{然后}…, \zh{总之}). Mais une belle structure ne suffit pas : en chinois, l'évaluation de l'expression écrite repose aussi sur la \cle{correction des caractères}. Un caractère mal écrit — un trait de trop, un radical confondu — peut changer complètement le sens d'une phrase. Cette section est donc consacrée à l'\cle{écriture correcte} des caractères clés de notre thème : l'importance de la pratique du sport.

\courssub{Observer avant d'écrire : les caractères dans un texte}

Lisons d'abord ce court texte. Repérez les caractères \zh{运}, \zh{动}, \zh{健}, \zh{康}, \zh{体}, \zh{炼} : comment sont-ils construits ? Quels éléments reviennent ?

\begin{textebox}[Le sport et la santé]
运动很重要。每天运动，身体才健康。\\
\emph{Yùndòng hěn zhòngyào. Měitiān yùndòng, shēntǐ cái jiànkāng.}\\
« Le sport est très important. En faisant du sport chaque jour, le corps est en bonne santé. »\\[4pt]
我喜欢锻炼身体。早上我跑步，下午我打球。\\
\emph{Wǒ xǐhuan duànliàn shēntǐ. Zǎoshang wǒ pǎobù, xiàwǔ wǒ dǎqiú.}\\
« J'aime entraîner mon corps. Le matin je cours, l'après-midi je joue au ballon. »\\[4pt]
运动让身体更健康，也让我们更快乐。\\
\emph{Yùndòng ràng shēntǐ gèng jiànkāng, yě ràng wǒmen gèng kuàilè.}\\
« Le sport rend le corps plus sain et nous rend aussi plus heureux. »
\end{textebox}

\textbf{Observations.} Le mot \zh{运动} (\emph{yùndòng}, « sport, mouvement ») revient trois fois : c'est le mot central du thème. On remarque que \zh{运} et \zh{远} se ressemblent, que \zh{健} et \zh{体} commencent tous les deux par la même forme à gauche (\zh{亻}), et que \zh{炼} contient le feu \zh{火}. Chacun de ces détails est une clé de mémorisation. Analysons maintenant chaque caractère un par un.

\courssub{Les six caractères clés : analyse détaillée}

\begin{definition}[运动 — yùndòng : sport, mouvement]
Le mot \zh{运动} se compose de deux caractères :
\begin{itemize}
\item \zh{运} (\emph{yùn}, bouger, transporter) : \cle{7 traits}. Radical \zh{辶} (clé du \cle{mouvement}) + composant phonétique \zh{云} (\emph{yún}, nuage), qui donne le son.
\item \zh{动} (\emph{dòng}, bouger) : \cle{6 traits}. Composant phonétique \zh{云} (\emph{yún}, le son) + radical \zh{力} (\emph{lì}, la \cle{force}), qui donne le sens : un mouvement demande de la force.
\end{itemize}
Sens littéral : « bouger en se déplaçant » → le sport, l'activité physique.
\end{definition}

\begin{definition}[健康 — jiànkāng : santé, sain]
\begin{itemize}
\item \zh{健} (\emph{jiàn}, sain, robuste) : \cle{10 traits}. Radical \zh{亻} (clé de l'\cle{homme}, variante de \zh{人}) à gauche + composant phonétique \zh{建} (\emph{jiàn}, construire) à droite. Image : un homme bien « construit » est robuste.
\item \zh{康} (\emph{kāng}, santé) : \cle{11 traits}. Radical \zh{广} (clé de l'\cle{abri}, de la maison) + composant phonétique \zh{庚} (\emph{gēng}). Image : la santé est un toit protecteur.
\end{itemize}
\zh{健康} est à la fois un nom (« la santé ») et un adjectif (« sain, en bonne santé »).
\end{definition}

\begin{definition}[体 — tǐ : corps]
\zh{体} : \cle{7 traits}. Radical \zh{亻} (homme) + \zh{本} (\emph{běn}, racine, origine). Image : le corps est la « racine » de l'homme. On le retrouve dans \zh{身体} (\emph{shēntǐ}, le corps), \zh{体育} (\emph{tǐyù}, l'éducation physique), \zh{体重} (\emph{tǐzhòng}, le poids).
\end{definition}

\begin{definition}[炼 — liàn : entraîner, forger, fondre]
\zh{炼} : \cle{9 traits}. Radical \zh{火} (clé du \cle{feu}) + composant phonétique \zh{东} (\emph{dōng}). Sens premier : fondre un métal au feu ; sens figuré : « forger » le corps, le caractère.
\end{definition}

\begin{definition}[锻炼 — duànliàn : s'entraîner, se forger]
\zh{锻} (\emph{duàn}, forger le métal) : radical \zh{钅} (métal) + composant phonétique \zh{段} (\emph{duàn}). \zh{炼} (\emph{liàn}) : radical \zh{火} (feu). L'image est magnifique : s'entraîner, c'est \cle{forger son corps comme le forgeron forge le fer} — par le feu, le martelage et la répétition. C'est le mot exact pour « faire de l'exercice physique ».
\end{definition}

\begin{center}
\begin{tabularx}{\linewidth}{|X|X|X|X|}
\hline
\rowcolor{popblueL} Caractère & Pinyin & Radical / Clé (sens) & Traduction \\
\hline
运 & yùn & \zh{辶} mouvement & bouger, transporter \\
\hline
动 & dòng & \zh{力} force & bouger, remuer \\
\hline
健 & jiàn & \zh{亻} homme & sain, robuste \\
\hline
康 & kāng & \zh{广} abri & santé \\
\hline
体 & tǐ & \zh{亻} homme & corps \\
\hline
炼 & liàn & \zh{火} feu & forger, entraîner \\
\hline
\end{tabularx}
\end{center}

\begin{popfigure}
\begin{tikzpicture}[node distance=0.8cm, every node/.style={font=\small, align=center}]
% Caractère 运
\node[draw=popblue, fill=popblueL, rounded corners, minimum width=2.4cm, align=center] (yun){\zh{运} yùn\\bouger};
\node[draw=poporange, fill=poporangeL, rounded corners, minimum width=1.6cm, above left=0.6cm and 0.2cm of yun, align=center] (zou){\zh{辶}\\mouvement};
\node[draw=popgreen, fill=popgreenL, rounded corners, minimum width=1.6cm, above right=0.6cm and 0.2cm of yun, align=center] (yunph){\zh{云} yún\\phonétique};
% Caractère 动
\node[draw=popblue, fill=popblueL, rounded corners, minimum width=2.4cm, right=3.5cm of yun, align=center] (dong){\zh{动} dòng\\bouger};
\node[draw=popgreen, fill=popgreenL, rounded corners, minimum width=1.6cm, above left=0.6cm and 0.2cm of dong, align=center] (dongph){\zh{云} yún\\phonétique};
\node[draw=poppurple, fill=poppurpleL, rounded corners, minimum width=1.6cm, above right=0.6cm and 0.2cm of dong, align=center] (li){\zh{力} lì\\force};
% Mot 运动
\node[draw=popdark, fill=popcream, rounded corners, minimum width=3.0cm, below=1.8cm of $(yun)!0.5!(dong)$, align=center] (sport){\zh{运动} yùndòng\\le sport};
% Flèches
\draw[-{Stealth}, thick, poporange] (zou) -- (yun);
\draw[-{Stealth}, thick, popgreen] (yunph) -- (yun);
\draw[-{Stealth}, thick, popgreen] (dongph) -- (dong);
\draw[-{Stealth}, thick, poppurple] (li) -- (dong);
\draw[-{Stealth}, thick, popdark] (yun) -- (sport);
\draw[-{Stealth}, thick, popdark] (dong) -- (sport);
\end{tikzpicture}
\legende{Décomposition de 运动 : un élément donne le sens (\zh{辶} mouvement, \zh{力} force), l'autre donne le son (\zh{云} yún).}
\end{popfigure}

\courssub{L'ordre des traits : les règles générales}

\begin{propriete}[Règles d'écriture des traits]
L'écriture chinoise obéit à des règles strictes, à respecter pour chaque caractère :
\begin{enumerate}
\item \cle{Haut avant bas} : on écrit d'abord la partie supérieure (ex. \zh{云} : les deux traits horizontaux avant le trait final).
\item \cle{Gauche avant droite} : la partie gauche d'abord (ex. \zh{体} : \zh{亻} puis \zh{本} ; \zh{健} : \zh{亻} puis \zh{建}).
\item \cle{Horizontal avant vertical} quand deux traits se croisent (ex. \zh{十} dans \zh{本} : horizontal puis vertical).
\item \cle{Extérieur avant intérieur}, puis on « ferme » le caractère (ex. \zh{广} dans \zh{康}).
\item \cle{Exception majeure} : le radical \zh{辶} (mouvement) s'écrit \textbf{en dernier}. Dans \zh{运}, on écrit d'abord \zh{云}, puis \zh{辶} en trois temps : le point, le pli, et la longue vague finale qui « porte » le caractère. Même règle pour \zh{远}, \zh{进}, \zh{这}, \zh{近}.
\end{enumerate}
\end{propriete}

\begin{methode}[Comment mémoriser durablement un caractère]
\begin{enumerate}
\item \cle{Identifier le radical} : il donne le domaine de sens (\zh{亻} = homme, \zh{火} = feu, \zh{辶} = mouvement, \zh{力} = force, \zh{广} = abri).
\item \cle{Identifier le composant phonétique} : il donne une indication de son (\zh{云} yún dans \zh{运} yùn et \zh{动} dòng ; \zh{建} jiàn dans \zh{健} jiàn ; \zh{庚} gēng dans \zh{康} kāng ; \zh{东} dōng dans \zh{炼} liàn ; \zh{本} běn dans \zh{体} tǐ).
\item Se raconter \cle{l'image} du caractère : \zh{休} = un homme \zh{亻} contre un arbre \zh{木} = le repos.
\item \cle{Écrire 5 fois} le caractère en respectant l'ordre des traits, en prononçant le pinyin à voix haute à chaque écriture.
\item Réécrire le caractère \cle{dans un mot} (\zh{运动}, \zh{健康}, \zh{锻炼}) puis dans une phrase complète.
\end{enumerate}
\end{methode}

\courssub{Caractères proches : ne pas confondre}

\begin{attention}[体 (tǐ, corps) ≠ 休 (xiū, se reposer)]
C'est l'erreur la plus fréquente des francophones. Les deux caractères commencent par \zh{亻} (homme), mais :
\begin{itemize}
\item \zh{体} = \zh{亻} + \zh{本} (racine) : le \cle{corps}, la racine de l'homme. Le trait horizontal du bas de \zh{本} est \textbf{court} (il ne dépasse pas).
\item \zh{休} = \zh{亻} + \zh{木} (arbre) : un homme \cle{adossé à un arbre} = le repos. Le trait horizontal du bas de \zh{木} \textbf{dépasse} des deux côtés.
\end{itemize}
Piège en contexte : \zh{身体} (\emph{shēntǐ}, le corps) et \zh{休息} (\emph{xiūxi}, se reposer). Écrire \zh{休体} est une faute grave : le sens devient incompréhensible.
\end{attention}

\begin{attention}[练 (liàn, s'entraîner) ≠ 炼 (liàn, forger) ; 运 (yùn) ≠ 远 (yuǎn)]
\begin{itemize}
\item \zh{练} et \zh{炼} sont \cle{homophones} (tous deux \emph{liàn}, 4e ton) mais pas synonymes : \zh{练} a le radical \zh{纟} (soie → répéter un geste, s'exercer : \zh{练习} \emph{liànxí}, exercices, s'entraîner) ; \zh{炼} a le radical \zh{火} (feu → forger, fondre : \zh{锻炼} \emph{duànliàn}). On dit \zh{锻炼身体} et \textbf{jamais} \zh{锻练身体}.
\item \zh{运} (\emph{yùn}, bouger, 4e ton) et \zh{远} (\emph{yuǎn}, lointain, 3e ton) partagent le radical \zh{辶} : seul l'intérieur change (\zh{云} yún contre \zh{元} yuán). Comparez : \zh{运动} (le sport) et \zh{很远} (\emph{hěn yuǎn}, très loin).
\end{itemize}
\end{attention}

\begin{popfigure}
\begin{tikzpicture}[node distance=0.6cm, every node/.style={font=\small, align=center}]
% Ligne 1: 体 vs 休
\node[draw=popblue, fill=popblueL, rounded corners, minimum width=2.8cm, align=center] (ti){\zh{体} tǐ\\亻 + 本\\le corps};
\node[draw=poporange, fill=poporangeL, rounded corners, minimum width=2.8cm, right=1.0cm of ti, align=center] (xiu){\zh{休} xiū\\亻 + 木\\se reposer};
% Ligne 2: 练 vs 炼
\node[draw=popblue, fill=popblueL, rounded corners, minimum width=2.8cm, below=0.8cm of ti, align=center] (lian1){\zh{练} liàn\\纟 (soie)\\s'exercer};
\node[draw=poporange, fill=poporangeL, rounded corners, minimum width=2.8cm, right=1.0cm of lian1, align=center] (lian2){\zh{炼} liàn\\火 (feu)\\forger};
% Ligne 3: 运 vs 远
\node[draw=popblue, fill=popblueL, rounded corners, minimum width=2.8cm, below=0.8cm of lian1, align=center] (yun2){\zh{运} yùn\\辶 + 云\\bouger};
\node[draw=poporange, fill=poporangeL, rounded corners, minimum width=2.8cm, right=1.0cm of yun2, align=center] (yuan){\zh{远} yuǎn\\辶 + 元\\lointain};
% Flèches doubles
\draw[{Stealth}-{Stealth}, thick, poppink] (ti) -- (xiu);
\draw[{Stealth}-{Stealth}, thick, poppink] (lian1) -- (lian2);
\draw[{Stealth}-{Stealth}, thick, poppink] (yun2) -- (yuan);
\end{tikzpicture}
\legende{Paires de caractères proches à distinguer : même silhouette, sens et radicaux différents.}
\end{popfigure}

\courssub{Exemples traduits et exercice résolu}

\begin{exemplebox}[Deux expressions à réutiliser dans vos rédactions]
\begin{itemize}
\item \zh{身体健康} — \emph{shēntǐ jiànkāng} — « être en bonne santé » (litt. « le corps est sain »). Exemple : \zh{每天运动，身体健康。} (\emph{Měitiān yùndòng, shēntǐ jiànkāng.}) « En faisant du sport chaque jour, on est en bonne santé. »
\item \zh{锻炼身体} — \emph{duànliàn shēntǐ} — « entraîner son corps, faire de l'exercice ». Exemple : \zh{我们应该多锻炼身体。} (\emph{Wǒmen yīnggāi duō duànliàn shēntǐ.}) « Nous devrions faire plus d'exercice. »
\end{itemize}
\end{exemplebox}

\begin{exoresolu}[Corriger les caractères fautifs]
Énoncé : un élève de Première a écrit la phrase suivante pour dire « Le sport rend notre corps plus sain » : \zh{远动让我们的休体更健康。} Relevez les trois erreurs de caractères et réécrivez la phrase correctement.
\tcblower
Solution détaillée :
\begin{enumerate}
\item \zh{远动} est faux : \zh{远} (\emph{yuǎn}) signifie « lointain ». Le mot « sport » s'écrit \zh{运动} avec \zh{云} à l'intérieur du radical \zh{辶}.
\item \zh{休体} est faux : \zh{休} (\emph{xiū}) signifie « se reposer » (\zh{亻} + \zh{木}). Le mot « corps » s'écrit \zh{身体} : \zh{体} = \zh{亻} + \zh{本}.
\item \zh{健康} est correct : \zh{健} (radical \zh{亻}) et \zh{康} (radical \zh{广}) sont bien écrits.
\end{enumerate}
Phrase corrigée : \zh{运动让我们的身体更健康。} (\emph{Yùndòng ràng wǒmen de shēntǐ gèng jiànkāng.}) « Le sport rend notre corps plus sain. »
\end{exoresolu}

\begin{savaistu}[Le radical 辶, la clé du mouvement]
Le radical \zh{辶} (appelé \emph{chuò} ou « clé de la marche ») est une forme abrégée de \zh{辵}, qui représentait des pas sur un chemin. Il indique le \cle{mouvement}, le déplacement. On le retrouve dans de nombreux caractères très fréquents : \zh{运} (\emph{yùn}, bouger, transporter), \zh{远} (\emph{yuǎn}, loin), \zh{近} (\emph{jìn}, proche), \zh{进} (\emph{jìn}, entrer), \zh{还} (\emph{hái}, encore / retourner), \zh{这} (\emph{zhè}, ceci). Astuce : il s'écrit toujours \textbf{en dernier}, comme une route que l'on trace sous le caractère pour le « mettre en marche ».
\end{savaistu}

\begin{aretenir}[L'essentiel de la section]
\begin{itemize}
\item \zh{运动} = \zh{辶} (mouvement) + \zh{云} (son) ; \zh{力} (force) + \zh{云} (son).
\item \zh{健康} : \zh{健} (10 traits, radical \zh{亻}, phonétique \zh{建}) + \zh{康} (11 traits, radical \zh{广}, phonétique \zh{庚}).
\item \zh{体} (7 traits) = \zh{亻} + \zh{本} ; \zh{炼} (9 traits) = \zh{火} + phonétique \zh{东}.
\item Ordre des traits : haut → bas, gauche → droite, horizontal → vertical ; le radical \zh{辶} s'écrit \textbf{en dernier}.
\item À ne pas confondre : \zh{体} (corps) / \zh{休} (repos) ; \zh{练} (s'exercer, radical \zh{纟}) / \zh{炼} (forger, radical \zh{火}) ; \zh{运} (bouger) / \zh{远} (loin).
\item Expressions clés : \zh{身体健康} (être en bonne santé), \zh{锻炼身体} (entraîner son corps).
\end{itemize}
\end{aretenir}

Maintenant que vous savez écrire et distinguer ces caractères, vous êtes prêts pour l'étape suivante : les utiliser dans la traduction du français vers le chinois (\cle{thème}), en veillant à chaque trait.

\coursec{Thème : français → chinois}

Après avoir appris à écrire correctement les caractères du thème du sport (\zh{运动}, \zh{健康}, \zh{跑步}, \zh{身体}…), nous pouvons maintenant passer à l'étape supérieure : le \cle{thème}, c'est-à-dire la traduction du \textbf{français vers le chinois}. C'est l'exercice le plus exigeant de la classe de Première, car il ne suffit pas de connaître les mots : il faut aussi \textbf{réorganiser la phrase} selon la logique chinoise. Le français et le chinois ne placent pas les mots au même endroit, surtout les compléments de temps et les verbes modaux. Cette section vous donne une méthode infaillible en quatre étapes, appliquée à quatre phrases modèles sur l'importance du sport.

\courssub{Observer avant de traduire : quatre phrases modèles}

Lisez d'abord attentivement les quatre paires de phrases ci-dessous. Surlignez mentalement ce qui \emph{change de place} entre le français et le chinois.

\begin{exemplebox}[Quatre phrases modèles à observer]
\textbf{Phrase 1.} Le sport est très important pour la santé.\\
\zh{运动对健康很重要。} \emph{Yùndòng duì jiànkāng hěn zhòngyào.} « Le sport, pour la santé, est très important. »

\textbf{Phrase 2.} Nous devrions faire du sport chaque jour.\\
\zh{我们应该每天运动。} \emph{Wǒmen yīnggāi měi tiān yùndòng.} « Nous devrions chaque jour faire du sport. »

\textbf{Phrase 3.} Parce que le sport est bon pour le corps, je cours tous les matins.\\
\zh{因为运动对身体好，所以我每天早上跑步。} \emph{Yīnwèi yùndòng duì shēntǐ hǎo, suǒyǐ wǒ měi tiān zǎoshang pǎobù.} « Parce que le sport est bon pour le corps, donc je chaque matin cours. »

\textbf{Phrase 4.} Le sport nous rend non seulement en bonne santé, mais aussi heureux.\\
\zh{运动不但让我们健康，而且让我们快乐。} \emph{Yùndòng búdàn ràng wǒmen jiànkāng, érqiě ràng wǒmen kuàilè.} « Le sport non seulement nous rend en bonne santé, mais encore nous rend heureux. »
\end{exemplebox}

\textbf{Observations guidées :}
\begin{itemize}
\item Dans la phrase 1, le français dit « important \emph{pour la santé} » ; le chinois dit \zh{对健康} (\emph{duì jiànkāng}, « pour la santé ») \textbf{avant} l'adjectif \zh{很重要}. La préposition \zh{对} introduit le complément qui précède l'adjectif.
\item Dans la phrase 2, « chaque jour » est à la fin en français ; en chinois, \zh{每天} se place \textbf{avant le verbe}, juste après le verbe modal \zh{应该}.
\item Dans la phrase 3, le chinois utilise volontiers le couple \zh{因为……所以……} (« parce que… donc… »), alors que le français n'emploie qu'un seul connecteur.
\item Dans la phrase 4, le couple « non seulement… mais aussi… » se traduit par \zh{不但……而且……}, et « rendre + adjectif » se traduit par \zh{让 + personne + adjectif}.
\end{itemize}

\courssub{La méthode de traduction en quatre étapes}

Traduire du français vers le chinois ne consiste jamais à remplacer mot à mot. Il faut d'abord \textbf{découper} la phrase française, puis \textbf{reconstruire} la phrase chinoise selon son propre ordre.

\begin{methode}[Réussir le thème : les quatre étapes]
\begin{enumerate}
\item \textbf{Traduire les mots-clés.} Repérez le vocabulaire de la phrase (sport = \zh{运动}, santé = \zh{健康}, corps = \zh{身体}, courir = \zh{跑步}, chaque jour = \zh{每天}) et écrivez-le en chinois sur votre brouillon.
\item \textbf{Placer le sujet.} Identifiez le sujet (qui ? quoi ?) et mettez-le en tête de phrase : \zh{我们}, \zh{运动}, \zh{我}…
\item \textbf{Placer les compléments de temps et de lieu AVANT le verbe.} C'est la grande différence avec le français : \zh{每天}, \zh{早上}, \zh{在学校} se placent entre le sujet et le verbe. Le verbe modal (\zh{应该}, \zh{会}, \zh{能}) se place aussi avant le verbe principal.
\item \textbf{Vérifier les caractères.} Relisez : chaque caractère est-il correct (traits, clés) ? La ponctuation chinoise est-elle pleine chasse （。，） ? Le sens correspond-il exactement à la phrase française ?
\end{enumerate}
\end{methode}

\begin{popfigure}
\begin{tikzpicture}[node distance=0.55cm, every node/.style={font=\small}]
\node[draw=popblue, fill=popblueL, rounded corners, align=center, text width=3.1cm] (fr) {\textbf{Phrase française}\\ « Nous devrions faire du sport chaque jour. »};
\node[draw=poporange, fill=poporangeL, rounded corners, align=center, text width=3.4cm, right=of fr] (an) {\textbf{Analyse}\\ Sujet : nous\\ Modal : devrions\\ Temps : chaque jour\\ Verbe : faire du sport};
\node[draw=poppurple, fill=poppurpleL, rounded corners, align=center, text width=3.4cm, right=of an] (re) {\textbf{Réorganisation}\\ S + modal + temps + V\\ \zh{我们} + \zh{应该} + \zh{每天} + \zh{运动}};
\node[draw=popgreen, fill=popgreenL, rounded corners, align=center, text width=3.1cm, right=of re] (zh) {\textbf{Phrase chinoise}\\ \zh{我们应该每天运动。}};
\draw[-{Stealth}, very thick, popdark] (fr) -- (an);
\draw[-{Stealth}, very thick, popdark] (an) -- (re);
\draw[-{Stealth}, very thick, popdark] (re) -- (zh);
\end{tikzpicture}
\legende{La chaîne de traduction : de la phrase française à la phrase chinoise en passant par l'analyse et la réorganisation.}
\end{popfigure}

\courssub{La règle d'or : la place des compléments de temps}

C'est le point grammatical central de cette section. En français, le complément de temps se met souvent à la fin (« je cours \emph{tous les matins} ») ou au début (« \emph{chaque jour}, nous faisons du sport »). En chinois, la place est fixe et obligatoire.

\begin{propriete}[Placement des compléments de temps en chinois]
\textbf{Structure :} Sujet + complément de temps + (verbe modal) + verbe + (complément).

Le complément de temps (\zh{每天}, \zh{早上}, \zh{现在}, \zh{明天}…) se place \textbf{toujours avant le verbe}, après le sujet. S'il y a un verbe modal (\zh{应该}, \zh{会}, \zh{想}, \zh{能}), l'ordre habituel est : Sujet + modal + temps + verbe, ou Sujet + temps + modal + verbe (les deux sont possibles, mais le temps ne vient jamais après le verbe).

\textbf{Exemple :} \zh{我每天早上跑步。} \emph{Wǒ měi tiān zǎoshang pǎobù.} « Je cours chaque matin. » (litt. « je chaque jour matin cours »).
\end{propriete}

\begin{center}
\begin{tabularx}{\linewidth}{|X|X|X|}
\hline
\rowcolor{popblueL} \textbf{Structure} & \textbf{Exemple (caractères + pinyin)} & \textbf{Traduction} \\
\hline
Sujet + temps + verbe & \zh{我每天运动。} \emph{Wǒ měi tiān yùndòng.} & Je fais du sport chaque jour. \\
\hline
Sujet + modal + temps + verbe & \zh{我们应该每天运动。} \emph{Wǒmen yīnggāi měi tiān yùndòng.} & Nous devrions faire du sport chaque jour. \\
\hline
Sujet + temps + lieu + verbe & \zh{我每天早上在学校跑步。} \emph{Wǒ měi tiān zǎoshang zài xuéxiào pǎobù.} & Je cours chaque matin à l'école. \\
\hline
Sujet + 对 + compl. + adjectif & \zh{运动对健康很重要。} \emph{Yùndòng duì jiànkāng hěn zhòngyào.} & Le sport est très important pour la santé. \\
\hline
因为 … 所以 … & \zh{因为运动对身体好，所以我每天跑步。} \emph{Yīnwèi yùndòng duì shēntǐ hǎo, suǒyǐ wǒ měi tiān pǎobù.} & Parce que le sport est bon pour le corps, je cours chaque jour. \\
\hline
不但 … 而且 … & \zh{运动不但让我们健康，而且让我们快乐。} \emph{Yùndòng búdàn ràng wǒmen jiànkāng, érqiě ràng wǒmen kuàilè.} & Le sport nous rend non seulement en bonne santé, mais aussi heureux. \\
\hline
\end{tabularx}
\end{center}

\courssub{Commentaire détaillé des quatre phrases modèles}

\textbf{Phrase 1 : « Le sport est très important pour la santé. » → \zh{运动对健康很重要。}}

Étape 1, mots-clés : sport \zh{运动}, santé \zh{健康}, important \zh{重要}, très \zh{很}. Étape 2, sujet : \zh{运动}. Étape 3 : le complément « pour la santé » se traduit par \zh{对健康} et se place \textbf{avant} l'adjectif. Étape 4 : on obtient \zh{运动对健康很重要。} Notez que \zh{很} n'exprime pas toujours « très » : il sert souvent de simple lien entre le sujet et l'adjectif.

\textbf{Phrase 2 : « Nous devrions faire du sport chaque jour. » → \zh{我们应该每天运动。}}

Mots-clés : nous \zh{我们}, devoir (conseil) \zh{应该}, faire du sport \zh{运动}, chaque jour \zh{每天}. Ordre chinois : Sujet \zh{我们} + modal \zh{应该} + temps \zh{每天} + verbe \zh{运动}. En français, « chaque jour » est à la fin ; en chinois, il remonte avant le verbe.

\textbf{Phrase 3 : « Parce que le sport est bon pour le corps, je cours tous les matins. » → \zh{因为运动对身体好，所以我每天早上跑步。}}

Le chinois aime les connecteurs en couple : \zh{因为} (parce que) dans la première proposition, \zh{所以} (donc) dans la seconde. En français, « parce que » suffit ; en chinois, garder les deux est la norme à l'écrit. « Tous les matins » = \zh{每天早上}, placé avant le verbe \zh{跑步}.

\textbf{Phrase 4 : « Le sport nous rend non seulement en bonne santé, mais aussi heureux. » → \zh{运动不但让我们健康，而且让我们快乐。}}

« Non seulement… mais aussi… » = \zh{不但……而且……}. « Rendre quelqu'un + adjectif » = \zh{让 + personne + adjectif} : \zh{让我们健康} (nous rend en bonne santé), \zh{让我们快乐} (nous rend heureux). La structure se répète en miroir dans les deux propositions, ce qui donne une phrase élégante et équilibrée, très appréciée en chinois écrit.

\begin{exemplebox}[Exemple supplémentaire traduit pas à pas]
« Faire du sport est bon pour notre corps. »\\
Étape 1 : faire du sport \zh{运动} ; notre corps \zh{我们的身体} ; bon \zh{好} ; pour \zh{对}.\\
Étape 2 : sujet = \zh{运动}.\\
Étape 3 : \zh{对我们的身体} avant l'adjectif \zh{很好}.\\
Étape 4 : \zh{运动对我们的身体很好。} \emph{Yùndòng duì wǒmen de shēntǐ hěn hǎo.} « Faire du sport est très bon pour notre corps. »
\end{exemplebox}

\begin{attention}[L'erreur n°1 des francophones : la place du verbe modal]
Ne dites jamais \zh{我们运动应该} : c'est un calque du français « nous faisons du sport devrions », totalement incorrect. Le verbe modal \zh{应该} se place \textbf{avant le verbe principal} : \zh{我们应该运动。} De même, ne placez jamais le complément de temps après le verbe : on ne dit pas \zh{我跑步每天早上}, mais \zh{我每天早上跑步。} Retenez le réflexe : en chinois, « quand » et « où » se disent \textbf{avant} l'action, jamais après.
\end{attention}

\begin{exoresolu}[Thème guidé : trois phrases à traduire]
\textbf{Énoncé.} Traduisez en chinois : 1) « La santé est très importante pour nous. » 2) « Je devrais courir chaque matin. » 3) « Parce que le sport nous rend heureux, nous aimons le sport. »
\tcblower
\textbf{Solution détaillée.}

1) Mots-clés : santé \zh{健康}, important \zh{重要}, pour nous \zh{对我们}. Sujet : \zh{健康}. Le complément \zh{对我们} précède l'adjectif : \zh{健康对我们很重要。} \emph{Jiànkāng duì wǒmen hěn zhòngyào.}

2) Mots-clés : je \zh{我}, devoir \zh{应该}, chaque matin \zh{每天早上}, courir \zh{跑步}. Ordre : Sujet + modal + temps + verbe : \zh{我应该每天早上跑步。} \emph{Wǒ yīnggāi měi tiān zǎoshang pǎobù.}

3) Connecteurs : \zh{因为……所以……} ; « rendre heureux » = \zh{让我们快乐} ; « aimer le sport » = \zh{喜欢运动}. On obtient : \zh{因为运动让我们快乐，所以我们喜欢运动。} \emph{Yīnwèi yùndòng ràng wǒmen kuàilè, suǒyǐ wǒmen xǐhuan yùndòng.}
\end{exoresolu}

\begin{exercice}[À vous de traduire]
Traduisez en chinois, en appliquant les quatre étapes de la méthode :
\begin{enumerate}
\item « Le sport est bon pour notre santé. »
\item « Nous devrions courir chaque jour à l'école. »
\item « Parce que je fais du sport chaque matin, je suis en très bonne santé. »
\end{enumerate}
\end{exercice}

\begin{corrige}[Exercice : thème]
1) \zh{运动对我们的健康很好。} \emph{Yùndòng duì wǒmen de jiànkāng hěn hǎo.} — Le complément \zh{对我们的健康} précède l'adjectif \zh{很好}.

2) \zh{我们应该每天在学校跑步。} \emph{Wǒmen yīnggāi měi tiān zài xuéxiào pǎobù.} — Ordre : sujet + modal \zh{应该} + temps \zh{每天} + lieu \zh{在学校} + verbe \zh{跑步}.

3) \zh{因为我每天早上运动，所以我很健康。} \emph{Yīnwèi wǒ měi tiān zǎoshang yùndòng, suǒyǐ wǒ hěn jiànkāng.} — Couple de connecteurs \zh{因为……所以……} ; le temps \zh{每天早上} est placé avant le verbe \zh{运动}.
\end{corrige}

\begin{aretenir}[L'essentiel du thème français → chinois]
\begin{itemize}
\item Méthode en quatre étapes : mots-clés → sujet → compléments de temps/lieu avant le verbe → vérification des caractères.
\item Ordre de base : \textbf{Sujet + (modal) + temps + (lieu) + verbe}. Le complément de temps ne se place jamais après le verbe.
\item Le verbe modal \zh{应该} précède toujours le verbe principal : \zh{我们应该运动。}
\item « Pour + nom » devant un adjectif = \zh{对 + nom + adjectif} : \zh{运动对健康很重要。}
\item Connecteurs en couple : \zh{因为……所以……} (parce que… donc…), \zh{不但……而且……} (non seulement… mais aussi…).
\item « Rendre + adjectif » = \zh{让 + personne + adjectif} : \zh{运动让我们快乐。}
\end{itemize}
\end{aretenir}

\coursec{Version : chinois → français}

Dans la section précédente, nous avons appris à traduire du français vers le chinois (le \emph{thème}) : il fallait penser en chinois, choisir les bons mots et respecter l'ordre des mots chinois. Nous faisons maintenant le chemin inverse : la \cle{version}, c'est-à-dire la traduction du chinois vers le français. L'exercice semble plus facile — on comprend « à peu près » le texte — mais il cache un piège redoutable : le \emph{mot à mot}. Un bon traducteur ne transporte pas des mots, il transporte des \emph{idées}, puis les habille à la française.

\courssub{Observer d'abord : le texte de la leçon}

Lisons d'abord le texte sans traduire, juste pour en saisir le sens général. De quoi parle-t-il ? Qui parle ? Quelle est l'idée principale ?

\begin{textebox}[Texte de version — L'importance du sport]
我觉得运动很重要。运动不但对身体好，而且对学习也有帮助。所以，我们学校的学生每天都锻炼。\\
Wǒ juéde yùndòng hěn zhòngyào. Yùndòng búdàn duì shēntǐ hǎo, érqiě duì xuéxí yě yǒu bāngzhù. Suǒyǐ, wǒmen xuéxiào de xuéshēng měitiān dōu duànliàn.\\
« Je pense que le sport est très important. Le sport est non seulement bon pour le corps, mais il aide aussi pour les études. C'est pourquoi les élèves de notre école s'entraînent chaque jour. »
\end{textebox}

Première lecture globale : le locuteur donne son \emph{opinion} (\zh{我觉得} « je pense que »), avance \emph{deux arguments} liés par \zh{不但……而且……} « non seulement... mais aussi », puis tire une \emph{conséquence} introduite par \zh{所以} « c'est pourquoi / donc ». Le texte est donc construit en trois temps : opinion → arguments → conclusion. Repérer cette architecture avant de traduire est la moitié du travail.

\courssub{Vocabulaire du texte}

\begin{center}
\begin{tabularx}{\linewidth}{|X|X|X|X|}
\hline
\rowcolor{popblueL} Caractères & Pinyin & Nature & Traduction \\
\hline
觉得 & juéde & verbe & penser, trouver, estimer \\
\hline
运动 & yùndòng & nom / verbe & sport ; faire du sport \\
\hline
重要 & zhòngyào & adjectif & important \\
\hline
不但……而且…… & búdàn... érqiě... & connecteurs & non seulement... mais aussi... \\
\hline
对 & duì & préposition & envers, pour, à l'égard de \\
\hline
身体 & shēntǐ & nom & corps, santé \\
\hline
学习 & xuéxí & nom / verbe & études ; étudier \\
\hline
有帮助 & yǒu bāngzhù & expression & être utile, aider \\
\hline
所以 & suǒyǐ & conjonction & donc, c'est pourquoi \\
\hline
学校 & xuéxiào & nom & école \\
\hline
学生 & xuéshēng & nom & élève, étudiant \\
\hline
每天 & měitiān & expression temporelle & chaque jour \\
\hline
锻炼 & duànliàn & verbe & s'entraîner, faire de l'exercice \\
\hline
\end{tabularx}
\end{center}

\courssub{Méthode : réussir une version en quatre étapes}

\begin{methode}[La version en quatre étapes]
\begin{enumerate}
\item \textbf{Découper en propositions.} Soulignez les signes de ponctuation chinois (。 ，) et les connecteurs (\zh{不但……而且……}, \zh{所以}, \zh{因为}) : ce sont les articulations du texte. Chaque proposition sera traduite comme un bloc.
\item \textbf{Traduire chaque bloc.} Pour chaque proposition, identifiez le sujet, le verbe, le complément, puis donnez un équivalent français de chaque élément. À ce stade, votre français peut être maladroit : ce n'est pas grave.
\item \textbf{Reformuler en français naturel.} Relisez votre brouillon comme si c'était un texte français : supprimez les répétitions, choisissez de beaux connecteurs (« c'est pourquoi » plutôt que « donc donc »), accordez les verbes et les adjectifs.
\item \textbf{Relire et vérifier.} Vérifiez que rien du texte chinois n'a été oublié (chaque caractère doit avoir un écho en français) et que rien n'a été ajouté. Comptez les phrases : trois phrases en chinois = trois phrases en français.
\end{enumerate}
\end{methode}

\courssub{Traduction guidée, proposition par proposition}

\textbf{Proposition 1 :} \zh{我觉得运动很重要。}

\emph{Wǒ juéde yùndòng hěn zhòngyào.}

Découpage : \zh{我} « je » + \zh{觉得} « je pense / je trouve » + \zh{运动} « le sport » + \zh{很} « très » + \zh{重要} « important ».

Traduction littérale : « Je pense sport très important. »

Traduction naturelle : « Je pense que le sport est très important. »

Remarque : \zh{觉得} introduit une opinion personnelle ; en français, on ajoute la conjonction « que » : \emph{je pense \textbf{que}...}. En chinois, il n'y a pas de mot pour « que » : la subordonnée suit directement.

\textbf{Proposition 2 :} \zh{运动不但对身体好，而且对学习也有帮助。}

\emph{Yùndòng búdàn duì shēntǐ hǎo, érqiě duì xuéxí yě yǒu bāngzhù.}

Découpage : \zh{运动} « le sport » + \zh{不但} « non seulement » + \zh{对身体好} « bon pour le corps » + \zh{而且} « mais aussi » + \zh{对学习也有帮助} « utile aussi pour les études ».

Traduction littérale : « Sport non seulement pour corps bon, mais aussi pour études aussi avoir aide. »

Traduction naturelle : « Le sport est non seulement bon pour le corps, mais il aide aussi pour les études. »

\begin{exemplebox}[La structure 对...有帮助]
La structure \cle{\zh{对} + nom + \zh{有帮助}} signifie « être utile pour... ».\\
\zh{运动对学习有帮助。} \emph{Yùndòng duì xuéxí yǒu bāngzhù.} « Le sport est utile pour les études. »\\
\zh{汉语对我们的未来有帮助。} \emph{Hànyǔ duì wǒmen de wèilái yǒu bāngzhù.} « Le chinois est utile pour notre avenir. »\\
\zh{睡眠对身体有帮助。} \emph{Shuìmián duì shēntǐ yǒu bāngzhù.} « Le sommeil est bon pour le corps. »\\
Attention à l'ordre : en chinois, \zh{对...} se place \emph{avant} le verbe ; en français, « pour... » se place \emph{après} : \zh{对身体好} = « bon \emph{pour le corps} ».
\end{exemplebox}

\textbf{Proposition 3 :} \zh{所以，我们学校的学生每天都锻炼。}

\emph{Suǒyǐ, wǒmen xuéxiào de xuéshēng měitiān dōu duànliàn.}

Découpage : \zh{所以} « c'est pourquoi » + \zh{我们学校的学生} « les élèves de notre école » (\zh{我们} + \zh{学校} + \zh{的} + \zh{学生}) + \zh{每天} « chaque jour » + \zh{都} « tous » + \zh{锻炼} « s'entraînent ».

Traduction littérale : « Donc, notre école de élèves chaque jour tous s'entraînent. »

Traduction naturelle : « C'est pourquoi les élèves de notre école s'entraînent chaque jour. »

Remarque : le petit mot \zh{都} « tous » est souvent absorbé dans le français (« les élèves... s'entraînent » suffit) ; on peut aussi le rendre par « tous » : « s'entraînent tous chaque jour ». Noter la nuance entre \zh{运动} (sport, activité générale) et \zh{锻炼} (exercice physique, entraînement ciblé).

\begin{popfigure}
\begin{tikzpicture}[font=\small]
\node[draw=popblue, fill=popblueL, rounded corners, align=center, text width=4.6cm] (c1) at (0,3) {我觉得运动很重要。};
\node[draw=popblue, fill=popblueL, rounded corners, align=center, text width=4.6cm] (c2) at (0,1.5) {运动不但对身体好，\\而且对学习也有帮助。};
\node[draw=popblue, fill=popblueL, rounded corners, align=center, text width=4.6cm] (c3) at (0,0) {所以，我们学校的学生\\每天都锻炼。};
\node[draw=poporange, fill=poporangeL, rounded corners, align=center, text width=5.2cm] (f1) at (8.6,3) {Je pense que le sport\\est très important.};
\node[draw=poporange, fill=poporangeL, rounded corners, align=center, text width=5.2cm] (f2) at (8.6,1.5) {Le sport est non seulement bon\\pour le corps, mais il aide aussi\\pour les études.};
\node[draw=poporange, fill=poporangeL, rounded corners, align=center, text width=5.2cm] (f3) at (8.6,0) {C'est pourquoi les élèves de\\notre école s'entraînent chaque jour.};
\draw[-{Stealth[length=3mm]}, very thick, popgreen] (c1) -- (f1);
\draw[-{Stealth[length=3mm]}, very thick, popgreen] (c2) -- (f2);
\draw[-{Stealth[length=3mm]}, very thick, popgreen] (c3) -- (f3);
\node[align=center, text=popdark] at (4.3,3.6) {\textbf{CHINOIS} $\longrightarrow$ \textbf{FRANÇAIS}};
\end{tikzpicture}
\legende{Correspondance entre les propositions chinoises et françaises : une phrase chinoise = une phrase française.}
\end{popfigure}

\courssub{Tableau récapitulatif : du chinois au français}

\begin{center}
\begin{tabularx}{\linewidth}{|X|X|X|}
\hline
\rowcolor{popblueL} Structure chinoise & Exemple (caractères + pinyin) & Traduction naturelle \\
\hline
我觉得 + phrase & 我觉得运动很重要。Wǒ juéde yùndòng hěn zhòngyào. & Je pense que le sport est très important. \\
\hline
不但 A，而且 B & 不但对身体好，而且对学习有帮助。búdàn duì shēntǐ hǎo, érqiě duì xuéxí yǒu bāngzhù & non seulement bon pour le corps, mais aussi utile pour les études \\
\hline
对 + nom + 好 & 对身体好 duì shēntǐ hǎo & bon pour le corps \\
\hline
对 + nom + 有帮助 & 对学习有帮助 duì xuéxí yǒu bāngzhù & utile pour les études \\
\hline
所以，... & 所以，我们每天都锻炼。Suǒyǐ, wǒmen měitiān dōu duànliàn. & C'est pourquoi nous nous entraînons chaque jour. \\
\hline
我们学校的学生 & wǒmen xuéxiào de xuéshēng & les élèves de notre école \\
\hline
每天都 + verbe & 每天都锻炼 měitiān dōu duànliàn & s'entraîner chaque jour / tous les jours \\
\hline
\end{tabularx}
\end{center}

\courssub{Le piège du mot à mot}

\begin{attention}[Ne traduisez jamais mot à mot]
L'erreur n°1 des francophones en version est le \textbf{calque} : transporter chaque mot chinois tel quel en français. Résultat : un français bizarre, voire incorrect. Trois pièges précis dans notre texte :
\begin{itemize}
\item \zh{因为……所以……} : en chinois, on peut dire littéralement « parce que... donc... » dans la même phrase. En français, c'est une \textbf{faute} : il faut choisir \emph{l'un des deux}. \zh{因为运动很重要，所以我们每天锻炼。} se traduit « Parce que le sport est important, nous nous entraînons chaque jour » \emph{ou} « Le sport est important, donc nous nous entraînons chaque jour » — jamais « parce que... donc... ».
\item \zh{很} ne signifie pas toujours « très » : devant un adjectif, il sert souvent de simple lien. \zh{运动很重要} peut se traduire « le sport est important » si l'insistance n'est pas forte.
\item \zh{我们学校的学生} : ne dites pas « les élèves de l'école de nous » ; retournez le groupe en français : « les élèves \textbf{de notre école} ». Le chinois empile les compléments avant le nom (avec \zh{的}), le français les place après.
\item \zh{都} placé avant le verbe (ex. \zh{每天都锻炼}) marque la totalité/fréquence ; ne le traduisez pas par « tous » placé n'importe où, mais intégrez-le naturellement : « s'entraînent \textbf{tous les jours} », « \textbf{tous} s'entraînent chaque jour ».
\end{itemize}
\end{attention}

\courssub{Exercice résolu}

\begin{exoresolu}[Version guidée]
Traduisez en français correct et naturel : \zh{因为锻炼对身体好，所以我每天下午都运动。}\\
\emph{Yīnwèi duànliàn duì shēntǐ hǎo, suǒyǐ wǒ měitiān xiàwǔ dōu yùndòng.}
\tcblower
\textbf{Étape 1 — Découpage.} Connecteurs repérés : \zh{因为} « parce que » et \zh{所以} « donc ». Deux propositions : \zh{锻炼对身体好} / \zh{我每天下午都运动}.\\
\textbf{Étape 2 — Traduction par blocs.} \zh{锻炼} « s'entraîner / faire de l'exercice » ; \zh{对身体好} « bon pour le corps » ; \zh{我} « je » ; \zh{每天下午} « chaque après-midi » ; \zh{都运动} « fais du sport ». Brouillon littéral : « Parce que faire de l'exercice pour le corps bon, donc je chaque après-midi fais du sport. »\\
\textbf{Étape 3 — Reformulation.} On supprime l'un des deux connecteurs (règle française) et on réordonne : « Parce que faire de l'exercice est bon pour le corps, je fais du sport chaque après-midi. » ou mieux : « Comme l'exercice est bon pour la santé, je fais du sport tous les après-midis. »\\
\textbf{Étape 4 — Vérification.} Deux propositions en chinois, deux en français ; aucun élément oublié (\zh{每天下午} = « chaque après-midi » est bien rendu). Traduction finale : « Comme l'exercice est bon pour le corps, je fais du sport chaque après-midi. »
\end{exoresolu}

\courssub{Entraînement}

\begin{exercice}[À vous de traduire]
Traduisez en français naturel, en suivant les quatre étapes de la méthode :
\begin{enumerate}
\item \zh{我觉得足球很有意思。}\emph{Wǒ juéde zúqiú hěn yǒu yìsi.}
\item \zh{运动不但对身体好，而且对心情也有帮助。}\emph{Yùndòng búdàn duì shēntǐ hǎo, érqiě duì xīnqíng yě yǒu bāngzhù.}
\item \zh{因为我们学校的学生都喜欢运动，所以每天下午操场上人很多。}\emph{Yīnwèi wǒmen xuéxiào de xuéshēng dōu xǐhuan yùndòng, suǒyǐ měitiān xiàwǔ cāochǎng shàng rén hěn duō.}
\end{enumerate}
\end{exercice}

\begin{corrige}[Exercice « À vous de traduire »]
\begin{enumerate}
\item « Je trouve que le football est très intéressant. » (\zh{有意思} = « intéressant, amusant » ; \zh{觉得} se rend ici naturellement par « trouver ».)
\item « Le sport est non seulement bon pour le corps, mais il est aussi utile pour le moral. » (\zh{心情} = « humeur, moral ».)
\item « Comme les élèves de notre école aiment tous le sport, il y a beaucoup de monde sur le terrain chaque après-midi. » — Attention : un seul connecteur en français (\zh{因为} ou \zh{所以}, pas les deux) ; \zh{操场上人很多} se reformule élégamment par « il y a beaucoup de monde sur le terrain ».
\end{enumerate}
\end{corrige}

\begin{aretenir}[L'essentiel de la version]
\begin{itemize}
\item La version se fait en quatre temps : \cle{découper} → \cle{traduire par blocs} → \cle{reformuler} → \cle{relire}.
\item \zh{我觉得} = « je pense que... » ; \zh{对...有帮助} = « utile pour... » ; \zh{不但……而且……} = « non seulement... mais aussi... » ; \zh{所以} = « c'est pourquoi ».
\item Interdiction absolue : le mot à mot. Le français final doit être \emph{correct et naturel}, comme si le texte avait été écrit directement en français.
\item \zh{因为……所以……} : en français, on choisit \emph{un seul} connecteur.
\item L'ordre du groupe nominal \zh{的} (ex. \zh{我们学校的学生}) s'inverse en français : « les élèves de notre école ».
\end{itemize}
\end{aretenir}

\coursec{Production guidée et grille d'évaluation}

Après avoir traduit du chinois vers le français (version), il est temps de franchir l'étape la plus exigeante et la plus valorisante à l'examen : \cle{produire} soi-même un texte en chinois. Cette dernière section vous donne un sujet précis, un plan imposé, une méthode de rédaction en trois étapes et une grille d'évaluation sur 20 points, identique dans son esprit à celle des épreuves de chinois LV2.

\courssub{Le sujet de production}

\begin{definition}[Sujet de production]
Rédigez un texte de \textbf{5 à 8 phrases} en chinois sur le thème : \zh{运动很重要} (Yùndòng hěn zhòngyào) — « Le sport est très important ». Vous devez utiliser le plan imposé, les connecteurs indiqués et la banque de mots obligatoires.
\end{definition}

\textbf{Plan imposé (5 étapes) :}

\begin{enumerate}
\item \textbf{Thèse} : commencez par \zh{我觉得...} (Wǒ juéde...) — « Je pense que... »
\item \textbf{Argument 1} : utilisez la structure de cause \zh{因为...所以...} (yīnwèi... suǒyǐ...) — « parce que... donc... »
\item \textbf{Argument 2} : utilisez la structure d'addition \zh{不但...而且...} (búdàn... érqiě...) — « non seulement... mais aussi... »
\item \textbf{Exemple personnel} : commencez par \zh{我每天都...} (Wǒ měitiān dōu...) — « Chaque jour, je... »
\item \textbf{Conclusion} : commencez par \zh{总之，...} (Zǒngzhī, ...) — « En résumé, ... »
\end{enumerate}

\textbf{Banque de mots obligatoires} (les six mots doivent apparaître dans votre texte) :

\begin{center}
\begin{tabularx}{\linewidth}{|X|X|X|X|}
\hline
\rowcolor{popblueL} Caractères & Pinyin & Nature & Traduction \\
\hline
运动 & yùndòng & nom / verbe & sport ; faire du sport \\
\hline
健康 & jiànkāng & adjectif / nom & sain ; la santé \\
\hline
身体 & shēntǐ & nom & le corps, la santé physique \\
\hline
锻炼 & duànliàn & verbe / nom & s'entraîner ; l'exercice physique \\
\hline
重要 & zhòngyào & adjectif & important \\
\hline
应该 & yīnggāi & verbe modal & devoir (il faut, on devrait) \\
\hline
\end{tabularx}
\end{center}

\courssub{Méthode de rédaction en trois étapes}

\begin{methode}[Comment rédiger votre texte en trois étapes]
\begin{enumerate}
\item \textbf{Écrire d'abord le plan en connecteurs.} Sur votre brouillon, écrivez uniquement le squelette du texte : \zh{我觉得... / 因为...所以... / 不但...而且... / 我每天都... / 总之，...}. Vous obtenez ainsi la colonne vertébrale de votre production : plus de risque d'oublier une étape du plan.
\item \textbf{Compléter chaque phrase une par une.} Pour chaque connecteur, rédigez une phrase courte et correcte (Sujet + Verbe + Complément). Vérifiez que chaque phrase contient au moins un mot de la banque obligatoire. Ne cherchez pas des phrases compliquées : en chinois, une phrase simple et juste vaut mieux qu'une phrase longue et fausse.
\item \textbf{Vérifier les caractères un à un.} Relisez votre texte en regardant chaque caractère : tous les traits sont-ils écrits ? Le radical est-il correct ? N'avez-vous pas confondu deux caractères proches (\zh{练} / \zh{炼}, \zh{休} / \zh{体}) ? C'est cette dernière relecture qui fait gagner des points.
\end{enumerate}
\end{methode}

\begin{propriete}[Critères de réussite au Bac et à l'examen blanc]
Une production de niveau Première est réussie si elle réunit trois conditions :
\begin{itemize}
\item des \textbf{caractères corrects} : traits complets, radicaux exacts, aucun caractère inventé ;
\item des \textbf{connecteurs variés} : au moins trois connecteurs différents parmi \zh{因为...所以...}, \zh{不但...而且...}, \zh{我每天都...}, \zh{总之} ;
\item \textbf{au moins une structure complexe} : \zh{因为...所以...} ou \zh{不但...而且...} correctement employée (les deux moitiés présentes, dans le bon ordre).
\end{itemize}
Un texte de 5 phrases parfaitement correct obtient une meilleure note qu'un texte de 10 phrases plein de fautes.
\end{propriete}

\courssub{Production type commentée}

Voici une production type, rédigée selon le plan imposé. Observez d'abord le texte, puis lisez les commentaires.

\begin{textebox}[运动很重要 — Production type]
我觉得运动很重要。\\
Wǒ juéde yùndòng hěn zhòngyào.\\
Je pense que le sport est très important.\\[4pt]
因为运动对身体好，所以我们应该每天锻炼。\\
Yīnwèi yùndòng duì shēntǐ hǎo, suǒyǐ wǒmen yīnggāi měitiān duànliàn.\\
Parce que le sport est bon pour le corps, nous devrions nous entraîner chaque jour.\\[4pt]
运动不但让我们健康，而且让我们快乐。\\
Yùndòng búdàn ràng wǒmen jiànkāng, érqiě ràng wǒmen kuàilè.\\
Le sport non seulement nous rend en bonne santé, mais aussi nous rend heureux.\\[4pt]
我每天都和朋友们一起踢足球。\\
Wǒ měitiān dōu hé péngyoumen yìqǐ tī zúqiú.\\
Chaque jour, je joue au football avec mes amis.\\[4pt]
总之，运动对我们很重要，我们应该多运动。\\
Zǒngzhī, yùndòng duì wǒmen hěn zhòngyào, wǒmen yīnggāi duō yùndòng.\\
En résumé, le sport est très important pour nous, nous devrions faire plus de sport.
\end{textebox}

\textbf{Commentaires sur la production type :}
\begin{itemize}
\item La phrase 1 pose la thèse avec \zh{我觉得} + adjectif \zh{重要} : structure Sujet + 很 + adjectif.
\item La phrase 2 utilise \zh{因为...所以...} : la cause d'abord, la conséquence ensuite, comme en français. Notez \zh{对身体好} (bon pour le corps) avec la préposition \zh{对}.
\item La phrase 3 utilise \zh{不但...而且...} : les deux moitiés encadrent la structure \zh{让我们} + adjectif (« nous rendre... »).
\item La phrase 4 donne l'exemple personnel : \zh{每天都} + activité concrète \zh{踢足球} (jouer au football), très naturelle au Cameroun.
\item La phrase 5 conclut avec \zh{总之} et reprend deux mots obligatoires : \zh{重要} et \zh{应该}.
\end{itemize}

\begin{exemplebox}[Conclusion modèle à réutiliser]
\zh{总之，运动对我们很重要，我们应该多运动。}\\
\emph{Zǒngzhī, yùndòng duì wǒmen hěn zhòngyào, wǒmen yīnggāi duō yùndòng.}\\
« En résumé, le sport est très important pour nous, nous devrions faire plus de sport. »\\[4pt]
Cette conclusion est un modèle : elle résume la thèse (\zh{很重要}) et formule un conseil (\zh{应该多运动}). Vous pouvez l'adapter : \zh{总之，健康很重要，我们应该多锻炼。} (En résumé, la santé est importante, nous devrions nous entraîner davantage.)
\end{exemplebox}

\begin{attention}[Erreurs fréquentes des francophones dans l'écriture des caractères]
\begin{itemize}
\item \textbf{Oublier le trait final du radical 辶 dans \zh{运}} : la clé du mouvement 辶 s'écrit en trois temps (point, virgule brisée, puis le long trait final ondulé qui « porte » le caractère). Sans ce dernier trait, le caractère est faux. Écrivez d'abord \zh{云}, puis la clé 辶 en dernier.
\item \textbf{Confondre \zh{练} et \zh{炼}} : \zh{锻炼} (s'entraîner) s'écrit avec \zh{炼} (clé du feu 火) et non avec \zh{练} (clé de la soie 纟, comme dans \zh{练习}, s'exercer). Les deux se prononcent \emph{liàn} (4e ton) : c'est un homophone, seul le radical les distingue.
\item \textbf{Oublier le trait horizontal dans \zh{体}} : \zh{身体} s'écrit avec \zh{体} (亻+ 本) ; ne pas confondre avec \zh{休} (se reposer, 亻+ 木).
\item \textbf{Placer le complément de temps après le verbe} : on écrit \zh{我每天都运动} et jamais « 我运动每天 ». En chinois, le complément de temps se place avant le verbe.
\end{itemize}
\end{attention}

\courssub{Grille d'évaluation sur 20 points}

Votre production sera notée selon cinq critères. Étudiez cette grille avant de rédiger : elle vous indique exactement où se gagnent les points.

\begin{popfigure}
\begin{tikzpicture}[x=1cm,y=1cm]
\fill[popblueL] (0,0) rectangle (10.5,0.8);
\node[anchor=west,font=\bfseries] at (0.2,0.4) {Critère};
\node[anchor=west,font=\bfseries] at (7.6,0.4) {Barème};
\fill[popcream] (0,-0.8) rectangle (10.5,0);
\node[anchor=west] at (0.2,-0.4) {Exactitude des caractères (traits, radicaux)};
\node[anchor=west] at (7.6,-0.4) {5 points};
\fill[white] (0,-1.6) rectangle (10.5,-0.8);
\node[anchor=west] at (0.2,-1.2) {Grammaire et ordre des mots};
\node[anchor=west] at (7.6,-1.2) {5 points};
\fill[popcream] (0,-2.4) rectangle (10.5,-1.6);
\node[anchor=west] at (0.2,-2.0) {Connecteurs (因为...所以..., 不但...而且...)};
\node[anchor=west] at (7.6,-2.0) {4 points};
\fill[white] (0,-3.2) rectangle (10.5,-2.4);
\node[anchor=west] at (0.2,-2.8) {Respect du plan imposé (5 étapes)};
\node[anchor=west] at (7.6,-2.8) {3 points};
\fill[popcream] (0,-4.0) rectangle (10.5,-3.2);
\node[anchor=west] at (0.2,-3.6) {Richesse du vocabulaire (banque de mots)};
\node[anchor=west] at (7.6,-3.6) {3 points};
\draw[popline,thick] (0,0.8) -- (10.5,0.8);
\draw[popline] (0,0) -- (10.5,0) (0,-0.8) -- (10.5,-0.8) (0,-1.6) -- (10.5,-1.6) (0,-2.4) -- (10.5,-2.4) (0,-3.2) -- (10.5,-3.2) (0,-4.0) -- (10.5,-4.0);
\draw[popline,thick] (0,0.8) -- (0,-4.0) (10.5,0.8) -- (10.5,-4.0) (7.4,0.8) -- (7.4,-4.0);
\fill[poporangeL] (0,-4.8) rectangle (10.5,-4.0);
\node[anchor=west,font=\bfseries] at (0.2,-4.4) {Total};
\node[anchor=west,font=\bfseries] at (7.6,-4.4) {20 points};
\draw[popline,thick] (0,-4.8) -- (10.5,-4.8) (0,-4.8) -- (0,-4.0) (10.5,-4.8) -- (10.5,-4.0);
\end{tikzpicture}
\legende{Grille d'évaluation de la production écrite sur 20 points}
\end{popfigure}

\begin{center}
\begin{tabularx}{\linewidth}{|X|X|X|}
\hline
\rowcolor{popblueL} Critère & Ce qui est vérifié & Erreur qui fait perdre des points \\
\hline
Caractères (5) & traits complets, radicaux exacts & \zh{运} sans le trait final de 辶 ; \zh{练} au lieu de \zh{炼} \\
\hline
Grammaire (5) & ordre Sujet + Verbe + Complément & complément de temps après le verbe \\
\hline
Connecteurs (4) & au moins 3 connecteurs variés & \zh{因为} sans \zh{所以} (structure incomplète) \\
\hline
Plan (3) & les 5 étapes dans l'ordre & exemple personnel avant les arguments \\
\hline
Vocabulaire (3) & les 6 mots obligatoires présents & moins de 4 mots de la banque utilisés \\
\hline
\end{tabularx}
\end{center}

\courssub{Auto-évaluation puis correction collective}

\begin{methode}[Auto-évaluer sa production avant de la rendre]
\begin{enumerate}
\item Comptez vos phrases : entre 5 et 8 ? Sinon, ajustez.
\item Surlignez les 6 mots obligatoires (\zh{运动}, \zh{健康}, \zh{身体}, \zh{锻炼}, \zh{重要}, \zh{应该}) : en manque-t-il un ?
\item Entourez vos connecteurs : avez-vous bien \zh{我觉得}, \zh{因为...所以...}, \zh{不但...而且...}, \zh{我每天都...}, \zh{总之} ?
\item Relisez chaque caractère un à un, en particulier \zh{运}, \zh{锻炼}, \zh{健康}, \zh{体}.
\item Attribuez-vous une note sur 20 avec la grille, puis échangez votre texte avec un voisin pour une double correction.
\end{enumerate}
\end{methode}

\begin{exoresolu}[Correction collective d'une production d'élève]
Énoncé : voici la production d'un élève de Première A. Corrigez-la et notez-la sur 20.

\zh{我觉得运动很重要。}\\
\zh{因为运动好，所以我们身体。}\\
\zh{我不但运动，而且健康。}\\
\zh{我每天运动都。}\\
\zh{总之，我们应该多运动。}

\tcblower

Solution détaillée :
\begin{itemize}
\item Phrase 1 : correcte. \zh{我觉得运动很重要。} (1 point de plan pour la thèse.)
\item Phrase 2 : \zh{因为运动好，所以我们身体。} est \textbf{fausse} : \zh{身体} est un nom, pas un verbe. Correction : \zh{因为运动对身体好，所以我们很健康。} (Parce que le sport est bon pour le corps, nous sommes en bonne santé.)
\item Phrase 3 : \zh{我不但运动，而且健康。} est maladroite : \zh{不但...而且...} doit encadrer deux éléments parallèles. Correction : \zh{运动不但让身体好，而且让我们健康。}
\item Phrase 4 : \zh{我每天运动都。} est \textbf{fausse} : \zh{都} doit précéder le verbe. Correction : \zh{我每天都运动。} ou \zh{我每天都踢足球。}
\item Phrase 5 : correcte, avec \zh{总之} et \zh{应该}.
\end{itemize}
Notation : caractères 5/5 (aucun caractère faux), grammaire 2/5 (deux phrases fautives), connecteurs 3/4 (\zh{不但...而且...} mal employé), plan 3/3 (5 étapes respectées), vocabulaire 2/3 (\zh{锻炼} absent). Total : \textbf{15/20}.
\end{exoresolu}

\begin{experience}[Atelier d'écriture en classe]
\begin{enumerate}
\item \textbf{Rédaction individuelle (20 min)} : chaque élève rédige son texte \zh{运动很重要} en suivant la méthode en trois étapes.
\item \textbf{Auto-évaluation (5 min)} : chacun se note avec la grille sur 20.
\item \textbf{Correction croisée (10 min)} : échange des copies par binômes ; le correcteur souligne les erreurs de caractères en rouge et les erreurs de connecteurs en bleu.
\item \textbf{Correction collective (10 min)} : le professeur corrige au tableau une production type volontairement fautive ; la classe propose les corrections phrase par phrase.
\end{enumerate}
\end{experience}

\begin{aretenir}[L'essentiel de la production guidée]
\begin{itemize}
\item Sujet : \zh{运动很重要}, 5 à 8 phrases, plan en 5 étapes : \zh{我觉得} → \zh{因为...所以...} → \zh{不但...而且...} → \zh{我每天都...} → \zh{总之}.
\item Méthode : d'abord le squelette en connecteurs, puis les phrases, enfin la vérification des caractères un à un.
\item Mots obligatoires : \zh{运动}, \zh{健康}, \zh{身体}, \zh{锻炼}, \zh{重要}, \zh{应该}.
\item Vigilance graphique : trait final de 辶 dans \zh{运}, clé du feu 火 dans \zh{炼}.
\item Grille sur 20 : caractères 5, grammaire 5, connecteurs 4, plan 3, vocabulaire 3.
\end{itemize}
\end{aretenir}

\begin{exercice}[À vous d'écrire]
Rédigez votre propre texte \zh{运动很重要} (5 à 8 phrases) en respectant le plan imposé et en utilisant les six mots obligatoires. Ajoutez au moins un élément personnel camerounais (le football au quartier, la course à l'école, les sports du lycée). Notez-vous ensuite avec la grille sur 20.
\end{exercice}

\begin{corrige}[Exercice — Production guidée]
Exemple de production attendue :

\zh{我觉得运动很重要。因为运动对身体好，所以我们应该每天锻炼。运动不但让我们健康，而且让我们快乐。我每天都和朋友们在杜阿拉踢足球。总之，运动对我们很重要，我们应该多运动。}

\emph{Wǒ juéde yùndòng hěn zhòngyào. Yīnwèi yùndòng duì shēntǐ hǎo, suǒyǐ wǒmen yīnggāi měitiān duànliàn. Yùndòng búdàn ràng wǒmen jiànkāng, érqiě ràng wǒmen kuàilè. Wǒ měitiān dōu hé péngyoumen zài Dùālā tī zúqiú. Zǒngzhī, yùndòng duì wǒmen hěn zhòngyào, wǒmen yīnggāi duō yùndòng.}

« Je pense que le sport est très important. Parce que le sport est bon pour le corps, nous devrions nous entraîner chaque jour. Le sport non seulement nous rend en bonne santé, mais aussi nous rend heureux. Chaque jour, je joue au football avec mes amis à Douala. En résumé, le sport est très important pour nous, nous devrions faire plus de sport. »

Vérification : 6 phrases, 5 étapes du plan respectées, 6 mots obligatoires présents, deux structures complexes (\zh{因为...所以...} et \zh{不但...而且...}), élément personnel camerounais (\zh{在杜阿拉踢足球}). Note visée : 18 à 20/20 si les caractères sont exacts.
\end{corrige}

\newpage

\coursec{Activité d'intégration}

\coursec{Leçon 12 — Expression écrite : caractères liés à l'importance de la pratique du sport}

\courssub{Objectifs de la leçon}

À la fin de cette leçon, tu seras capable de :

\begin{itemize}
\item produire un \cle{texte argumentatif simple} en chinois (affiche, message, court paragraphe) sur l'importance de la pratique du sport ;
\item réutiliser correctement le \cle{vocabulaire du thème} : \zh{运动} (yùndòng), \zh{健康} (jiànkāng), \zh{身体} (shēntǐ), \zh{重要} (zhòngyào), \zh{锻炼} (duànliàn), \zh{好处} (hǎochù), \zh{跑步} (pǎobù), \zh{精神} (jīngshén) ;
\item employer les \cle{connecteurs logiques} \zh{因为……所以……} (yīnwèi… suǒyǐ…), \zh{不但……而且……} (búdàn… érqiě…) et \zh{总之} (zǒngzhī) pour structurer ton raisonnement ;
\item écrire \cle{sans faute} les caractères du thème, en maîtrisant l'ordre des traits, les radicaux (clés) et en distinguant les caractères graphiquement proches (\zh{己}/\zh{已}, \zh{休}/\zh{体}, \zh{远}/\zh{运}) ;
\item réaliser des exercices de \cle{thème} (français → chinois) et de \cle{version} (chinois → français) sur des phrases simples du thème ;
\item évaluer et corriger une production écrite à l'aide d'une \cle{grille d'évaluation} critériée.
\end{itemize}

\courssub{Situation déclenchante : La Semaine du sport au lycée}

\begin{textebox}[Annonce officielle du lycée — 学校通知]
学校通知\\
\textit{Xuéxiào tōngzhī.}\\
« Annonce de l'école »\\[4pt]
下周是我们学校的体育周。\\
\textit{Xià zhōu shì wǒmen xuéxiào de tǐyù zhōu.}\\
« La semaine prochaine, c'est la semaine du sport de notre école. »\\[4pt]
请每个小组做一张海报：写一个口号，写三个理由，最后写一个结论。\\
\textit{Qǐng měi ge xiǎozǔ zuò yì zhāng hǎibào : xiě yí ge kǒuhào, xiě sān ge lǐyóu, zuìhòu xiě yí ge jiélùn.}\\
« Chaque groupe doit faire une affiche : écrire un slogan, trois arguments, et enfin une conclusion. »\\[4pt]
请用正确的汉字！\\
\textit{Qǐng yòng zhèngquè de Hànzì !}\\
« Utilisez des caractères chinois corrects, s'il vous plaît ! »
\end{textebox}

\begin{experience}[Analyse de la situation — Travail en groupe de 3]
\begin{enumerate}
\item \textbf{Compréhension.} Réponds en français : (a) Quel événement est annoncé ? (b) Que doit contenir chaque affiche ? (c) Quelle exigence porte sur les caractères ? Souligne dans le texte : \zh{海报} (hǎibào, affiche), \zh{口号} (kǒuhào, slogan), \zh{理由} (lǐyóu, argument), \zh{结论} (jiélùn, conclusion).
\item \textbf{Contexte.} Une délégation d'élèves de l'école partenaire de \zh{北京} (Běijīng, Pékin) sera présente. L'affiche doit être compréhensible, correcte et motivante.
\end{enumerate}
\end{experience}

\courssub{Modèle de texte commenté : structure d'une affiche argumentative}

\begin{methode}[Rédiger une affiche « Sport et Santé » en 5 étapes]
\begin{enumerate}
\item \textbf{Le slogan (口号)} : Court, percutant, structure \cle{Sujet + 很 + Adjectif}. Exemple : \zh{运动很重要！} (Yùndòng hěn zhòngyào !)
\item \textbf{Argument 1 — Cause / Conséquence} : Utilise \cle{因为……所以……}. Exemple : \zh{因为运动对身体好，所以我们要天天锻炼。}
\item \textbf{Argument 2 — Accumulation d'avantages} : Utilise \cle{不但……而且……}. Exemple : \zh{运动不但快乐，而且强身健体。}
\item \textbf{Argument 3 — Structure libre (but, possibilité)} : \zh{可以} (kěyǐ), \zh{能够} (nénggòu), \zh{有助于} (yǒuzhùyú). Exemple : \zh{运动可以提高学习效率。}
\item \textbf{La conclusion (结论)} : Ouvre avec \cle{总之}, termine par une incitation (\zh{吧} ba, \zh{一起} yìqǐ). Exemple : \zh{总之，为了健康，大家一起运动吧！}
\end{enumerate}
\end{methode}

\begin{exoresolu}[Affiche modèle commentée — « Bougeons pour notre santé ! » 动起来，为了健康！]
\textbf{Production complète :}\\[4pt]
\zh{运动很重要！}\\
\textit{Yùndòng hěn zhòngyào !}\\
« Le sport, c'est important ! »\\[6pt]
\zh{因为运动对我们的身体好，所以我们应该每天锻炼。}\\
\textit{Yīnwèi yùndòng duì wǒmen de shēntǐ hǎo, suǒyǐ wǒmen yīnggāi měitiān duànliàn.}\\
« Parce que le sport est bon pour notre corps, nous devrions nous entraîner chaque jour. »\\[6pt]
\zh{运动不但很有意思，而且让我们更健康。}\\
\textit{Yùndòng búdàn hěn yǒu yìsi, érqiě ràng wǒmen gèng jiànkāng.}\\
« Le sport n'est pas seulement amusant, il nous rend aussi plus en forme. »\\[6pt]
\zh{运动还可以帮助我们学习：运动后，精神更好。}\\
\textit{Yùndòng hái kěyǐ bāngzhù wǒmen xuéxí : yùndòng hòu, jīngshén gèng hǎo.}\\
« Le sport peut aussi aider nos études : après le sport, l'esprit est meilleur. »\\[6pt]
\zh{总之，为了我们的健康，大家一起运动吧！}\\
\textit{Zǒngzhī, wèile wǒmen de jiànkāng, dàjiā yìqǐ yùndòng ba !}\\
« En résumé, pour notre santé, faisons tous du sport ensemble ! »\\[6pt]
\textbf{Analyse linguistique :}
\begin{itemize}
\item \textbf{Slogan} : Respect de l'ordre S + 很 + Adj. Ponctuation \zh{！} (pleine chasse).
\item \textbf{Arg. 1} : \zh{因为……所以……} complet. \zh{应该} (yīnggāi) exprime l'obligation morale/conseil. \zh{对……好} (duì… hǎo) = « bon pour… ».
\item \textbf{Arg. 2} : \zh{不但……而且……} lie deux qualités. \zh{让} (ràng) + objet + adj. = « rendre qqn + adj. ».
\item \textbf{Arg. 3} : Variation syntaxique. \zh{还可以} (hái kěyǐ) « peut aussi ». \zh{帮助} (bāngzhù) + objet + verbe. Deux-points \zh{：} pour expliciter.
\item \textbf{Conclusion} : \zh{总之} (zǒngzhī) résume. \zh{为了} (wèile) « dans le but de ». \zh{吧} (ba) suggestion collective.
\end{itemize}
\end{exoresolu}

\courssub{Structures et connecteurs utiles — Tableau de référence}

\begin{propriete}[Connecteurs logiques pour l'argumentation (Niveau HSK 2-3)]
\begin{center}
\begin{tabularx}{\linewidth}{|X|X|X|}
\hline
\rowcolor{popblueL} Structure & Exemple (Caractères + Pinyin) & Traduction / Emploi \\
\hline
S + 很 + Adj. & \zh{运动很重要。} (Yùndòng hěn zhòngyào.) & Affirmation simple. \textbf{Attention :} pas de verbe « être » (是), 很 lie sujet et adjectif. \\
\hline
因为 A，所以 B & \zh{因为下雨，所以我不去跑步。} (Yīnwèi xiàyǔ, suǒyǐ wǒ bù qù pǎobù.) & Cause (A) → Conséquence (B). En chinois écrit, les deux termes sont souvent présents. \\
\hline
S + 不但 A，而且 B & \zh{他不但会打篮球，而且会游泳。} (Tā búdàn huì dǎ lánqiú, érqiě huì yóuyǒng.) & Addition d'idées (A et B sont vrais). A et B ont même fonction syntaxique (verbes, adjectifs...). \\
\hline
总之，S + V. & \zh{总之，健康最重要。} (Zǒngzhī, jiànkāng zuì zhòngyào.) & Marque la conclusion / synthèse. Place en début de phrase, suivi d'une virgule \zh{，}. \\
\hline
为了 + But, S + V. & \zh{为了健康，我每天跑步。} (Wèile jiànkāng, wǒ měitiān pǎobù.) & Structure de finalité (but antéposé). Très courant dans les exhortations. \\
\hline
\end{tabularx}
\end{center}
\end{propriete}

\courssub{Écriture correcte des caractères du thème : Traits, Radicaux, Caractères proches}

\begin{propriete}[Analyse graphique des caractères clés du thème]
\begin{center}
\begin{tabularx}{\linewidth}{|X|X|X|X|X|X|}
\hline
\rowcolor{popblueL} Caractère & Pinyin & Radical (Clé) & Nb traits & Ordre des traits (règles clés) & Pièges / Caractères proches \\
\hline
\zh{运} & yùn & \zh{辶} (chuò, marche) & 7 & 1. \zh{云} (yún) : horizontales, verticale, crochet. 2. \zh{辶} : point, horizontale, pliure finale. & \textbf{Confusion \zh{远} (yuǎn)} : intérieur \zh{元} (yuán) vs \zh{云} (yún). \zh{远} = loin ; \zh{运} = mouvement/transporter. \\
\hline
\zh{动} & dòng & \zh{力} (lì, force) & 6 & 1. \zh{云} (simplifié haut). 2. \zh{力} : horizontale, crochet vertical, point. & Composé de \zh{运} + \zh{动} = \zh{运动}. \zh{动} = bouger, action. \\
\hline
\zh{体} & tǐ & \zh{亻} (rén, homme) & 7 & 1. \zh{亻} : trait oblique, verticale. 2. \zh{本} (běn) : horizontale, verticale, horizontales. & \textbf{Confusion \zh{休} (xiū)} : \zh{休} = homme + arbre (repos) ; \zh{体} = homme + racine/origine (corps). \zh{体育} (tǐyù), \zh{身体} (shēntǐ). \\
\hline
\zh{健} & jiàn & \zh{亻} (rén, homme) & 10 & 1. \zh{亻}. 2. \zh{建} (jiàn, construire) : haut \zh{廿}, bas \zh{攵} (pū, frapper). & \zh{健康} (jiànkāng). Phonétique \zh{建}. Sens : homme fort/bâti. \\
\hline
\zh{康} & kāng & \zh{广} (guǎng, toit) & 11 & 1. \zh{广} : point, horizontale, crochet. 2. \zh{米} (mǐ, riz). 3. \zh{甘} (gān, doux) : horizontales, verticale, point. & \zh{康} = toit + nourriture abondante → bien-être. Ne pas confondre avec \zh{庆} (qìng, célébrer). \\
\hline
\zh{锻} & duàn & \zh{钅} (jīn, métal) & 12 & 1. \zh{钅} (5 traits : oblique, horizontale, verticale, point, horizontale). 2. \zh{段} (duàn, segment). & \zh{锻炼} (duànliàn). Sens originel : forger le métal → s'entraîner dur. \\
\hline
\zh{炼} & liàn & \zh{火} (huǒ, feu) & 9 & 1. \zh{火} (4 traits). 2. \zh{东} (dōng, est) : horizontales, verticale, point, horizontale. & \zh{炼} = feu + est (phonétique) → raffiner, forger. \zh{锻炼} = forger le corps/l'esprit. \\
\hline
\end{tabularx}
\end{center}
\end{propriete}

\begin{attention}[Erreurs graphiques fréquentes des francophones]
\begin{itemize}
\item \textbf{Ordre des traits} : Écrire \zh{辶} (clé marche) en 1 seul trait continu (point → horizontale → pliure) et \textbf{en dernier} (règle : enveloppe se ferme après le contenu).
\item \textbf{Radical \zh{亻} (rén)} : Le trait oblique part du haut vers le bas gauche ; la verticale ne touche pas le trait oblique au départ.
\item \textbf{\zh{己} (jǐ, soi) vs \zh{已} (yǐ, déjà)} : \zh{己} 3e trait ne dépasse pas la verticale (boucle fermée) ; \zh{已} 3e trait dépasse vers la droite (boucle ouverte). Présent dans \zh{起} (qǐ) vs \zh{以} (yǐ).
\item \textbf{Ponctuation} : N'utilise \textbf{jamais} la ponctuation française (, . ; : ! ?) dans un texte chinois. Utilise \zh{，} \zh{。} \zh{；} \zh{：} \zh{！} \zh{？} \zh{、} \zh{«} \zh{»}.
\end{itemize}
\end{attention}

\begin{experience}[Atelier d'écriture guidée — 15 min]
\begin{enumerate}
\item Sur papier quadrillé (格子纸), écris 5 fois chaque caractère du tableau ci-dessus en respectant l'ordre des traits.
\item Entoure en rouge le radical (clé) de chaque caractère.
\item Écris les paires de confusions : \zh{运 / 远}, \zh{体 / 休}, \zh{己 / 已}, \zh{锻 / 短} (duǎn, court). Explique à l'oral la différence de composant intérieur.
\end{enumerate}
\end{experience}

\courssub{Vocabulaire thématique — Sport, Santé, Bien-être}

\begin{center}
\begin{tabularx}{\linewidth}{|X|X|X|X|}
\hline
\rowcolor{popblueL} Caractères & Pinyin & Nature & Traduction \\
\hline
运动 & yùndòng & nom / verbe & sport, exercice physique ; faire du sport \\
\hline
体育 & tǐyù & nom & éducation physique, sport (scolaire) \\
\hline
健康 & jiànkāng & nom / adj. & santé ; sain, en bonne santé \\
\hline
身体 & shēntǐ & nom & corps, condition physique \\
\hline
重要 & zhòngyào & adjectif & important \\
\hline
锻炼 & duànliàn & verbe & s'entraîner, se forger (corps, volonté) \\
\hline
好处 & hǎochù & nom & avantage, bienfait, côté positif \\
\hline
跑步 & pǎobù & verbe & courir, faire du jogging \\
\hline
打球 & dǎqiú & verbe & jouer au ballon (basket, foot, volley...) \\
\hline
游泳 & yóuyǒng & verbe & nager \\
\hline
精神 & jīngshén & nom / adj. & esprit, moral, vitalité ; en forme, dynamique \\
\hline
应该 & yīnggāi & verbe modal & devoir, il faut (conseil moral) \\
\hline
为了 & wèile & préposition & pour, dans le but de (antéposé) \\
\hline
让 & ràng & verbe & faire en sorte que, rendre (qn + adj.) \\
\hline
坚持 & jiānchí & verbe & persévérer, s'accrocher \\
\hline
\end{tabularx}
\end{center}

\courssub{Entraînement : Thème (Français → Chinois) et Version (Chinois → Français)}

\begin{exercice}[Thème : Traduire en chinois (caractères + pinyin)]
\begin{enumerate}
\item Le sport est très important pour la santé.
\item Parce que je veux être en bonne santé, donc je cours chaque matin.
\item Le sport non seulement renforce le corps, mais aussi améliore le moral.
\item Pour bien étudier, nous devons persévérer dans l'exercice physique.
\item En résumé, bougeons ensemble pour un corps sain !
\end{enumerate}
\end{exercice}

\begin{corrige}[Exercice Thème]
\begin{enumerate}
\item \zh{运动对健康很重要。} (Yùndòng duì jiànkāng hěn zhòngyào.) \textit{Structure : Sujet + 对 + Objet + 很 + Adj.}
\item \zh{因为我想健康，所以我每天早上跑步。} (Yīnwèi wǒ xiǎng jiànkāng, suǒyǐ wǒ měitiān zǎoshang pǎobù.) \textit{因为……所以…… ; 每天早上.}
\item \zh{运动不但强身，而且让精神更好。} (Yùndòng búdàn qiángshēn, érqiě ràng jīngshén gèng hǎo.) \textit{\zh{强身} (qiángshēn) = renforcer le corps ; \zh{更} (gèng) = plus.}
\item \zh{为了学习好，我们应该坚持锻炼。} (Wèile xuéxí hǎo, wǒmen yīnggāi jiānchí duànliàn.) \textit{\zh{为了} en tête ; \zh{坚持} + V.}
\item \zh{总之，大家一起运动，为了健康的身体！} (Zǒngzhī, dàjiā yìqǐ yùndòng, wèile jiànkāng de shēntǐ !) \textit{\zh{大家} (dàjiā) tout le monde ; \zh{健康的身体} (corps sain).}
\end{enumerate}
\end{corrige}

\begin{exercice}[Version : Traduire en français]
\begin{enumerate}
\item \zh{因为跑步对身体好，所以我每天坚持。}
\item \zh{打篮球不但有意思，而且能交朋友。}
\item \zh{游泳可以锻炼全身的肌肉。}
\item \zh{总之，健康是幸福的基础，我们要重视体育。}
\item \zh{为了明天的比赛，我们今晚早点睡吧。}
\end{enumerate}
\end{exercice}

\begin{corrige}[Exercice Version]
\begin{enumerate}
\item Parce que la course à pied est bonne pour le corps, je persévère chaque jour.
\item Jouer au basket n'est pas seulement amusant, mais permet aussi de se faire des amis. (\zh{交朋友} jiāo péngyou)
\item La natation permet d'exercer les muscles de tout le corps. (\zh{全身} quánshēn, \zh{肌肉} jīròu)
\item En résumé, la santé est le fondement du bonheur, nous devons accorder de l'importance au sport/à l'EPS. (\zh{幸福} xìngfú, \zh{基础} jīchǔ, \zh{重视} zhòngshì)
\item Pour le match de demain, couchons-nous tôt ce soir ! (\zh{比赛} bǐsài, \zh{早点} zǎodiǎn, \zh{睡} shuì)
\end{enumerate}
\end{corrige}

\courssub{Production guidée : Création de l'affiche (Travail de groupe — 1h30)}

\begin{methode}[Déroulement de l'activité d'intégration (2h au total)]
\begin{enumerate}
\item \textbf{Préparation lexicale (15 min)} : En groupe, liste 10 mots du thème (vocabulaire + structures) sur une fiche brouillon. Vérifie caractères et pinyin.
\item \textbf{Rédaction collaborative (45 min)} : Rédige l'affiche complète : Slogan + 3 arguments (1 \zh{因为……所以……}, 1 \zh{不但……而且……}, 1 libre) + Conclusion (\zh{总之}). Utilise la ponctuation chinoise.
\item \textbf{Évaluation par les pairs (20 min)} : Échange l'affiche avec un autre groupe. Utilise la \textbf{Grille d'évaluation} ci-dessous. Annote au crayon : \textcolor{popgreen}{✓} correct, \textcolor{poporange}{~} à améliorer, \textcolor{poppink}{✗} erreur (caractère, connecteur, ponctuation).
\item \textbf{Mise au propre \& Présentation orale (40 min)} : Recopie en grand format (A3) en soignant la calligraphie (ordre des traits, alignement). Un élève lit l'affiche à haute voix (prononciation, tons). Classe + Professeur : correction collective au tableau des erreurs résiduelles.
\end{enumerate}
\end{methode}

\begin{center}
\begin{tabularx}{\linewidth}{|X|X|}
\hline
\rowcolor{popblueL} \textbf{Critère d'évaluation} & \textbf{Points} \\
\hline
Exactitude des caractères du thème (traits, radicaux, distinctions 运/远, 体/休, 锻/短) & / 6 \\
\hline
Slogan correct (Structure Sujet + 很 + Adj. + \zh{！}) & / 2 \\
\hline
Argument 1 : \zh{因为……所以……} complet et bien placé & / 2 \\
\hline
Argument 2 : \zh{不但……而且……} complet, parallélisme respecté & / 2 \\
\hline
Argument 3 : Structure variée correcte (\zh{可以/能够/有助于/为了}) & / 2 \\
\hline
Conclusion : \zh{总之} en tête + incitation (\zh{吧}/\zh{一起}) & / 2 \\
\hline
Ponctuation chinoise pleine chasse (\zh{，。！：}) + Présentation soignée & / 2 \\
\hline
Expression orale : Prononciation, tons, fluidité, regard public & / 2 \\
\hline
\textbf{TOTAL} & \textbf{/ 20} \\
\hline
\end{tabularx}
\end{center}

\courssub{Approfondissement socioculturel : Sport au Cameroun et en Chine}

\begin{savaistu}[Le sport, pont entre Yaoundé et Pékin]
\begin{itemize}
\item \textbf{Cameroun} : Le football (\zh{足球} zúqiú) est roi. Les \zh{狮子 indomptables} (不屈的狮子 Bùqū de Shīzi) sont une fierté nationale. Le \zh{手球} (shǒuqiú, handball) et l'athlétisme (\zh{田径} tiánjìng) brillent aussi. La \zh{体育周} (tǐyù zhōu) scolaire inclut souvent la danse traditionnelle (\zh{传统舞蹈} chuántǒng wǔdǎo).
\item \textbf{Chine} : Sports nationaux : \zh{乒乓球} (pīngpāngqiú, tennis de table), \zh{羽毛球} (yǔmáoqiú, badminton), \zh{武术} (wǔshù, arts martiaux / kung-fu), \zh{太极拳} (tàijíquán, tai-chi). Depuis les JO de Pékin 2008, le sport scolaire (\zh{体育课} tǐyùkè) est quotidien (course matinale \zh{晨跑} chénpǎo, exercices oculaires \zh{眼保健操} yǎn bǎojiàn cāo).
\item \textbf{Échange} : Les lycées partenaires organisent des matchs amicaux (\zh{友谊赛} yǒuyì sài). Apprendre le vocabulaire du sport permet de communiquer avec les élèves chinois invités : « \zh{你喜欢踢足球吗？} » (Nǐ xǐhuan tī zúqiú ma ? — Tu aimes jouer au foot ?).
\end{itemize}
\end{savaistu}

\courssub{Bilan et devoirs personnels}

\begin{aretenir}[Ce qu'il faut retenir de la Leçon 12]
\begin{itemize}
\item \textbf{Vocabulaire} : Maîtriser l'écriture, le pinyin (tons !) et le sens des 15 mots du tableau.
\item \textbf{Grammaire} : Trois structures argumentatives indispensables : \zh{因为……所以……}, \zh{不但……而且……}, \zh{总之} (+ \zh{为了} en tête de phrase). Respecter l'ordre des mots chinois (Sujet + 很 + Adj ; Complément de but \zh{为了} avant le verbe).
\item \textbf{Caractères} : Savoir écrire de mémoire les 7 caractères analysés (\zh{运动体健康锻炼}) + reconnaître les pièges graphiques (\zh{运/远}, \zh{体/休}, \zh{己/已}, \zh{锻/段}).
\item \textbf{Compétence écrite} : Produire un texte court structuré (Slogan - Arguments - Conclusion) avec ponctuation chinoise.
\item \textbf{Compétence orale} : Lire son texte avec les 4 tons corrects, rythme et intonation d'affiche (impératif / exhortation).
\end{itemize}
\end{aretenir}

\begin{exercice}[Devoirs individuels — À rendre au prochain cours]
\begin{enumerate}
\item \textbf{Rédaction} : Seul(e), écris une deuxième affiche sur le thème \zh{健康饮食} (Jiànkāng yǐnshí — Alimentation saine) / « Mangeons sainement ! ». Structure imposée : Slogan + 3 arguments (1 \zh{因为……所以……}, 1 \zh{不但……而且……}, 1 \zh{为了……}) + Conclusion \zh{总之}. Mots à intégrer : \zh{蔬菜} (shūcài), \zh{水果} (shuǐguǒ), \zh{少吃} (shǎo chī), \zh{垃圾食品} (lājī shípǐn), \zh{营养} (yíngyǎng).
\item \textbf{Révision caractères} : Écris 10 fois chaque caractère du tableau « Analyse graphique » sur ton cahier d'écriture (格子本). Note en marge le radical et le nombre de traits.
\item \textbf{Préparation orale} : Entraîne-toi à lire l'affiche du groupe et ta propre affiche « Alimentation » en t'enregistrant (téléphone). Vérifie les tons 3 (basculant) et les tons neutres (\zh{吧} ba, \zh{的} de).
\end{enumerate}
\end{exercice}

\newpage

\coursec{Méthodes et exercices résolus}

\courssub{Savoir-faire 1 : Lire un modèle de texte commenté}

Avant de rédiger, observe un texte modèle. Lis-le d'abord en entier, puis repère le plan : opinion générale → pratique personnelle → bienfaits → conclusion.

\begin{textebox}[Le sport et moi — texte modèle]
我非常喜欢运动。\\
\emph{Wǒ fēicháng xǐhuan yùndòng.}\\
« J'aime beaucoup le sport. »\\
我每天早上在学校的操场跑步，星期六还和朋友一起踢足球。\\
\emph{Wǒ měitiān zǎoshang zài xuéxiào de cāochǎng pǎobù, xīngqīliù hái hé péngyou yìqǐ tī zúqiú.}\\
« Chaque matin, je cours sur le terrain de l'école, et le samedi je joue encore au football avec mes amis. »\\
因为运动对身体很好，所以我很健康。\\
\emph{Yīnwèi yùndòng duì shēntǐ hěn hǎo, suǒyǐ wǒ hěn jiànkāng.}\\
« Parce que le sport est très bon pour le corps, je suis en bonne santé. »\\
运动还可以让我放松心情。\\
\emph{Yùndòng hái kěyǐ ràng wǒ fàngsōng xīnqíng.}\\
« Le sport me permet aussi de me détendre. »\\
我们应该每天运动！\\
\emph{Wǒmen yīnggāi měitiān yùndòng!}\\
« Nous devrions faire du sport chaque jour ! »
\end{textebox}

\begin{center}
\begin{tabularx}{\linewidth}{|X|X|X|X|}
\hline
\rowcolor{popblueL} Caractères & Pinyin & Nature & Traduction \\
\hline
运动 & yùndòng & nom / verbe & sport ; faire du sport \\
\hline
跑步 & pǎobù & verbe & courir \\
\hline
踢足球 & tī zúqiú & verbe + objet & jouer au football \\
\hline
打篮球 & dǎ lánqiú & verbe + objet & jouer au basket-ball \\
\hline
操场 & cāochǎng & nom & terrain de sport (école) \\
\hline
体育场 & tǐyùchǎng & nom & stade \\
\hline
身体 & shēntǐ & nom & corps, santé \\
\hline
健康 & jiànkāng & adjectif / nom & en bonne santé ; santé \\
\hline
重要 & zhòngyào & adjectif & important \\
\hline
放松 & fàngsōng & verbe & détendre, se détendre \\
\hline
心情 & xīnqíng & nom & humeur, état d'esprit \\
\hline
休息 & xiūxi & verbe & se reposer \\
\hline
比赛 & bǐsài & nom / verbe & match, compétition \\
\hline
应该 & yīnggāi & verbe modal & devoir (conseil) \\
\hline
\end{tabularx}
\end{center}

\begin{methode}[Rédiger un texte simple sur la pratique du sport]
Pour écrire un petit texte (5 à 8 phrases) sur l'importance du sport, suis toujours ces étapes :
\begin{enumerate}
\item \textbf{Poser le sujet} : commence par une phrase générale avec \zh{运动很重要} (le sport est important) ou \zh{我喜欢运动} (j'aime le sport).
\item \textbf{Dire qui, quand, où} : utilise l'ordre chinois \cle{Sujet + temps + lieu + verbe}. Exemple : \zh{我每天早上在学校的操场跑步。} \emph{Wǒ měitiān zǎoshang zài xuéxiào de cāochǎng pǎobù.} « Je cours chaque matin sur le terrain de l'école. »
\item \textbf{Développer avec des connecteurs} : \zh{因为} (parce que), \zh{所以} (donc), \zh{而且} (de plus), \zh{但是} (mais), \zh{先……然后……} (d'abord… ensuite…).
\item \textbf{Donner les bienfaits} : \zh{对身体好} (bon pour le corps), \zh{让我们更健康} (nous rend plus en bonne santé), \zh{可以放松心情} (permet de se détendre).
\item \textbf{Conclure} : une phrase d'appel ou d'opinion : \zh{我们应该多运动！} (Nous devrions faire plus de sport !)
\end{enumerate}
Réflexe : une phrase chinoise = un sujet + un verbe, sans infinitif flottant ni calque du français.
\end{methode}

\begin{popfigure}
\begin{tikzpicture}[node distance=0.35cm]
\node[draw=popblue, fill=popblueL, rounded corners, minimum width=2.1cm, minimum height=0.8cm, align=center] (s) {\zh{我}\\ Sujet};
\node[draw=popgreen, fill=popgreenL, rounded corners, minimum width=2.1cm, minimum height=0.8cm, align=center, right=of s] (t) {\zh{星期六}\\ Temps};
\node[draw=poporange, fill=poporangeL, rounded corners, minimum width=2.1cm, minimum height=0.8cm, align=center, right=of t] (l) {\zh{在体育场}\\ Lieu};
\node[draw=poppurple, fill=poppurpleL, rounded corners, minimum width=2.1cm, minimum height=0.8cm, align=center, right=of l] (v) {\zh{踢足球}\\ Verbe};
\draw[-{Stealth}, thick, popdark] (s) -- (t);
\draw[-{Stealth}, thick, popdark] (t) -- (l);
\draw[-{Stealth}, thick, popdark] (l) -- (v);
\end{tikzpicture}
\legende{L'ordre de la phrase chinoise : Sujet + Temps + Lieu + Verbe. « Le samedi, je joue au football au stade. »}
\end{popfigure}

\begin{exoresolu}[Rédaction guidée : « Le sport et moi »]
Énoncé : En t'aidant des mots \zh{喜欢、每天、跑步、因为、健康}, écris un texte de 5 phrases sur ta pratique du sport.
\tcblower
Solution détaillée :
\begin{enumerate}
\item Sujet général : \zh{我非常喜欢运动。} \emph{Wǒ fēicháng xǐhuan yùndòng.} « J'aime beaucoup le sport. »
\item Qui, quand, où : \zh{我每天早上在学校的操场跑步。} \emph{Wǒ měitiān zǎoshang zài xuéxiào de cāochǎng pǎobù.} « Chaque matin, je cours sur le terrain de sport de l'école. » (temps \zh{每天早上} avant le lieu \zh{在学校的操场}, verbe à la fin.)
\item Cause : \zh{因为跑步对身体很好，所以我很健康。} \emph{Yīnwèi pǎobù duì shēntǐ hěn hǎo, suǒyǐ wǒ hěn jiànkāng.} « Parce que la course est très bonne pour le corps, je suis en bonne santé. » (couple \zh{因为……所以……}.)
\item Bienfait supplémentaire : \zh{运动还可以让我放松心情。} \emph{Yùndòng hái kěyǐ ràng wǒ fàngsōng xīnqíng.} « Le sport me permet aussi de me détendre. »
\item Conclusion : \zh{我们应该每天运动！} \emph{Wǒmen yīnggāi měitiān yùndòng!} « Nous devrions faire du sport chaque jour ! »
\end{enumerate}
Conclusion : le texte respecte le plan sujet → faits → cause → bienfait → appel, avec des connecteurs variés.
\end{exoresolu}

\courssub{Savoir-faire 2 : Écrire correctement les caractères du thème}

\begin{methode}[Maîtriser traits, radicaux et caractères proches du thème du sport]
\begin{enumerate}
\item \textbf{Identifier le radical (clé)} : il donne le sens. \zh{跑} contient la clé du pied \zh{足} (à gauche) ; \zh{健} et \zh{体} contiennent la clé de l'homme \zh{亻} ; \zh{泳} (nager, dans \zh{游泳}) contient la clé de l'eau \zh{氵}.
\item \textbf{Respecter l'ordre des traits} : de haut en bas, de gauche à droite, l'horizontal avant le vertical, l'extérieur avant l'intérieur, on ferme la boîte en dernier (\zh{国} : on trace l'intérieur \zh{玉} avant de fermer le cadre).
\item \textbf{Compter les traits} : \zh{身} (corps) a 7 traits ; \zh{跑} en a 12 ; \zh{健} en a 10 ; \zh{康} en a 11.
\item \textbf{Distinguer les caractères proches} : \zh{体} (tǐ, corps) et \zh{休} (xiū, se reposer) ; \zh{跑} (pǎo, courir) et \zh{跳} (tiào, sauter) ; \zh{场} (chǎng, terrain) et \zh{扬} (yáng, lever).
\item \textbf{Mémoriser par mots entiers} : apprends \zh{身体、健康、运动、操场、跑步} comme des blocs, pas caractère par caractère.
\end{enumerate}
\end{methode}

\begin{exoresolu}[Caractères proches : choisir le bon]
Énoncé : Complète chaque phrase avec le bon caractère : (a) \zh{体 / 休} : 运动以后，我们要 \zh{\_\_} 息一下。(b) \zh{跑 / 跳} : 他每天早上在操场 \zh{\_\_} 步。(c) Recopie \zh{健康} en respectant l'ordre des traits.
\tcblower
Solution détaillée :
\begin{enumerate}
\item (a) La phrase signifie « Après le sport, nous devons nous reposer un peu. » Le mot est \zh{休息} (xiūxi, se reposer) : il faut \zh{休} (homme \zh{亻} + arbre \zh{木} : l'homme s'appuie contre un arbre pour se reposer). \zh{体} signifie « corps » et ne convient pas.
\item (b) « Il court chaque matin sur le terrain » : le mot est \zh{跑步} (pǎobù, courir) : il faut \zh{跑}, clé du pied \zh{足}. \zh{跳} signifie « sauter » : \zh{跳步} n'est pas le mot attendu pour « courir ».
\item (c) \zh{健} : 10 traits — d'abord la clé \zh{亻} (2 traits, de haut en bas), puis la partie droite \zh{建} de haut en bas. \zh{康} : 11 traits — d'abord la clé \zh{广} (3 traits), puis l'intérieur de haut en bas, l'horizontal avant le vertical.
\end{enumerate}
Conclusion : pour choisir entre deux caractères proches, raisonne sur le sens du radical et sur le mot complet.
\end{exoresolu}

\courssub{Savoir-faire 3 : Traduire du français vers le chinois (thème)}

\begin{methode}[Réussir le thème sur le thème du sport]
\begin{enumerate}
\item \textbf{Repérer le sujet et le verbe} de la phrase française, puis reconstruire dans l'ordre chinois : Sujet + temps + lieu + verbe + complément d'objet.
\item \textbf{Traduire « faire du sport »} par \zh{做运动} ou \zh{运动}, « jouer au foot » par \zh{踢足球} (litt. « frapper le ballon »), « jouer au basket » par \zh{打篮球}.
\item \textbf{Placer les compléments de temps et de lieu avant le verbe} : « Je joue au foot avec mes amis le samedi » → \zh{我星期六和朋友一起踢足球。}
\item \textbf{Utiliser les mots-outils justes} : \zh{很} devant un adjectif prédicat (\zh{很健康}), \zh{得} après le verbe pour le complément de manière (\zh{跑得很快}), \zh{了} pour l'accompli.
\item \textbf{Relire} : vérifier les tons du pinyin, les caractères et le classificateur éventuel (\zh{一场比赛}, un match).
\end{enumerate}
\end{methode}

\begin{exoresolu}[Thème : trois phrases]
Énoncé : Traduis en chinois : (1) « Le sport est très important pour la santé. » (2) « Mes amis et moi jouons au football le samedi au stade. » (3) « Il court très vite. »
\tcblower
Solution détaillée :
\begin{enumerate}
\item (1) « Pour la santé » se dit \zh{对健康} ; « important » est \zh{重要}. Réponse : \zh{运动对健康很重要。} \emph{Yùndòng duì jiànkāng hěn zhòngyào.} (structure : Sujet + \zh{对} + objet + \zh{很} + adjectif.)
\item (2) Ordre chinois : sujet \zh{我和朋友} + temps \zh{星期六} + lieu \zh{在体育场} + verbe \zh{踢足球}. Réponse : \zh{我和朋友星期六在体育场踢足球。} \emph{Wǒ hé péngyou xīngqīliù zài tǐyùchǎng tī zúqiú.} « Jouer au foot » = \zh{踢足球}, jamais un calque de « jouer » avec \zh{玩}.
\item (3) Complément de manière : verbe + \zh{得} + adjectif. Réponse : \zh{他跑得很快。} \emph{Tā pǎo de hěn kuài.} Attention au \zh{得} (et non \zh{的}) après le verbe.
\end{enumerate}
Conclusion : chaque phrase suit l'ordre Sujet + temps + lieu + verbe, avec le bon mot-outil.
\end{exoresolu}

\courssub{Savoir-faire 4 : Traduire du chinois vers le français (version)}

\begin{methode}[Réussir la version sur un texte sportif]
\begin{enumerate}
\item \textbf{Lire le texte entier} une première fois pour saisir le sens global avant de traduire.
\item \textbf{Découper en groupes de sens} : repère sujet, temps, lieu, verbe ; ne traduis jamais mot à mot.
\item \textbf{Identifier les mots-outils} : \zh{因为……所以……} = « parce que… donc… » (en français, un seul des deux suffit souvent) ; \zh{对……好} = « bon pour… » ; \zh{让} = « faire que / permettre ».
\item \textbf{Rédiger en français naturel} : accorde les adjectifs, conjugue les verbes, choisis des formulations idiomatiques.
\item \textbf{Vérifier} qu'aucune information du texte chinois n'a été omise (temps, lieu, personnes).
\end{enumerate}
\end{methode}

\begin{exoresolu}[Version : un court texte]
Énoncé : Traduis en français : \zh{喀麦隆的学生很喜欢运动。他们放学以后常常一起踢足球。因为运动对身体好，所以老师也让每个学生多运动。}
\tcblower
Solution détaillée :
\begin{enumerate}
\item Phrase 1 : \zh{喀麦隆的学生} = « les élèves du Cameroun » ; \zh{很喜欢运动} = « aiment beaucoup le sport ». → « Les élèves camerounais aiment beaucoup le sport. »
\item Phrase 2 : \zh{放学以后} = « après la classe » (complément de temps placé avant le verbe en chinois, à replacer naturellement en français) ; \zh{常常一起踢足球} = « jouent souvent au football ensemble ». → « Après les cours, ils jouent souvent au football ensemble. »
\item Phrase 3 : couple \zh{因为……所以……} ; \zh{对身体好} = « est bon pour le corps » ; \zh{让每个学生多运动} = « encourage chaque élève à faire plus de sport » (verbe \zh{让} = faire/permettre). → « Parce que le sport est bon pour la santé, les professeurs encouragent aussi chaque élève à faire plus d'exercice. »
\end{enumerate}
Traduction complète : « Les élèves camerounais aiment beaucoup le sport. Après les cours, ils jouent souvent au football ensemble. Comme le sport est bon pour la santé, les professeurs encouragent aussi chaque élève à faire plus d'exercice. »
Conclusion : une bonne version restitue tout le contenu dans un français correct et naturel.
\end{exoresolu}

\courssub{Savoir-faire 5 : Utiliser les connecteurs et la grille d'évaluation}

\begin{center}
\begin{tabularx}{\linewidth}{|X|X|X|}
\hline
\rowcolor{popblueL} Structure & Exemple en caractères + pinyin & Traduction \\
\hline
因为……所以…… & 因为运动对身体好，所以我很健康。 \emph{Yīnwèi yùndòng duì shēntǐ hǎo, suǒyǐ wǒ hěn jiànkāng.} & Parce que le sport est bon pour le corps, je suis en bonne santé. \\
\hline
而且 & 运动很有意思，而且很便宜。 \emph{Yùndòng hěn yǒu yìsi, érqiě hěn piányi.} & Le sport est intéressant, et de plus il ne coûte pas cher. \\
\hline
但是 & 我喜欢跑步，但是我不喜欢游泳。 \emph{Wǒ xǐhuan pǎobù, dànshì wǒ bù xǐhuan yóuyǒng.} & J'aime courir, mais je n'aime pas nager. \\
\hline
先……然后…… & 我们先跑步，然后踢足球。 \emph{Wǒmen xiān pǎobù, ránhòu tī zúqiú.} & Nous courons d'abord, puis nous jouons au foot. \\
\hline
如果……就…… & 如果天气好，我们就去运动。 \emph{Rúguǒ tiānqì hǎo, wǒmen jiù qù yùndòng.} & S'il fait beau, nous allons faire du sport. \\
\hline
\end{tabularx}
\end{center}

\begin{methode}[Structurer son texte et s'auto-évaluer]
\begin{enumerate}
\item \textbf{Connecteurs à retenir} : \zh{因为……所以……} (cause), \zh{而且} (addition), \zh{但是} (opposition), \zh{先……然后……} (chronologie), \zh{如果……就……} (condition).
\item \textbf{Plan type} : introduction (opinion générale) → pratique personnelle (qui, quand, où) → bienfaits → conclusion (conseil).
\item \textbf{Grille d'auto-évaluation} avant de rendre :
\begin{itemize}
\item Ordre des mots respecté (temps et lieu avant le verbe) ?
\item Caractères du thème écrits sans confusion (\zh{体/休}, \zh{跑/跳}) ?
\item Mots-outils corrects (\zh{很} devant adjectif, \zh{得} de manière, \zh{了} d'accompli) ?
\item Classificateurs justes (\zh{一场比赛}, \zh{一个操场}) ?
\item Au moins deux connecteurs différents utilisés ?
\end{itemize}
\end{enumerate}
\end{methode}

\begin{exoresolu}[Corriger un texte d'élève]
Énoncé : Un élève a écrit : \zh{我踢足球喜欢。我跑步在操场每天。运动很好，因为健康。} Relève et corrige les trois fautes, puis rédige le texte correct.
\tcblower
Solution détaillée :
\begin{enumerate}
\item \zh{我踢足球喜欢。} : calque du français « j'aime jouer au foot ». En chinois, le verbe \zh{喜欢} précède son complément : \zh{我喜欢踢足球。} \emph{Wǒ xǐhuan tī zúqiú.}
\item \zh{我跑步在操场每天。} : le temps et le lieu doivent précéder le verbe. Correction : \zh{我每天在操场跑步。} \emph{Wǒ měitiān zài cāochǎng pǎobù.}
\item \zh{运动很好，因为健康。} : la cause est mal construite ; il faut \zh{因为……所以……} avec une proposition complète : \zh{因为运动对身体好，所以我很健康。} \emph{Yīnwèi yùndòng duì shēntǐ hǎo, suǒyǐ wǒ hěn jiànkāng.}
\end{enumerate}
Texte corrigé : \zh{我喜欢踢足球。我每天在操场跑步。因为运动对身体好，所以我很健康。} « J'aime jouer au football. Je cours chaque jour sur le terrain. Parce que le sport est bon pour le corps, je suis en bonne santé. »
Conclusion : les trois fautes types sont le calque de l'ordre français, la mauvaise place des compléments et le connecteur incomplet.
\end{exoresolu}

\begin{attention}[Les 5 erreurs qui coûtent des points à l'examen]
\begin{enumerate}
\item \textbf{Ordre des mots calqué du français} : le temps et le lieu se placent AVANT le verbe. Écris \zh{我星期六在体育场踢足球。} et jamais \zh{我踢足球在体育场星期六。}
\item \textbf{Confusion de caractères proches} : \zh{体} (corps) contre \zh{休} (se reposer), \zh{跑} (courir) contre \zh{跳} (sauter), \zh{场} (terrain) contre \zh{扬} (lever). Vérifie le radical avant d'écrire.
\item \textbf{Mots-outils confondus} : \zh{的} (possession : \zh{我的操场}), \zh{得} (manière après le verbe : \zh{跑得很快}), \zh{地} (adverbe : \zh{慢慢地跑}). Un seul caractère mal choisi et la phrase est fausse.
\item \textbf{Tons et pinyin négligés} : \zh{tǐyù} (\zh{体育}, sport) n'est pas \zh{tíyǔ} ; une marque de ton fausse ou absente est comptée comme une faute dans l'écriture du pinyin.
\item \textbf{Classificateurs oubliés ou faux} : on dit \zh{一场比赛} (un match), \zh{一个操场} (un terrain), jamais \zh{一比赛}. Et « jouer au foot » se traduit \zh{踢足球}, pas par un calque avec \zh{玩}.
\end{enumerate}
\end{attention}

\begin{exercice}[À toi de jouer : production guidée]
\begin{enumerate}
\item Recopie les mots \zh{健康、操场、跑步、休息} en respectant l'ordre des traits, puis donne le radical de chaque caractère souligné en classe.
\item Thème : traduis en chinois : (a) « Nous faisons du sport tous les jours à l'école. » (b) « Le football est très important pour les élèves camerounais. » (c) « Elle court très vite. »
\item Version : traduis en français : \zh{我哥哥很喜欢打篮球。他每天下午和朋友们在体育场打球。因为运动让他更健康，所以他每天都运动。}
\item Rédaction : écris un texte de 6 à 8 phrases intitulé \zh{运动和我} (Le sport et moi) en utilisant au moins trois connecteurs différents et la grille d'auto-évaluation.
\end{enumerate}
\end{exercice}

\begin{corrige}[Exercice : production guidée]
\begin{enumerate}
\item Radicaux : \zh{健} → clé de l'homme \zh{亻} ; \zh{康} → clé \zh{广} ; \zh{操} → clé de la main \zh{扌} ; \zh{场} → clé de la terre \zh{土} ; \zh{跑} → clé du pied \zh{足} ; \zh{休} → clé de l'homme \zh{亻} ; \zh{息} → clé du cœur \zh{心}.
\item (a) \zh{我们每天在学校做运动。} \emph{Wǒmen měitiān zài xuéxiào zuò yùndòng.} (b) \zh{足球对喀麦隆的学生很重要。} \emph{Zúqiú duì Kāmàilóng de xuésheng hěn zhòngyào.} (c) \zh{她跑得很快。} \emph{Tā pǎo de hěn kuài.}
\item « Mon grand frère aime beaucoup jouer au basket-ball. Chaque après-midi, il joue au ballon avec ses amis au stade. Parce que le sport le rend plus en bonne santé, il fait du sport tous les jours. »
\item Exemple de production attendue : introduction (\zh{我非常喜欢运动。}), pratique (\zh{我每天早上在操场跑步。}), addition (\zh{而且我星期六和朋友踢足球。}), cause (\zh{因为运动对身体好，所以我很健康。}), bienfait (\zh{运动还让我放松心情。}), conclusion (\zh{我们应该多运动！}).
\end{enumerate}
\end{corrige}

\newpage

\coursec{Exercices}

\courssub{Je vérifie mes connaissances}

\begin{exercice}[QCM : vocabulaire du sport]
Pour chaque question, entoure la bonne réponse.

1. « Le sport est important pour la santé » se dit :
\begin{itemize}
\item a. \zh{运动对健康很重要。}
\item b. \zh{健康对运动很重要。}
\item c. \zh{运动很重要对健康。}
\end{itemize}

2. Le mot \zh{锻炼} (duànliàn) signifie :
\begin{itemize}
\item a. se reposer
\item b. s'entraîner
\item c. manger
\end{itemize}

3. Quel connecteur exprime la cause ?
\begin{itemize}
\item a. \zh{而且}
\item b. \zh{所以}
\item c. \zh{因为}
\end{itemize}

4. « Nous devrions faire du sport chaque jour » se dit :
\begin{itemize}
\item a. \zh{我们应该每天运动。}
\item b. \zh{我们每天应该运动。}
\item c. \zh{应该我们每天运动。}
\end{itemize}
\end{exercice}

\begin{exercice}[Vrai ou faux ? Justifie ta réponse]
Dis si chaque affirmation est vraie ou fausse, puis justifie en une phrase en français.

1. Dans \zh{体}, le radical est \zh{亻} (personne).
2. \zh{应该} (yīnggāi) exprime une obligation ou un conseil.
3. Le mot \zh{因为} peut s'utiliser seul, sans \zh{所以}, dans une phrase complète.
4. \zh{远} (yuǎn, loin) et \zh{运} (yùn, transporter/bouger) ont exactement le même sens.
\end{exercice}

\begin{exercice}[Complète le vocabulaire]
Complète chaque phrase avec le mot qui convient, choisi dans la liste : \zh{运动} — \zh{健康} — \zh{身体} — \zh{锻炼} — \zh{重要} — \zh{应该}.

1. 每天跑步对 \_\_\_\_\_\_ 很好。 (Měitiān pǎobù duì \_\_\_\_\_\_ hěn hǎo.)
2. 青少年 \_\_\_\_\_\_ 多运动。 (Qīngshàonián \_\_\_\_\_\_ duō yùndòng.)
3. 踢足球是很好的 \_\_\_\_\_\_ 。 (Tī zúqiú shì hěn hǎo de \_\_\_\_\_\_.)
4. 学习很 \_\_\_\_\_\_ ，运动也很 \_\_\_\_\_\_ 。 (Xuéxí hěn \_\_\_\_\_\_, yùndòng yě hěn \_\_\_\_\_\_.)
\end{exercice}

\begin{exercice}[Associe les caractères au pinyin]
Relie chaque caractère ou mot à sa transcription pinyin.

\begin{center}
\begin{tabularx}{\linewidth}{|X|X|}
\hline
\rowcolor{popblueL} Caractères & Pinyin \\
\hline
1. 运动 & a. duànliàn \\
\hline
2. 锻炼 & b. yīnggāi \\
\hline
3. 健康 & c. yùndòng \\
\hline
4. 应该 & d. zhòngyào \\
\hline
5. 重要 & e. jiànkāng \\
\hline
\end{tabularx}
\end{center}
\end{exercice}

\courssub{Je m'entraîne}

\begin{exercice}[Discrimination : caractères proches]
Complète chaque phrase avec le caractère qui convient : \zh{体} ou \zh{休} ; \zh{练} ou \zh{炼} ; \zh{运} ou \zh{远}.

1. 身 \_\_\_ 健康很重要。
2. 锻 \_\_\_ 对身体好。
3. 我们学校离我家很 \_\_\_ 。
4. 他每天 \_\_\_ 习汉字。
5. 周末我们应该 \_\_\_ 息，也应该运动。
6. \_\_\_ 动让我们心情愉快。
\end{exercice}

\begin{exercice}[Texte à trous : mots-outils et connecteurs]
Complète le texte suivant avec les mots de la liste : \zh{因为} — \zh{所以} — \zh{而且} — \zh{不但} — \zh{也} — \zh{的}.

\begin{textebox}[Texte à compléter]
\_\_\_\_\_\_ 运动对身体好，\_\_\_\_\_\_ 我每天都锻炼。\\
Yīnwèi yùndòng duì shēntǐ hǎo, suǒyǐ wǒ měitiān dōu duànliàn.\\[4pt]
运动 \_\_\_\_\_\_ 让我们健康，\_\_\_\_\_\_ 让我们心情愉快。\\
Yùndòng bùdàn ràng wǒmen jiànkāng, érqiě ràng wǒmen xīnqíng yúkuài.\\[4pt]
我 \_\_\_\_\_\_ 朋友 \_\_\_\_\_\_ 喜欢踢足球。\\
Wǒ de péngyou yě xǐhuan tī zúqiú.
\end{textebox}
\end{exercice}

\begin{exercice}[Remise en ordre]
Remets les phrases dans le bon ordre pour reconstituer un texte argumentatif (thèse, arguments, conclusion). Numérote-les de 1 à 4.

\begin{itemize}
\item[a.] \zh{所以，年轻人应该多运动。} (Suǒyǐ, niánqīngrén yīnggāi duō yùndòng.)
\item[b.] \zh{运动对我们的健康很重要。} (Yùndòng duì wǒmen de jiànkāng hěn zhòngyào.)
\item[c.] \zh{而且，运动让我们心情愉快。} (Érqiě, yùndòng ràng wǒmen xīnqíng yúkuài.)
\item[d.] \zh{因为运动不但对身体好，还能帮助我们学习。} (Yīnwèi yùndòng bùdàn duì shēntǐ hǎo, hái néng bāngzhù wǒmen xuéxí.)
\end{itemize}
\end{exercice}

\begin{exercice}[Traduction thème : français → chinois]
Traduis les phrases suivantes en chinois (caractères et pinyin).

1. Le sport est très important pour la santé.
2. Nous devrions nous entraîner chaque jour.
3. Parce que le sport est bon pour le corps, je joue au football le week-end.
4. Le sport nous rend non seulement en bonne santé, mais aussi joyeux.
\end{exercice}

\courssub{Je me prépare au Bac}

\begin{exercice}[Compréhension de texte et questions de langue]
Lis le texte suivant, puis réponds aux questions.

\begin{textebox}[运动和学习]
很多年轻人觉得学习太忙，没有时间运动。\\
Hěn duō niánqīngrén juéde xuéxí tài máng, méiyǒu shíjiān yùndòng.\\[4pt]
其实，运动对学习也有好处。因为运动不但对身体好，而且让我们心情愉快。\\
Qíshí, yùndòng duì xuéxí yě yǒu hǎochù. Yīnwèi yùndòng bùdàn duì shēntǐ hǎo, érqiě ràng wǒmen xīnqíng yúkuài.\\[4pt]
心情好的时候，我们学习得更好。所以，我们每天应该运动半个小时。\\
Xīnqíng hǎo de shíhou, wǒmen xuéxí de gèng hǎo. Suǒyǐ, wǒmen měitiān yīnggāi yùndòng bàn ge xiǎoshí.\\[4pt]
\emph{Traduction : Beaucoup de jeunes trouvent que les études prennent trop de temps et qu'ils n'ont pas le temps de faire du sport. En réalité, le sport est aussi bénéfique pour les études. Parce que le sport est non seulement bon pour le corps, mais il nous rend aussi de bonne humeur. Quand nous sommes de bonne humeur, nous étudions mieux. C'est pourquoi nous devrions faire du sport une demi-heure chaque jour.}
\end{textebox}

\textbf{Questions de compréhension} (réponds en français) :
\begin{enumerate}
\item Pourquoi beaucoup de jeunes ne font-ils pas de sport ?
\item Quels sont les deux bienfaits du sport mentionnés dans le texte ?
\item Combien de temps devrions-nous faire du sport chaque jour, selon l'auteur ?
\end{enumerate}

\textbf{Questions de langue} :
\begin{enumerate}
\item Relève dans le texte la structure \zh{不但……而且……} et traduis-la.
\item Explique le rôle de \zh{因为……所以……} dans le texte.
\item Trouve dans le texte un verbe modal exprimant le conseil et donne sa traduction.
\end{enumerate}
\end{exercice}

\begin{exercice}[Thème et version type Bac]
\textbf{A. Thème.} Traduis en chinois (caractères et pinyin) :

« Parce que le sport est bon pour la santé, nous devrions nous entraîner chaque jour. »

« Les jeunes devraient faire plus de sport le week-end. »

\textbf{B. Version.} Traduis en français le texte suivant :

\begin{textebox}[Texte à traduire]
运动不但对身体好，而且让我们心情愉快。所以，年轻人应该多运动。\\
Yùndòng bùdàn duì shēntǐ hǎo, érqiě ràng wǒmen xīnqíng yúkuài. Suǒyǐ, niánqīngrén yīnggāi duō yùndòng.
\end{textebox}
\end{exercice}

\begin{exercice}[Expression écrite guidée : 运动和学习]
\textbf{Sujet :} Tu es élève en classe de Première. Ton école organise une semaine du sport. Écris un court texte en chinois intitulé \zh{运动和学习} (Le sport et les études) pour le journal de l'école.

\textbf{Consignes :}
\begin{enumerate}
\item Écris au moins \textbf{6 phrases} (environ 60 à 80 caractères chinois).
\item Utilise au moins \textbf{deux connecteurs} étudiés dans la leçon : \zh{因为……所以……}, \zh{不但……而且……}, \zh{也}.
\item Utilise au moins \textbf{quatre mots du vocabulaire du thème} : \zh{运动}, \zh{健康}, \zh{身体}, \zh{锻炼}, \zh{重要}, \zh{应该}.
\item Soigne l'écriture des caractères : respecte les radicaux et l'ordre des traits (haut → bas, gauche → droite).
\item Termine par une conclusion qui donne ton avis ou un conseil.
\end{enumerate}

\textbf{Grille d'évaluation :} exactitude des caractères (4 pts) ; correction grammaticale et ordre des mots (4 pts) ; emploi des connecteurs (4 pts) ; richesse du vocabulaire du thème (4 pts) ; cohérence du texte, introduction et conclusion (4 pts). Total : 20 pts.
\end{exercice}

\newpage

\coursec{Corrigés des exercices}

\begin{corrige}[Exercice 1 — QCM : vocabulaire du sport]
\textbf{1. Réponse a.} \zh{运动对健康很重要。} (Yùndòng duì jiànkāng hěn zhòngyào.) — « Le sport est très important pour la santé. » Structure : A + 对 + B + 很 + adjectif. La réponse b inverse A et B (sens contraire) ; la réponse c place le complément \zh{对健康} après l'adjectif, ce qui est impossible en chinois : le complément introduit par \zh{对} se place \textbf{avant} l'adjectif.

\textbf{2. Réponse b.} \zh{锻炼} (duànliàn) = « s'entraîner, faire de l'exercice physique ». « Se reposer » se dit \zh{休息} (xiūxi) ; « manger » se dit \zh{吃} (chī).

\textbf{3. Réponse c.} \zh{因为} (yīnwèi) = « parce que », exprime la \textbf{cause}. \zh{所以} (suǒyǐ) exprime la conséquence ; \zh{而且} (érqiě) exprime l'addition.

\textbf{4. Réponse a.} \zh{我们应该每天运动。} (Wǒmen yīnggāi měitiān yùndòng.) Ordre des mots : Sujet + verbe modal \zh{应该} + complément de temps \zh{每天} + verbe. Le modal se place juste après le sujet, jamais en tête de phrase (réponse c fautive) ; le complément de temps se place après le modal (réponse b moins naturelle, ordre déconseillé à ce niveau).
\end{corrige}

\begin{corrige}[Exercice 2 — Vrai ou faux ?]
\textbf{1. Vrai.} Dans \zh{体} (tǐ, corps), le radical est \zh{亻}, la clé de la personne (forme abrégée de \zh{人} à gauche) ; la partie droite \zh{本} (běn) donne une indication phonétique.

\textbf{2. Vrai.} \zh{应该} (yīnggāi) est un verbe modal qui exprime le devoir, le conseil ou ce qui est souhaitable : « devoir, il faudrait ». Exemple : \zh{我们应该多运动。} (Wǒmen yīnggāi duō yùndòng.) « Nous devrions faire plus de sport. »

\textbf{3. Vrai.} Contrairement au français, le chinois accepte \zh{因为} seul (la conséquence étant sous-entendue) ou \zh{所以} seul. Exemple : \zh{因为下雨了，我不去。} ou simplement \zh{因为下雨。} dans la conversation. Mais dans un texte écrit soigné, la paire \zh{因为……所以……} est préférable.

\textbf{4. Faux.} \zh{远} (yuǎn) signifie « loin, lointain » (adjectif), tandis que \zh{运} (yùn) signifie « transporter, bouger » (verbe), présent dans \zh{运动} (sport). Ils partagent seulement l'élément phonétique \zh{元} et la clé du mouvement \zh{辶}, mais leurs sens sont différents.
\end{corrige}

\begin{corrige}[Exercice 3 — Complète le vocabulaire]
1. 每天跑步对\cle{健康}很好。 (Měitiān pǎobù duì jiànkāng hěn hǎo.) « Courir chaque jour est très bon pour la santé. »

2. 青少年\cle{应该}多运动。 (Qīngshàonián yīnggāi duō yùndòng.) « Les adolescents devraient faire plus de sport. »

3. 踢足球是很好的\cle{运动}。 (Tī zúqiú shì hěn hǎo de yùndòng.) « Jouer au football est un très bon sport. » (On accepte aussi \zh{锻炼} : « un très bon exercice ».)

4. 学习很\cle{重要}，运动也很\cle{重要}。 (Xuéxí hěn zhòngyào, yùndòng yě hěn zhòngyào.) « Les études sont très importantes, le sport est aussi très important. »

\emph{Note :} \zh{身体} (shēntǐ, le corps) n'était pas attendu ici, mais la phrase 1 accepterait \zh{身体} avec le sens « bon pour le corps ».
\end{corrige}

\begin{corrige}[Exercice 4 — Associe les caractères au pinyin]
1 — c : \zh{运动} = yùndòng (sport, faire du sport).\\
2 — a : \zh{锻炼} = duànliàn (s'entraîner).\\
3 — e : \zh{健康} = jiànkāng (santé ; en bonne santé).\\
4 — b : \zh{应该} = yīnggāi (devoir, il faudrait).\\
5 — d : \zh{重要} = zhòngyào (important).

\emph{Point de vigilance :} attention au ton de \zh{重} dans \zh{重要} : zhòng (4\up{e} ton), à ne pas confondre avec chóng dans \zh{重复} (chóngfù, répéter).
\end{corrige}

\begin{corrige}[Exercice 5 — Discrimination : caractères proches]
1. 身\cle{体}健康很重要。 (Shēntǐ jiànkāng hěn zhòngyào.) « La santé du corps est très importante. » \zh{体} = corps (radical \zh{亻} + \zh{本}).

2. 锻\cle{炼}对身体好。 (Duànliàn duì shēntǐ hǎo.) « L'entraînement est bon pour le corps. » \zh{炼} avec la clé du feu \zh{火} : « forger, entraîner ».

3. 我们学校离我家很\cle{远}。 (Wǒmen xuéxiào lí wǒ jiā hěn yuǎn.) « Notre école est très loin de chez moi. » \zh{远} = loin.

4. 他每天\cle{练}习汉字。 (Tā měitiān liànxí Hànzì.) « Il s'exerce aux caractères chinois chaque jour. » \zh{练} avec la clé de la soie \zh{纟} : « s'exercer, répéter » (\zh{练习}).

5. 周末我们应该\cle{休}息，也应该运动。 (Zhōumò wǒmen yīnggāi xiūxi, yě yīnggāi yùndòng.) « Le week-end, nous devrions nous reposer et aussi faire du sport. » \zh{休} = se reposer (une personne \zh{亻} contre un arbre \zh{木}).

6. \cle{运}动让我们心情愉快。 (Yùndòng ràng wǒmen xīnqíng yúkuài.) « Le sport nous met de bonne humeur. » \zh{运} = bouger, transporter.

\emph{Moyen mnémotechnique :} \zh{休} (亻+木) : une personne appuyée contre un arbre se repose ; \zh{体} (亻+本) : le corps est la « racine » (本) de la personne.
\end{corrige}

\begin{corrige}[Exercice 6 — Texte à trous : mots-outils et connecteurs]
\cle{因为}运动对身体好，\cle{所以}我每天都锻炼。\\
(Yīnwèi yùndòng duì shēntǐ hǎo, suǒyǐ wǒ měitiān dōu duànliàn.)\\
« Parce que le sport est bon pour le corps, je m'entraîne chaque jour. »

运动\cle{不但}让我们健康，\cle{而且}让我们心情愉快。\\
(Yùndòng bùdàn ràng wǒmen jiànkāng, érqiě ràng wǒmen xīnqíng yúkuài.)\\
« Le sport non seulement nous rend en bonne santé, mais aussi nous met de bonne humeur. »

我\cle{的}朋友\cle{也}喜欢踢足球。\\
(Wǒ de péngyou yě xǐhuan tī zúqiú.)\\
« Mon ami aime aussi jouer au football. »

\emph{Justification :} la paire \zh{因为……所以……} articule cause et conséquence ; \zh{不但……而且……} ajoute un second argument ; \zh{的} est la particule de détermination (mon ami) ; \zh{也} (aussi) se place avant le verbe.
\end{corrige}

\begin{corrige}[Exercice 7 — Remise en ordre]
Ordre correct : \textbf{b (1) — d (2) — c (3) — a (4)}.

1. \zh{运动对我们的健康很重要。} (Yùndòng duì wǒmen de jiànkāng hěn zhòngyào.) — \textbf{Thèse} : le sport est très important pour notre santé.

2. \zh{因为运动不但对身体好，还能帮助我们学习。} (Yīnwèi yùndòng bùdàn duì shēntǐ hǎo, hái néng bāngzhù wǒmen xuéxí.) — \textbf{Premier argument} (cause) : parce que le sport est non seulement bon pour le corps, mais aide aussi à étudier.

3. \zh{而且，运动让我们心情愉快。} (Érqiě, yùndòng ràng wǒmen xīnqíng yúkuài.) — \textbf{Argument supplémentaire} : de plus, le sport nous met de bonne humeur.

4. \zh{所以，年轻人应该多运动。} (Suǒyǐ, niánqīngrén yīnggāi duō yùndòng.) — \textbf{Conclusion} : c'est pourquoi les jeunes devraient faire plus de sport.

\emph{Indice méthodologique :} repérer les connecteurs — \zh{因为} ouvre l'argumentation, \zh{而且} ajoute, \zh{所以} conclut.
\end{corrige}

\begin{corrige}[Exercice 8 — Traduction thème : français → chinois]
1. \zh{运动对健康很重要。} (Yùndòng duì jiànkāng hěn zhòngyào.)\\
Structure : A + 对 + B + 很 + adjectif. Ne pas dire \zh{运动很重要对健康}.

2. \zh{我们应该每天锻炼。} (Wǒmen yīnggāi měitiān duànliàn.)\\
Ordre : Sujet + 应该 + complément de temps + verbe.

3. \zh{因为运动对身体好，所以我周末踢足球。} (Yīnwèi yùndòng duì shēntǐ hǎo, suǒyǐ wǒ zhōumò tī zúqiú.)\\
« Jouer au football » = \zh{踢足球} (litt. « frapper le ballon du pied ») ; le complément de temps \zh{周末} se place après le sujet.

4. \zh{运动不但让我们健康，而且让我们心情愉快。} (Yùndòng bùdàn ràng wǒmen jiànkāng, érqiě ràng wǒmen xīnqíng yúkuài.)\\
Structure d'addition : 不但……而且…… ; « rendre + adjectif » se traduit par \zh{让} + personne + adjectif.
\end{corrige}

\begin{corrige}[Exercice 9 — Compréhension de texte et questions de langue]
\textbf{Questions de compréhension}

1. Beaucoup de jeunes ne font pas de sport parce qu'ils trouvent que les études prennent trop de temps (\zh{学习太忙}) et qu'ils n'ont pas le temps de faire du sport (\zh{没有时间运动}).

2. Les deux bienfaits mentionnés : le sport est bon pour le corps (\zh{对身体好}) et il nous met de bonne humeur (\zh{让我们心情愉快}). Le texte ajoute que, de bonne humeur, on étudie mieux.

3. Selon l'auteur, nous devrions faire du sport une demi-heure chaque jour (\zh{每天应该运动半个小时}).

\textbf{Questions de langue}

1. \zh{运动不但对身体好，而且让我们心情愉快。} (Yùndòng bùdàn duì shēntǐ hǎo, érqiě ràng wǒmen xīnqíng yúkuài.) « Le sport est non seulement bon pour le corps, mais il nous met aussi de bonne humeur. » Structure d'addition : 不但……而且……

2. \zh{因为……所以……} articule la cause et la conséquence : \zh{因为} introduit la raison (le sport est bénéfique), \zh{所以} introduit la conclusion logique (nous devrions faire du sport chaque jour). C'est la charnière argumentative du texte.

3. Le verbe modal de conseil est \zh{应该} (yīnggāi) : « devoir, il faudrait » (dans \zh{我们每天应该运动半个小时}).
\end{corrige}

\begin{corrige}[Exercice 10 — Thème et version type Bac]
\textbf{A. Thème}

Phrase 1 : \zh{因为运动对健康好，所以我们应该每天锻炼。}\\
(Yīnwèi yùndòng duì jiànkāng hǎo, suǒyǐ wǒmen yīnggāi měitiān duànliàn.)\\
Variante acceptée : \zh{运动对健康有好处} (le sport est bénéfique pour la santé).

Phrase 2 : \zh{年轻人周末应该多运动。}\\
(Niánqīngrén zhōumò yīnggāi duō yùndòng.)\\
Variante : \zh{青少年应该在周末多运动。}

\textbf{B. Version}

« Le sport est non seulement bon pour le corps, mais il nous met aussi de bonne humeur. C'est pourquoi les jeunes devraient faire plus de sport. »

\textbf{Barème indicatif (sur 10 pts) :} thème : 3 pts par phrase (exactitude des caractères 1 pt, grammaire et ordre des mots 1 pt, choix du vocabulaire 1 pt) ; version : 4 pts (fidélité du sens 2 pts, qualité du français 2 pts).

\textbf{Points de vigilance :} ne pas inverser \zh{因为}/\zh{所以} ; placer \zh{应该} juste après le sujet ; \zh{多} (davantage) précède directement le verbe \zh{运动} ; en version, ne pas traduire \zh{所以} par « donc » en début de phrase si « c'est pourquoi » est plus élégant.
\end{corrige}

\begin{corrige}[Exercice 11 — Expression écrite guidée : 运动和学习]
\textbf{Proposition de rédaction modèle}

\begin{textebox}[运动和学习 — texte modèle]
很多年轻人觉得学习太忙，没有时间运动。\\
Hěn duō niánqīngrén juéde xuéxí tài máng, méiyǒu shíjiān yùndòng.\\[4pt]
其实，运动对我们的健康很重要。因为运动不但对身体好，而且让我们心情愉快。\\
Qíshí, yùndòng duì wǒmen de jiànkāng hěn zhòngyào. Yīnwèi yùndòng bùdàn duì shēntǐ hǎo, érqiě ràng wǒmen xīnqíng yúkuài.\\[4pt]
心情好的时候，我们学习得更好。\\
Xīnqíng hǎo de shíhou, wǒmen xuéxí de gèng hǎo.\\[4pt]
所以，我们应该每天锻炼半个小时。运动和学习都很重要！\\
Suǒyǐ, wǒmen yīnggāi měitiān duànliàn bàn ge xiǎoshí. Yùndòng hé xuéxí dōu hěn zhòngyào!\\[4pt]
\emph{Traduction : Beaucoup de jeunes trouvent que les études prennent trop de temps et qu'ils n'ont pas le temps de faire du sport. En réalité, le sport est très important pour notre santé. Parce que le sport est non seulement bon pour le corps, mais il nous met aussi de bonne humeur. Quand nous sommes de bonne humeur, nous étudions mieux. C'est pourquoi nous devrions nous entraîner une demi-heure chaque jour. Le sport et les études sont tous les deux très importants !}
\end{textebox}

\textbf{Analyse du texte modèle :} 7 phrases, environ 75 caractères ; connecteurs utilisés : \zh{因为……所以……}, \zh{不但……而且……}, \zh{其实} ; vocabulaire du thème : \zh{运动}, \zh{健康}, \zh{身体}, \zh{锻炼}, \zh{重要}, \zh{应该} (6 mots) ; conclusion personnelle avec \zh{都} et un point d'exclamation.

\textbf{Application de la grille (20 pts) :}
\begin{itemize}
\item Exactitude des caractères (4 pts) : vérifier surtout \zh{锻炼} (clés \zh{钅} et \zh{火}), \zh{健康}, \zh{体} (et non \zh{休}).
\item Correction grammaticale (4 pts) : place de \zh{应该} après le sujet ; complément \zh{对……} avant l'adjectif ; \zh{得} dans \zh{学习得更好}.
\item Connecteurs (4 pts) : au moins deux paires correctement employées.
\item Vocabulaire du thème (4 pts) : au moins quatre mots imposés.
\item Cohérence (4 pts) : introduction (constat), arguments, conclusion avec avis ou conseil.
\end{itemize}

\textbf{Points de vigilance pour les élèves :} ne pas écrire \zh{我很运动} (\zh{运动} n'est pas un adjectif) ; ne pas oublier \zh{的} dans \zh{我们的健康} ; respecter l'ordre des traits pour \zh{健} (d'abord \zh{亻}, puis la partie droite de gauche à droite) ; éviter de traduire mot à mot « faire du sport » par un verbe \zh{做} : on dit \zh{运动} ou \zh{做运动}, jamais \zh{做锻炼} pour « faire du sport » en général.
\end{corrige}

\newpage

\coursec{Fiche bilan}

\begin{popfigure}
\begin{tikzpicture}[mindmap, grow cyclic, every node/.style={concept, text=white, font=\small},
root concept/.append style={concept color=popblue, font=\bfseries},
level 1/.append style={level distance=3.2cm, sibling angle=60},
level 1 concept/.append style={font=\small}]
\node[root concept, align=center]{Leçon 12\\ 运动的重要性}
child[concept color=poporange]{ node[align=center]{Lexique du sport\\ 运动 锻炼 健康 比赛} }
child[concept color=popgreen]{ node[align=center]{Structures\\ 对……有好处\\ 越来越} }
child[concept color=poppurple]{ node[align=center]{Connecteurs\\ 首先 而且 所以 总之} }
child[concept color=poppink]{ node[align=center]{Écriture\\ radicaux 辶 扌 月} }
child[concept color=popgold]{ node[align=center]{Thème\\ français vers chinois} }
child[concept color=popblue!70]{ node[align=center]{Version\\ chinois vers français} };
\end{tikzpicture}
\legende{Carte mentale de la leçon : l'importance de la pratique du sport}
\end{popfigure}

\begin{bilan}
\textbf{Lexique clé à connaître par cœur}
\begin{itemize}
\item \zh{运动} (yùndòng) : le sport, l'exercice ; \zh{锻炼} (duànliàn) : s'entraîner ; \zh{健康} (jiànkāng) : la santé
\item \zh{身体} (shēntǐ) : le corps ; \zh{比赛} (bǐsài) : le match ; \zh{足球} (zúqiú) : le football
\item \zh{重要} (zhòngyào) : important ; \zh{好处} (hǎochu) : l'avantage ; \zh{每天} (měitiān) : chaque jour
\item \zh{坚持} (jiānchí) : persévérer ; \zh{越来越} (yuè lái yuè) : de plus en plus
\end{itemize}

\textbf{Structures grammaticales essentielles}
\begin{center}
\begin{tabularx}{\linewidth}{|X|X|X|}
\hline
\rowcolor{popblueL} Structure & Exemple & Traduction \\
\hline
A 对 B 有好处 & 运动对健康有好处。Yùndòng duì jiànkāng yǒu hǎochu. & Le sport est bon pour la santé. \\
\hline
越来越 + adjectif & 越来越多的人喜欢运动。Yuè lái yuè duō de rén xǐhuan yùndòng. & De plus en plus de gens aiment le sport. \\
\hline
应该 + verbe & 我们应该每天锻炼。Wǒmen yīnggāi měitiān duànliàn. & Nous devrions nous entraîner chaque jour. \\
\hline
不仅……而且…… & 运动不仅有意思，而且有用。Yùndòng bùjǐn yǒu yìsi, érqiě yǒuyòng. & Le sport est non seulement amusant, mais aussi utile. \\
\hline
\end{tabularx}
\end{center}

\textbf{Connecteurs pour le texte} : \zh{首先} (d'abord), \zh{而且} (de plus), \zh{所以} (donc), \zh{但是} (mais), \zh{总之} (en conclusion).

\textbf{Écriture des caractères}
\begin{itemize}
\item \zh{运} : clé du mouvement 辶 (3 traits, tracée en dernier) + \zh{云} ; ne pas confondre avec \zh{远} (yuǎn, loin).
\item \zh{健} : clé de l'homme 亻à gauche, puis \zh{建} à droite (gauche avant droite).
\item \zh{锻} : clé du métal 钅à gauche ; attention à \zh{炼} (clé du feu 火).
\item \zh{赛} : toit 宀 en haut, puis trois horizontales, \zh{贝} en bas (haut vers bas).
\end{itemize}

\textbf{Méthodes express}
\begin{enumerate}
\item Thème : repérer sujet + temps + verbe ; placer le complément de temps avant le verbe ; vérifier les tons et les classificateurs.
\item Version : découper la phrase en groupes de sens, traduire d'abord le verbe, puis remettre les compléments dans l'ordre français.
\item Production : plan en 3 parties (introduction avec \zh{现在}, développement avec \zh{首先}/\zh{而且}, conclusion avec \zh{总之}).
\end{enumerate}
\end{bilan}

\begin{aretenir}[Checklist avant le Bac]
\begin{itemize}
\item[$\square$] Je connais le lexique du sport et de la santé (caractères, pinyin, tons).
\item[$\square$] Je sais écrire \zh{运动}, \zh{锻炼}, \zh{健康}, \zh{比赛} sans erreur de traits.
\item[$\square$] Je distingue les caractères proches : \zh{运}/\zh{远}, \zh{锻}/\zh{段}, \zh{休}/\zh{体}.
\item[$\square$] Je maîtrise la structure \zh{对……有好处} avec la bonne place de \zh{对}.
\item[$\square$] J'utilise correctement \zh{越来越} suivi d'un adjectif.
\item[$\square$] Je place le complément de temps (\zh{每天}, \zh{常常}) avant le verbe.
\item[$\square$] Je sais relier mes idées avec \zh{首先}, \zh{而且}, \zh{所以}, \zh{总之}.
\item[$\square$] Je vérifie les tons dans le pinyin de chaque phrase traduite.
\item[$\square$] Je n'oublie pas la ponctuation chinoise pleine chasse (，。).
\item[$\square$] Je relis mon texte : sujet, verbe, ordre des mots, caractères.
\item[$\square$] Je sais produire un texte de 5 à 8 phrases sur l'importance du sport.
\item[$\square$] Je gère mon temps : 10 minutes de brouillon, puis copie propre.
\end{itemize}
\end{aretenir}

\findecours

\end{document}
